\documentclass[11pt,a4paper]{article}
\usepackage[T1]{fontenc}
\usepackage{lmodern}
\usepackage[margin=2.4cm]{geometry}
\usepackage{amsmath,amssymb}
\usepackage{booktabs,longtable,array}
\usepackage{enumitem}
\usepackage{xcolor}
\usepackage{float}
\usepackage{parskip}
\usepackage[colorlinks=true,linkcolor=blue!50!black,urlcolor=blue!50!black,citecolor=blue!50!black]{hyperref}
% Shared macros for SR_algorithms.tex and every file in sections/.
% Only these packages are available (TinyTeX): amsmath, amssymb, booktabs, longtable, enumitem,
% xcolor, hyperref, float, geometry, tikz, caption. NOT available: algorithm2e, algorithmicx, microtype.
% Pseudocode = numbered enumerate lists.

\definecolor{srccol}{RGB}{90,90,140}
\definecolor{opencol}{RGB}{190,30,30}

% Inline code / identifiers (underscores are safe inside).
% \path (url package, loaded by hyperref) allows line breaks inside long identifiers.
\makeatletter\g@addto@macro\UrlBreaks{\do\:\do\-\do\,}\makeatother
\Urlmuskip=0mu plus 1mu
\urlstyle{tt}
\newcommand{\code}[1]{\ifmmode\text{\path{#1}}\else\path{#1}\fi}
% Source citation for an equation or statement: \src{SRK/fp_likelihood.py:fpt_full}
\newcommand{\src}[1]{{\small\textcolor{srccol}{[\path{#1}]}}}
% Anything not verified against the source, or a question for the author. Collected in the appendix.
\newcommand{\openq}[1]{{\color{opencol}[\textbf{OPEN:} #1]}}

% Model notation (physical)
\newcommand{\xc}{x_c}
\newcommand{\hext}{h_{\mathrm{ext}}}
\newcommand{\eps}{\varepsilon}
% Dimensionless groups (as named in the code)
\newcommand{\reta}{\rho_\eta}
\newcommand{\rbeta}{\rho_\beta}
\newcommand{\reps}{\rho_\varepsilon}
% Probability / survival
\newcommand{\dd}{\mathrm{d}}
\newcommand{\E}{\mathbb{E}}

% Per-algorithm template headings (use in this order; omit none, write "None known." if empty)
\newcommand{\tplPurpose}{\paragraph{Purpose.}}
\newcommand{\tplMath}{\paragraph{Mathematics.}}
\newcommand{\tplImpl}{\paragraph{Discretisation and implementation.}}
\newcommand{\tplLimits}{\paragraph{Assumptions and validity limits.}}
\newcommand{\tplWhere}{\paragraph{Where it lives today.}}
\newcommand{\tplEvidence}{\paragraph{Validation evidence.}}
\newcommand{\tplIssues}{\paragraph{Known issues and planned changes.}}
\newcommand{\tplRefs}{\paragraph{References.}}

% A verify marker (spec 004 FR-020, constitution II): \verify{results JSON}{dotted path}{format spec}{text}. It prints
% the text; a checker recomputes format(value at the path, spec) from the stored result and compares it with the text,
% as tools/verify_readme.py does for the README's <!-- verify:path:key:spec -->text<!-- /verify --> (path elements:
% dict keys or [i] list indices; an empty spec is str(value)).
\newcommand{\verify}[4]{#4}

\emergencystretch=2em
\setlength{\parskip}{0.45em}
\setcounter{tocdepth}{2}

\title{sraging algorithms}
\author{Draft for review}
\date{3 October 2026}

\begin{document}
\sloppy
\maketitle

\begin{abstract}
\noindent This document describes the algorithms of the \code{sraging} library. It grows by one section for
every algorithm that is merged (constitution XII): no algorithm is merged without its section. Each section
says what the code computes, with equations, and cites the ported source as \src{file.py:line}. Anything the
code does that is a known defect is listed with its divergence-register identifier and the test that pins it.
Measured results (coverage, CPU time, nats from benchmark runs) are not quoted here unless the source file
and section they come from are named (constitution II); where a figure matters, this document points to the
register or to the reference record instead of repeating it.
\end{abstract}

\tableofcontents
\clearpage

\section{The calibrator \texttt{profile\_v033}}\label{sec:profile-v033}

\subsection{What it is, and what ``verbatim'' means}\label{sec:pv-intro}

\code{profile_v033} is the reference calibrator \code{mle_profile_jax} v0.3.3 behind named objects: a cohort goes
in, a \code{FitResult} (maximum-likelihood point $\hat\theta$, one confidence claim per parameter, seven
truth-free checks, provenance) comes out. The \emph{numerics are a verbatim port}. Every function that
computes a number lives in \code{src/sraging/compat/v033/} and is a copy of the frozen reference in
\code{tests/oracle/v033/}; \code{VERBATIM.toml} declares, file by file, the only edits allowed (import
rewrites, removal of environment hooks, \code{runtime.configure()} in place of an \code{os.environ} default), and
\code{tests/unit/test_verbatim.py} fails if any other statement differs. Value and gradient bytes and the claims
of a named panel of cohorts are compared side by side with the oracle in \code{tests/parity/} and
\code{tests/acceptance/}. Nothing in this section is therefore a design: it is a description of what the
frozen code does, including its known defects (Section~\ref{sec:pv-defects}). Paths below are relative to
\code{src/sraging/}; \code{compat/v033/} holds the ported numerics and \code{calibrators/profile_v033.py} the thin
wrapper.

Parameters are $\theta=(\eta,\beta,\eps,\kappa)$ with $\xc=1$ fixed. The estimator works on
$z=\log\theta$. Inside \code{compat} $\theta$ is a positional vector in the order
$(\eta,\beta,\eps,\kappa)$; the wrapper converts it to a named \code{SRParams} and asserts its length
\src{calibrators/profile_v033.py:fit}.

\subsection{The computation end to end}\label{sec:pv-pipeline}

\code{profile_v033.fit} calls \code{runtime.check()} (Section~\ref{sec:pv-runtime}), then
\code{estimator.fit} \src{compat/v033/estimator.py:547}, then \code{claims.build}
\src{compat/v033/claims.py:41}, then the seven checks (Section~\ref{sec:pv-checks}). \code{estimator.fit} does,
in order:
\begin{enumerate}[nosep]
  \item \textbf{Scale and grid} (\S\ref{sec:pv-grid}): Kaplan--Meier mean $L$; the automatic grid
    $(\tau_{\max},\Delta\tau)$; the dt-ceiling rule when a recording resolution $\delta$ is given.
  \item \textbf{Compile} the value function and the value-and-gradient function of the log-likelihood over $z$
    (\S\ref{sec:pv-lik}, \S\ref{sec:pv-jax}).
  \item \textbf{Seed} (\S\ref{sec:pv-seed}): a deterministic curve fit to the Kaplan--Meier curve.
  \item \textbf{Maximum likelihood} (\S\ref{sec:pv-mle}): two bounded L-BFGS-B starts, keep the better.
  \item \textbf{Profile each parameter} (\S\ref{sec:pv-profile}): walk outward on each side, repair, polish,
    repair again if the polish changed anything.
  \item \textbf{Interval and verdict} (\S\ref{sec:pv-ci}): \code{profile_ci} turns the node list into an
    interval and a verdict; \code{claims.build} maps verdicts to claim shapes (\S\ref{sec:pv-claims}).
\end{enumerate}
If every MLE start fails, \code{estimator.fit} returns the seed with no intervals; the wrapper then reports
four \code{none} claims with status \code{ok} (D-015) \src{compat/v033/estimator.py:593-595}. sraging keeps the
status for parity but prints \code{RESULT FAILED} and exits with code 5 (FR-018); it happens in none of the
1990 stored v0.3.3 benchmark records.

\subsection{Scale, grid and the dt ceiling}\label{sec:pv-grid}

Let $t_i$ be the ages, $e_i\in\{0,1\}$ the event flags, $\delta$ the recording resolution (or none for exact
times). The time scale is the Kaplan--Meier mean
\begin{equation}
  L=\int_0^{t_{(n)}}\hat S_{\mathrm{KM}}(t)\,\dd t \quad\text{(trapezoid rule over the KM vertices)},
\end{equation}
\src{compat/v033/fp_likelihood.py:125} so $L$ stops at the last age and depends on censoring; $\tau=t/L$.
The solver horizon and step come from the data:
\begin{align}
  \tau_{\max}&=\min\!\big(40,\ \max(2,\ 1.15\,t_{(n)}/L)\big), &
  \Delta\tau_{\rm auto}&=\max\!\big(2\times10^{-3},\ \tau_{\max}/1200\big),
  \label{eq:pv-autogrid}
\end{align}
\src{compat/v033/fp_likelihood.py:141} and the dt-ceiling rule \src{compat/v033/estimator.py:530}
\begin{equation}
  \Delta\tau=\begin{cases}\min\big(\Delta\tau_{\rm auto},\ \tfrac{\delta/L}{3}\big), & \delta>0\ \text{and}\ \delta/L<\Delta\tau_{\rm auto},\\[2pt]
  \Delta\tau_{\rm auto}, & \text{otherwise.}\end{cases}
  \label{eq:pv-ceiling}
\end{equation}
The remaining grid settings are fixed: $N=900$ space cells, $G_{\rm tol}=10^{-11}$, scheme \code{bdf2},
\code{sampling_dt}$=\delta$ \src{compat/v033/estimator.py:541-543}. The environment hook
\code{SRK_DT_CEILING} of the reference \code{auto_grid} is deleted in the port (register D-103); the
rule~\eqref{eq:pv-ceiling} is applied in \code{solver_grid} instead, which is the behaviour of the reference
with the variable unset as it was in the v0.3.3 runs.
$L$ is used for more than the likelihood's numerical unit: it also sets the grid, the corner trigger
$\rho_\beta<10^{-3}$ and the seed box (register D-010). Its Kaplan--Meier code does not process deaths before
censorings at a tie, so $L$ depends on the row order of tied ages (D-007).

\subsection{The likelihood}\label{sec:pv-lik}

\paragraph{Model.} With $y=X/\xc$, the damage process is
$\dd y=(\rho_\eta\tau-\rho_\beta\,y/(y+\kappa))\,\dd\tau+\sqrt{2\rho_\eps}\,\dd W$, reflecting at $y=0$, death at
the first passage of $y=1$, with
\begin{equation}
  \rho_\eta=L^2\eta,\qquad \rho_\beta=L\beta,\qquad \rho_\eps=L\eps,\qquad\xc=1
  \label{eq:pv-groups}
\end{equation}
\src{compat/v033/fp_jax.py:355}. The density $p(y,\tau)$ of surviving damage obeys
$\partial_\tau p=-\partial_y(a\,p)+\rho_\eps\,\partial_y^2p$ with drift $a(y,\tau)=\rho_\eta\tau-\rho_\beta y/(y+\kappa)$,
$p(y,0)=\delta_0$, and survival $S(\tau)=\int_0^1p\,\dd y$.

\paragraph{March.} On $N$ cells of width $h=1/N$ with centres $y_j=(j+\tfrac12)h$, the removal term
$r_j=\rho_\beta y_j/(y_j+\kappa)$ is time-invariant. Scharfetter--Gummel face coefficients use the
Bernoulli function $B(z)=z/(e^z-1)$, argument clipped to $[-500,500]$, with
$P_j=(\rho_\eta\tau_n-r_j)\,h/\rho_\eps$ and $c=\rho_\eps/h^2$. The tridiagonal operator $A(\tau_n)$ has
\begin{equation}
  \text{sub}_i=c\,B(-P_{i-1}),\qquad \text{diag}_i=-c\big(B(-P_i)+B(P_{i-1})\big),\qquad \text{super}_i=c\,B(P_i),
\end{equation}
with $\text{diag}_0=-cB(-P_0)$ at the reflecting wall and no flux out of the last cell other than through the
absorbing face \src{compat/v033/fp_jax.py:141-159}. Time stepping is backward Euler for the first
\code{nstart}$=20$ steps (damping the $\delta$-function start), then BDF2:
\begin{equation}
  \big(\tfrac32 I-\Delta\tau\,A(\tau_{n+1})\big)\,p^{n+1}=2p^n-\tfrac12p^{n-1},
  \qquad p^{n+1}\leftarrow\max(p^{n+1},0),\qquad S_{n+1}=h\sum_jp^{n+1}_j,
\end{equation}
\src{compat/v033/fp_jax.py:167-173, 201-202}. The march stops once $S<G_{\rm tol}$ and $\tau>0.3$
\src{compat/v033/fp_jax.py:206}. The survival array is then made monotone (running minimum, clipped to
$[0,1]$) and the density is $p=-\nabla_\tau S$, clipped at zero \src{compat/v033/fp_jax.py:312-324}
(\code{_truncate}).

\paragraph{Terms.} Survival is read by linear interpolation of the $\tau$ grid (value 1 to the left, the
last value to the right). For individual $i$, with $\tau_i=t_i/L$ and $\tau_i^+=(t_i+\delta)/L$:
\begin{align}
  \ell(\theta)&=\sum_{i:\,e_i=1}\log^{\ge}\!\big(P_i\big)+\sum_{i:\,e_i=0}\log^{\ge}\!\big(S(\tau_i)\big),\qquad
  \log^{\ge}(x)=\max\!\big(\log\max(x,10^{-305}),\,-700\big),\\
  P_i&=\max\big(S(\tau_i)-S(\tau_i^+),\,0\big)\ \ (\delta\ \text{given}),\\
  P_i&=p(\tau_i)/L\ \ (\text{exact times; }0\text{ beyond the stop index}),
\end{align}
and $\ell=-\infty$ if the sum is not finite \src{compat/v033/fp_jax.py:360-370}. The floor constants are
\code{_LOG_FLOOR}$=-700$ and \code{_P_FLOOR}$=e^{-700}$ \src{compat/v033/fp_jax.py:74-75}. Two consequences
are ported as they are: the interval mass is a \emph{difference} of two survival values, which underflows to
exactly $0$ when $S\approx1$ and such a term then scores $-700$ or a rounding-dependent value (D-004); and a
region where every term sits on the floor has an exactly zero gradient, which the optimiser handles with its
fallbacks (\S\ref{sec:pv-mle}).

\subsection{Differentiating the march: the JAX twin and its custom VJP}\label{sec:pv-jax}

\code{fp_jax.make_loglik} returns two compiled functions over $z=\log\theta$: a value function with the
early-exit march, and \code{jax.value_and_grad} of a second program that uses a hand-written reverse sweep
\src{compat/v033/fp_jax.py:327-375}. They are different compiled programs, so their values can differ in the
last bit; the code never mixes them in one comparison and keeps separate memo caches
\src{compat/v033/estimator.py:72-142}.

The reverse sweep (\code{_march_vjp}) replaces automatic differentiation of the scan.
For step $n$ let $M_n=\alpha_nI-\Delta\tau A(\tau_{n+1})$ ($\alpha_n=1$ for backward Euler, $\tfrac32$ for BDF2),
$r_n$ the right-hand side, and $q_{n+1}=M_n^{-1}r_n$, so $p^{n+1}=\max(q_{n+1},0)$. The forward pass stores only
$q_{n+1}$ per step. The backward pass carries two adjoint vectors $\lambda_{n+1},\lambda_{n+2}$,
\begin{equation}
  \bar q=\big(h\,\bar S_{n+1}+a_{n+1}\lambda_{n+1}+b_{n+2}\lambda_{n+2}\big)\odot\mathbf 1[q_{n+1}>0],
  \qquad \lambda_n=M_n^{-\top}\bar q,
\end{equation}
($a=2$, $b=-\tfrac12$ under BDF2; $a=1$, $b=0$ under backward Euler; the indicator has value $\tfrac12$ at a tie,
as \code{lax.max} defines it), then pulls the band cotangent $-\lambda_n\,q_{n+1}^{\top}$ back through the
recomputed coefficient function to $(\rho_\eta,\rho_\eps,r)$ \src{compat/v033/fp_jax.py:215-235}. One transposed
tridiagonal solve per step; steps after the early-stop index $k$ carry their cotangent onto $S_k$
(\code{_fold_tail}). When the stored $q$ would exceed $64$~MiB (\code{Q_BUDGET_BYTES}
\src{compat/v033/fp_jax.py:86}) the sweep stores a checkpoint every $\lceil\sqrt{n_{\rm steps}}\rceil$ steps
and re-marches each block (\code{_march_vjp_ckpt}), which changes memory and CPU but not the gradient.
The Bernoulli pair $(B(P),B(-P))$ uses one \code{expm1} ($B(-a)=a+B(a)$) with an analytic derivative
$B'(z)=B(z)(1-B(-z))/z$ and an even series below $|z|=1.5$ \src{compat/v033/fp_jax.py:97-138}.

\subsection{The seed}\label{sec:pv-seed}

\code{fp_curvefit_seed} \src{compat/v033/fp_curvefit.py:276} fits the Kaplan--Meier curve, deterministically.
The curve is resampled at $60$ survival levels uniform from $1-10^{-6}$ down to $2\times10^{-3}$
(\code{NGRID}, \code{_empirical_on_taugrid}), and $z=\log(\rho_\eta,\rho_\beta,\rho_\eps,\kappa)$ is fitted by
bounded trust-region least squares of $G(\tau;z)-S_{\rm emp}$ on a coarser solver grid
(\code{N_FIT}$=900$, \code{DT_FIT}$=2\times10^{-3}$, \code{TAU_MAX}$=16$), from two starts, the production start
$(3.5,0.2,0.10,10^{-3})$ and the removal start $(1,6,0.5,0.3)$ \src{compat/v033/fp_curvefit.py:59-62}; the
lower survival-only RMS wins. Bounds are $\rho\in[10^{-4},\,300]$ and $\kappa\in[10^{-3},5]$
\src{compat/v033/fp_curvefit.py:27-30, 130}. The result is converted to physical units with $\xc=1$ and the
Kaplan--Meier $L$. The cap $\rho\le300$ is why large, sharp cohorts are seeded far from the truth (D-001).
The extra ``surface'' starts and the composite (hazard, CV, mean) residuals exist in the ported code but are off
in this call: \code{surface_start=False} and all weights zero by default.

\subsection{Maximum likelihood: two bounded L-BFGS-B starts}\label{sec:pv-mle}

The objective is $-\ell$ on $z$, with a box check on $\theta$: $-\ell=+\infty$ outside
\begin{equation}
  \eta,\beta,\eps\in[10^{-12},10^3],\qquad\kappa\in[10^{-3},\kappa_{\max}=300]
\end{equation}
\src{compat/v033/estimator.py:67-69, 114-127}. Two starts are run from the seed
$\theta_s$: $\log\theta_s$ and the same with $\eps$ multiplied by $4$ \src{compat/v033/estimator.py:577-579}.
Each is a bounded L-BFGS-B on the exact gradient with \code{maxiter}$=1000$, \code{maxfun}$=2000$,
\code{ftol}$=10^{-13}$, \code{gtol}$=10^{-6}$ \src{compat/v033/estimator.py:51}. The best finite result is kept
\src{compat/v033/estimator.py:583-591}.

Safety net (\code{_lbfgs} \src{compat/v033/estimator.py:145-188}). A gradient method cannot leave a region where the
gradient is exactly zero. If the start value is not finite or its gradient is identically zero, the search
first tries L-BFGS-B from the seed (when one is given and it has a nonzero gradient) and then falls back to
Nelder--Mead on the value function; a fallback never returns something worse than its start. No check is made
that the MLE stage converged: an \code{ABNORMAL} or false-\code{ftol} stop is accepted (D-002).

\subsection{The profile walk, repair and polish}\label{sec:pv-profile}

For each parameter $p$ (index $i$) with free set $F$ (the other three), the profile
$\ell_p(v)=\max_{z_F}\ell(z_F,\ \log v)$ is evaluated at nodes $v$ placed by an adaptive walk, then repaired and
polished. The edge threshold is $\mathrm{thr}=\chi^2_{1,0.95}$ from SciPy \src{compat/v033/estimator.py:598}.

\paragraph{Node solve} (\code{_profile_ll_at} \src{compat/v033/estimator.py:257}). Warm-started bounded L-BFGS-B
over $z_F$ (\code{maxiter}$=400$ in the walk, $600$ in repair and polish; \code{maxfun}$=800$, same tolerances).
If the warm start ends in a degenerate corner, the node is re-solved from the seed's nuisance values and the
better solution is kept. The corner triggers are $\rho_\beta=L\beta<10^{-3}$, or a free parameter within
$10^{-6}$ (in $\log$) of a box edge \src{compat/v033/estimator.py:191-212}. On a walk side an adaptive rule
decides whether the restart runs: a $\rho_\beta$ trigger always restarts; otherwise after $2$ consecutive
restarts gaining less than $10^{-3}$ nats the side stops restarting until the trigger set changes or the walk
sees a deviance jump, when the skipped restart is run first (\code{_want_restart}, \code{obj.pending}
\src{compat/v033/estimator.py:215-254}).

\paragraph{Walk} (\code{_walk_side} \src{compat/v033/estimator.py:312-394}). From the MLE value $\hat v$ on each side,
with $\ell_{\rm ref}=\hat\ell$, step $s$ in decades, initial $s=0.05$:
\begin{enumerate}[nosep]
  \item $v'=v\cdot10^{\pm s}$, snapped to $1.01\times$ the box edge if it would pass it (reason \code{box_edge}
    if already there).
  \item Solve the node; deviance $D=2(\ell_{\rm ref}-\ell_p(v'))$ and jump if $D-D_{\rm prev}>4$ or not finite.
  \item On a jump: re-arm the restart rule; if $s>10^{-4}$ then $s\leftarrow\max(s/3,10^{-4})$ and
    retry (a \emph{rejection}); otherwise record the node and stop with reason \code{cliff}.
  \item Otherwise accept the node; stop with reason \code{dev_stop} if $D\ge10$; else $s\leftarrow\min(3s,2)$ if
    $D<1$, or $\min(2s,2)$ if $D-D_{\rm prev}<0.3$ (cap $2$ decades); stop at the box edge.
\end{enumerate}
A side ends after $n_{\rm acc}=\texttt{per\_side}$ accepted nodes or $24$ rejections, both reported as
\code{max_nodes}. The v0.3.3 change (P1) is that rejections no longer use up node slots
\src{compat/v033/estimator.py:62-64, 344}. The budget is
$\texttt{per\_side}=\max\big(24,\lfloor(n_{\rm profile}-1)/2\rfloor+3\big)$ \src{compat/v033/estimator.py:601},
which is $24$ for any $n_{\rm profile}<43$, so the setting $n_{\rm profile}=11$ has no effect (D-012).

\paragraph{Repair} (\code{_repair} \src{compat/v033/estimator.py:397-424}). On the value-sorted nodes, up to $3$
passes: for a node $j$ left (right) of the best node, if some node further out is higher by more than $0.2$
nats, re-solve $j$ from that node's solution and keep the result if higher.

\paragraph{Polish} (\code{_polish_edges} \src{compat/v033/estimator.py:460-527}). On each side, the first node outward with
$D\ge\mathrm{thr}$ places the interval edge. It is re-solved from fresh starts: the MLE's nuisance values, the
seed's, the next node outward, and the opposite $\kappa$ corners ($\kappa$ at its floor, at $1$ and at its
ceiling, with $\beta$ rescaled to keep $\beta/(1+\kappa)$; \code{_kappa_corner_starts}
\src{compat/v033/estimator.py:427-457}). A gain above $10^{-3}$ nats is kept; if the node now falls below
$\mathrm{thr}$ the edge moves outward (at most $12$ moves per side), and if that passes the last node the walk
is continued once. A polish can only raise node values. If the polish changed anything the repair runs once more
\src{compat/v033/estimator.py:621-628}.

\paragraph{Re-centring diagnostic.} If some node has $\ell_p>\hat\ell+0.01$ the parameter is recorded as
\code{recentred} with \code{gain_nats}$=\max\ell_p-\hat\ell$ \src{compat/v033/estimator.py:631-640}. The point
$\hat\theta$ is \emph{not} refreshed: it stays at the stale MLE (D-003).

\subsection{From profile to interval and verdict}\label{sec:pv-ci}

\code{profile_ci} \src{compat/v033/fp_gof.py:17-128} takes the nodes $(v_k,\ell_k)$, discards non-finite
values (fewer than $3$ remaining gives verdict \code{insufficient}), and forms the deviance relative to the
\emph{profile} maximum,
$d_k=2(\max_m\ell_m-\ell_k)$, and $\mathrm{thr}=3.8415$ at $95\%$ (table \code{_CHI2_1},
\src{compat/v033/fp_gof.py:14}; the walk and polish use the SciPy value, which differs in the fifth digit; D-016).
The interval is the whole set $\{d\le\mathrm{thr}\}$, not the component around the maximum: the edges are
the first crossings outward from the leftmost and rightmost accepted nodes, interpolated linearly in
$\sqrt d$ against $\log v$ (linear in $v$ if a value is non-positive); \code{n_components} counts runs of
accepted nodes. The verdict is
\begin{center}
\begin{tabular}{@{}ll@{}}
\toprule
Crossing found & Verdict \\
\midrule
both sides & \code{identified} \\
above only ($v_{\rm hi}$ finite) & \code{one-sided:<} (upper bound) \\
below only & \code{one-sided:>} (lower bound) \\
none & \code{flat} \\
fewer than $3$ finite nodes & \code{insufficient} \\
\bottomrule
\end{tabular}
\end{center}
(the function's docstring states the two one-sided cases the other way round; D-013; the code is as tabulated.)
Because $d$ is taken against the profile maximum, a re-centred profile gives an interval around a different
centre than $\hat\theta$ (D-003).

\subsection{Claims}\label{sec:pv-claims}

\code{claims.build} \src{compat/v033/claims.py:41-75} maps each parameter's verdict and interval to a claim,
with fixed constants \code{VERDICT_TO_SHAPE}: \code{identified}$\to$\code{closed},
\code{one-sided:<}$\to$\code{upper_bound}, \code{one-sided:>}$\to$\code{lower_bound}, \code{flat} and
\code{insufficient}$\to$\code{none} \src{compat/v033/claims.py:24-30}. Then:
\begin{enumerate}[nosep]
  \item A \code{closed} claim missing a finite end is demoted to \code{upper_bound} if only the upper end is
    finite, \code{lower_bound} if only the lower, else \code{none}; an \code{upper_bound} without a finite upper end
    and a \code{lower_bound} without a finite lower end become \code{none}.
  \item Ends are clipped to the box $\mathrm{BOUNDS}$ (\code{claims.py:13-18}): $\eta,\beta,\eps\in[10^{-12},10^3]$,
    $\kappa\in[10^{-3},300]$. Lower ends apply for \code{closed} and \code{lower_bound}, upper ends for
    \code{closed} and \code{upper_bound}.
  \item For every shape other than \code{none} the point is $\hat\theta_p$ clipped to the same box; level $0.95$;
    \code{n_components} is the \code{Claim} default $1$ (it is never filled from \code{profile_ci}; D-014).
\end{enumerate}
Only these four shapes on $\eta,\beta,\eps,\kappa$ exist; the wrapper adds a human-readable label to each claim
(box-limited raw-$\beta$ upper bounds carry a note; D-009) without changing a number
\src{calibrators/profile_v033.py:fit}. The diagnostics dictionary is the estimator's diagnostics without the
profile node lists, plus \code{loglik} and \code{dt_ceiling}.

\subsection{Explicit settings}\label{sec:pv-settings}

Every setting is explicit (FR-011); nothing is read from the environment (D-103). The values below are the
entries of \code{profile_v033.SETTINGS} \src{calibrators/profile_v033.py:26-38}; the third column gives where the
number is consumed in the ported code.

\begin{center}\small
\begin{longtable}{@{}p{3.2cm}p{4.4cm}p{7.2cm}@{}}
\toprule
Setting & Value & Defined / consumed at \\
\midrule
\endhead
\code{kappa_max} & $300$ & \code{profile_v033.py:27}; box upper edge for $\kappa$, \code{estimator.py:67-69} \\
\code{levels} & $(0.95,)$ & \code{profile_v033.py:28}; \code{estimator.py:598}, \code{claims.py:22} \\
\code{n_profile} & $11$ (inert, D-012; see \S\ref{sec:pv-profile}) & \code{profile_v033.py:29}; \code{estimator.py:547, 601} \\
\code{polish} & \code{True} & \code{profile_v033.py:30}; \code{estimator.py:621} \\
\code{box} & $\eta,\beta,\eps:[10^{-12},10^3]$; $\kappa:[10^{-3},300]$ & \code{profile_v033.py:31}; \code{estimator.py:67-69} \\
\code{dt_ceiling_rule} & $\delta/L/3$ when $\delta/L<\Delta\tau_{\rm auto}$ & \code{profile_v033.py:32}; \code{estimator.py:537-540} \\
\code{N} & $900$ & \code{profile_v033.py:33}; \code{estimator.py:541} \\
\code{Gtol} & $10^{-11}$ & \code{profile_v033.py:34}; \code{estimator.py:541}, \code{fp_jax.py:206} \\
\code{scheme} & \code{bdf2} & \code{profile_v033.py:35}; \code{estimator.py:541} (the JAX march is always BDF2 after $20$ BE steps, \code{fp_jax.py:328}) \\
\code{pjrt_nproc} & $1$ & \code{profile_v033.py:36}; \code{runtime.py:18} \\
\code{x64} & \code{True} & \code{profile_v033.py:37}; \code{runtime.py:23-33} \\
\bottomrule
\end{longtable}
\end{center}

Constants that are not wrapper settings but fix the result are listed with their source lines in the
sections above: auto-grid constants $(40,2,1.15,2\times10^{-3},1200)$ \code{fp_likelihood.py:141,155-156};
$\tau>0.3$ early-stop guard \code{fp_jax.py:206,315}; log floor $-700$ \code{fp_jax.py:74}; walk constants
\code{_DEV_STOP}$=10$, \code{_DEV_MAX_STEP}$=4$, \code{_STEP0_DEX}$=0.05$, \code{_STEP_MIN_DEX}$=10^{-4}$,
\code{_STEP_MAX_DEX}$=2$, \code{_REPAIR_MIN_NATS}$=0.2$, \code{_MAX_REJECTIONS}$=24$ (\code{estimator.py:41-46, 64});
optimiser budgets \code{_LBFGS_MLE}, \code{_LBFGS_NODE} (\code{estimator.py:51-52}); restart and polish
tolerances (\code{estimator.py:55-60}); re-centring tolerance $0.01$ nats (\code{estimator.py:631}). The
reference descriptor's declared cost (\code{REFERENCE_COST_SECONDS}, \code{profile_v033.py:42}) is a quotation
of the reference, not a measurement made here, and is not reproduced.

\subsection{Runtime requirements}\label{sec:pv-runtime}

\code{runtime.py} sets and checks two facts the ported numerics need. (1) \emph{float64}:
\code{jax_enable_x64} is set at import, and \code{check()} refuses to fit if it is off, because the $-700$ log
floor and $G_{\rm tol}=10^{-11}$ need double precision \src{compat/v033/runtime.py:23-52}. (2)
\emph{One CPU thread}: \code{PJRT_NPROC=1} pins XLA's CPU client to a single thread. It is
\emph{overwritten}, not defaulted, so a shell variable cannot change the runtime
\src{compat/v033/runtime.py:18, 33}; it is read when the backend initialises, so \code{check()} raises
\code{RuntimeConfigError} if the backend was already initialised under another value
\src{compat/v033/runtime.py:44-60}. In practice: import \code{sraging} before any JAX computation, in a fresh
process. Byte identity with the oracle holds only for the pinned environment on the host that runs the test
(Python~3.14.6, JAX~0.11.1, NumPy~2.4.3, SciPy~1.17.1; see \code{tests/oracle/v033/pip_freeze.txt}). The
runtime state is recorded in the fit's provenance.

\subsection{The truth-free checks}\label{sec:pv-checks}

\code{warnings.evaluate} \src{warnings.py} runs the checks after every fit (five in spec 001; seven since spec 002); they never change a
number. The conditions and thresholds of the first four are verbatim from the reference entry point
\code{tools/fit_cohort_jax.py} (the module header states the lines); \code{tests/parity/test_warnings_parity.py}
compares outcomes with the reference functions on stored records.
\begin{center}\small
\begin{tabular}{@{}p{2.6cm}p{8.6cm}p{3.2cm}@{}}
\toprule
Code & Fires when & Threshold \\
\midrule
\code{BASIN_GAIN} & some re-centred parameter has $\mathtt{gain\_nats}>1.92$: the walk found a point that much better than the reported best fit & $1.92$ nats (\code{warnings.py:16}) \\
\code{BASIN_OUTSIDE} & a \code{closed} claim has its point outside its own interval, $\mathrm{lower}(1-10^{-9})\le\mathrm{point}\le\mathrm{upper}(1+10^{-9})$ violated & relative $10^{-9}$ (\code{warnings.py:38}) \\
\code{CENSOR_TIE} & $\delta$ given, some censoring, the censoring age is the maximal age, and the number of deaths at that age (\code{isclose}) is at least $0.5\%$ of $n$ & $0.005\,n$ (\code{warnings.py:17, 53}) \\
\code{DELTA_GRID} & $\delta$ given and some age is not a multiple of $\delta$ ($|t/\delta-\mathrm{round}(t/\delta)|>10^{-6}$) & any, tolerance $10^{-6}$ (\code{warnings.py:67}) \\
\code{BOX_EDGE} & $\hat\theta$ for some parameter on an \emph{upper} edge of the prior box ($\eta,\beta,\eps$ at $10^3$, $\kappa$ at $\kappa_{\max}=300$); sraging's own (FR-017b), not in the reference & relative $10^{-3}$ \\
\bottomrule
\end{tabular}
\end{center}
The first two detect the wrong-basin failures of Section~\ref{sec:pv-defects}; the third detects the
censor-tie bias (D-006); the fourth is a data check; the fifth flags a fit that ran up the ridge of D-020 to the box (D-018): on
the 1990 stored v0.3.3 records it fires on at least 215 (a \code{none} claim hides $\hat\theta$, so not every record
can be scored), and a closed $\eta$ interval of a flagged fit covers the true $\eta$ in 8 of 214 cases. A quiet
\code{BOX_EDGE} does not make an interval right. No check detects the $N=900$ barrier of D-005. A check
that does not fire is not evidence that a fit is right; the checks record both whether they ran and whether they
fired.

\subsection{What \texttt{profile\_v033} does not do}\label{sec:pv-notcovered}

From the Not-Covered section of the spec (\code{specs/001-reference-fit/spec.md}, constitution XII):
\begin{itemize}[nosep]
  \item \textbf{Model extensions.} No external (Makeham) hazard, no fit window, no heterogeneity, and $\xc$ is
    fixed at $1$ (no $\xc\ne1$ fitting): the model is the four-parameter $\theta$ above. Windowed observations
    are refused (D-104).
  \item \textbf{Claim shapes.} Only \code{closed}, \code{upper_bound}, \code{lower_bound}, \code{none} on
    $\eta,\beta,\eps,\kappa$. No combination claims ($\beta/(1+\kappa)$, $\bar r$, $\bar r/\eta$, $\sqrt\eta/\eps$,
    $\beta/(\eta(\kappa+1))$) and no \code{negligible} or \code{open above: box-limited} shape; both arrive in
    spec~003. The labels on claims (D-009) cover the gap for users.
  \item \textbf{Leverage check.} Deferred to spec~005: the reference likelihood returns no per-individual terms.
  \item \textbf{Times in $L_{\rm SR}$.} No $L$-scaled quantity is reported; the numerics scale by the
    Kaplan--Meier $L$ as v0.3.3 does (D-010).
  \item \textbf{Engines and uncertainty.} FP only (no finite-$n$ engine, spec~004); no replicates.
  \item \textbf{Out of scope entirely:} SRtools compatibility, unit conversions, plots, data cleaning, Langevin,
    interventions, MCMC, and speed parity (spec~002). The default-versus-\code{v033} comparison test of
    constitution VII has nothing to compare until the first numerics divergence (spec~002).
  \item \textbf{Known limitations inherited deliberately:} wrong fits on large sharp human cohorts are warned
    about, not fixed; the censor-tie bias is warned about, not corrected; the $N=900$ barrier is not detected by
    any truth-free check; byte identity holds only in the pinned environment on the testing host.
\end{itemize}

\subsection{Known defects, ported as-is, and sraging differences}\label{sec:pv-defects}

Entries D-001 to D-021 are properties of the reference that sraging keeps on purpose (status
\emph{ported-as-is}; impact ``none: not fixed''); D-101 to D-105 and D-018's warning are deliberate differences of sraging from
the reference. The register is \code{divergences/register.toml}; the last column is its \code{regression_test}.
Measured consequences are in the register entry and its cited source, not repeated here.

\begin{center}\footnotesize
\begin{longtable}{@{}p{1.1cm}p{7.3cm}p{6.4cm}@{}}
\toprule
ID & Description & Pinned by \\
\midrule
\endhead
D-001 & Curve-fit seed caps $\rho$ at $300$ (\code{RHO_CEIL}), so large sharp (human) cohorts are seeded far from the truth; flagged by \code{BASIN_GAIN}/\code{BASIN_OUTSIDE}. & \code{tests/unit/test_asis_defects.py::test_d001_seed_rho_cap} \\
D-002 & No MLE-stage convergence check; an \code{ABNORMAL} or false-\code{ftol} L-BFGS-B result is accepted. & \code{tests/acceptance/test_results.py::test_basin_cohorts_reproduced} \\
D-003 & Recentring is diagnostic only: the point stays stale and can lie outside its own interval. & \code{tests/acceptance/test_results.py::test_basin_cohorts_reproduced} \\
D-004 & Interval-mass cancellation: $S(t)-S(t+\delta)$ underflows to $0$ when $S\approx1$; head deaths score at the floor or a rounding-dependent value. & \code{tests/unit/test_asis_defects.py::test_d004_interval_mass_is_a_difference} \\
D-005 & At $N=900$ the fit grid invents a barrier on the p15 cohorts; no truth-free warning detects it. & \code{tests/acceptance/test_results.py::test_basin_cohorts_reproduced} \\
D-006 & Censor-tie deaths are scored over the whole end-of-study recording bin (biases $\eta$ upward); flagged by \code{CENSOR_TIE}. & \code{tests/unit/test_warnings.py::test_censor_tie_at_exactly_half_a_percent_fires_and_just_below_does_not} \\
D-007 & Kaplan--Meier tie order: with a tied death and censoring the result depends on row order; $L$ sets the grid. & \code{tests/unit/test_asis_defects.py::test_d007_kaplan_meier_depends_on_row_order} \\
D-008 & \code{fpt_full}: a \code{scheme} other than \code{bdf2} silently runs backward Euler (the match is case-insensitive, \code{str(scheme).lower()=="bdf2"}, \code{fp_likelihood.py:116}); not reachable through the public API. & \code{tests/unit/test_asis_defects.py::test_d008_unknown_scheme_runs_backward_euler} \\
D-009 & $\kappa_{\max}=300$ box: raw-$\beta$ upper bounds may be set by the box rather than the data; an \code{upper_bound} claim carries the box label, but a \code{closed} $\beta$ whose walk ended at the box (\code{grid.beta.up = box_edge}) carries none. & \code{tests/unit/test_labels.py::test_raw_beta_upper_bound_carries_the_box_note} \\
D-010 & $L_{\rm KM}$ is load-bearing beyond the likelihood's numerical unit (grid, corner trigger, seed box, box extent). & \code{tests/unit/test_verbatim.py::test_port_is_verbatim} \\
D-011 & The dt-ceiling rule $\delta/L/3$ rests on one measurement and dominates the cost of the reference; re-measured in spec~002. & \code{tests/unit/test_verbatim.py::test_port_is_verbatim} \\
D-012 & $n_{\rm profile}=11$ is inert: the walk allows at least $24$ nodes per side whatever $n_{\rm profile}$ is. & \code{tests/unit/test_asis_defects.py::test_d012_n_profile_is_inert} \\
D-013 & The \code{profile_ci} docstring states the one-sided verdicts the wrong way round; the code (\code{one-sided:<}$\to$\code{upper_bound}) is what runs. & \code{tests/unit/test_asis_defects.py::test_d013_one_sided_mapping_follows_the_code} \\
D-014 & \code{Claim.n_components} is never filled from \code{profile_ci}: every claim says $1$ and a disconnected region is reported as its hull. & \code{tests/unit/test_asis_defects.py::test_d014_n_components_is_always_one} \\
D-015 & When every MLE start fails, the result is status \code{ok} with four \code{none} claims; sraging prints \code{RESULT FAILED} and exits 5. & \code{tests/unit/test_asis_defects.py::test_d015_all_starts_failed_maps_to_four_none_claims} \\
D-016 & Two $\chi^2$ thresholds: polish and walk use $3.841459$ (SciPy) against the MLE, \code{profile_ci} uses $3.8415$ (table) against the profile maximum. & \code{tests/unit/test_asis_defects.py::test_d016_two_chi2_thresholds} \\
D-017 & \code{make_loglik} normalises only a scalar \code{sampling_dt}; an array is used without shape or positivity checks. & \code{tests/unit/test_asis_defects.py::test_d017_sampling_dt_scalar_branch_only} \\
D-018 & $\hat\theta$ on an upper box edge: the fit ran up the ridge of D-020 and its intervals sit too high (the true $\eta$ lies below the interval in 143 of 147 stored fits with $\beta$ on the edge and a closed $\eta$); flagged by \code{BOX_EDGE} (8 of 214 flagged fits cover $\eta$). & \code{tests/unit/test_asis_defects.py::test_d018_closed_interval_ends_at_the_beta_box_edge} \\
D-019 & Results depend on the host (machine and software environment, not separable here): the stored v0.3.3 records of the 10 named panel cohorts (WEXAC) differ from the parity Mac's oracle run on all 10, by up to $3.2$ dex on the 4 basin cohorts and at most $0.15$ dex on the others, with claim shapes flipping on two. & \code{tests/unit/test_asis_defects.py::test_d019_named_cohorts_differ_between_hosts} \\
D-020 & (model) A Gompertz-like ridge: $\eta,\beta,\eps$ grow together with $\eta/\eps$ nearly fixed; on the Drosophila cohort the log-likelihood changes by $0.42$ nats from $\beta=30$ to $3\times10^4$. & \code{tests/unit/test_asis_defects.py::test_d020_ridge_keeps_eta_over_epsilon} \\
D-021 & A one-ulp change of $\log\theta$ moves the log-likelihood by 60--200 times what its gradient predicts (up to $1.1\times10^{-6}$ nats); line searches can stall on this noise. & \code{tests/unit/test_asis_defects.py::test_d021_loglik_moves_at_the_ulp_scale} \\
\midrule
D-101 & (refused input) CSV rows the reference reader skips or recodes are refused, with every offending row listed. & \code{tests/unit/test_csv_reader.py::test_every_offending_row_is_listed_in_one_error} \\
D-102 & (refused input) The recording resolution $\delta$ must be stated; continuous times must be said explicitly. & \code{tests/unit/test_cohort_core.py::test_delta_rule_explains_why} \\
D-103 & (diverged) Environment overrides are not read; their v0.3.3 values are explicit settings (Section~\ref{sec:pv-settings}). & \code{tests/unit/test_env_ignored.py::test_overrides_do_not_change_bytes} \\
D-104 & (refused input) Windowed observations are refused by the tester adapter. & \code{tests/unit/test_adapter_offline.py::test_window_is_refused} \\
D-105 & (diverged) With $\delta$, a death at age $0$ (found at the first check) is kept and scored over $[0,\delta)$; the reference reader drops it. The estimator handles it exactly as the oracle's does, but where the first interval carries almost no risk the term is the $-700$ floor (D-004). & \code{tests/unit/test_cohort_core.py::test_a_death_at_age_zero_with_delta_is_the_first_interval} \\
\bottomrule
\end{longtable}
\end{center}

The two behaviours of the claim path noted in Sections~\ref{sec:pv-pipeline} and~\ref{sec:pv-claims} are registered as D-015 (all starts failed) and D-014 (\code{n_components}).

\subsection{Relation to the v0 reference document}\label{sec:pv-v0}

This section is derived from \code{docs/planning/algorithms/} (sections on the solver, grid, JAX likelihood and
profile estimator). Where the ported code differs from that document in form, the code governs: the dt ceiling
is applied in \code{solver_grid} (\code{estimator.py:530-544}) rather than by writing and restoring an environment
variable; \code{fpt_full} compares \code{str(scheme).lower()=="bdf2"} (\code{fp_likelihood.py}, as in the oracle),
so the scheme match is case-insensitive, as D-008 now states, where the v0 document describes it as an exact comparison;
and citations here use line numbers of the ported files, not of the oracle.

% Section for specs/002-fp-engine (T062). Paths are relative to src/sraging/ unless they start with tests/ or SRK/.
\providecommand{\todo}[1]{{\color{opencol}[\textbf{TODO:} #1]}}
\newcommand{\ulp}{u}
\newcommand{\Fabs}{F}

\section{The FP engine (\texttt{sraging.engine})}\label{sec:fp}

\subsection{Overview}\label{sec:fp-overview}

\code{sraging.engine} is one Fokker--Planck march and one observation model. It replaces SRK's \code{fpt_full},
\code{fp_loglik}, \code{fp_loglik_het_physical} and \code{fp_piecewise} (which duplicated the stepper) and
v0.3.3's JAX twin. Everything is switched by static options; a feature that is off contributes \emph{no traced
operation}. Two numerics versions exist:
\begin{itemize}[nosep]
  \item \textbf{\code{"v033"}} with every feature off, grid policy \code{"v033"} and scheme \code{"bdf2"}: the
    engine's value and value-and-gradient programs have the \emph{same canonical program text} as
    \code{compat/v033/fp_jax.py:make_loglik} (StableHLO with symbols renamed in order of appearance,
    \src{engine/program.py:canonical_text}) and give the same value and gradient bytes
    \src{engine/stepper.py:make_march}. The gate is a program comparison, not a timing.
  \item \textbf{\code{"v1"}} (the default) $=$ \code{"v033"} $+$ D-004, D-106, D-107, D-108 (Sections
    \ref{sec:fp-obs}, \ref{sec:fp-guard}). Everything else (D-105, D-109 to D-115) is common to both numerics.
\end{itemize}
D-115, D-116 and D-117 are ported as they are in both numerics (Section~\ref{sec:fp-limits}).

\paragraph{Parameter layout.} The jitted kernels take one vector $z=\log(\text{free parameters})$ in the order
carried by a \code{ParamLayout} \src{engine/layout.py:ParamLayout}. The default layout is
$(\eta,\beta,\eps,\kappa)$, i.e.\ exactly v0.3.3's $z$. Other names that may be free: \code{xc},
\code{h_ext}, \code{het_v}, \code{seg:<id>}; inputs that snap to the grid (entry ages, window, $\delta$, segment
edges, a segment's $\xc$) have no gradient and are refused (FR-008). Everything not free is \emph{fixed} for the
life of the evaluator (compiled in as a constant); parameters that disagree with a fixed value are refused.
No public function takes a bare positional vector.

\paragraph{Evaluator API.} \code{Evaluator(cohort, free=, fixed=, numerics="v1", grid=GridPolicy(...),
scheme="bdf2", segments=, heterogeneity=)} compiles once \src{engine/evaluator.py:Evaluator}. Public methods
take a named \code{SRParams}: \code{value}, \code{value_and_grad} (gradient by name), \code{terms} (per-row log
terms and an auxiliary dictionary), \code{weighted_grad} ($\nabla\sum_iw_i\ell_i$), \code{survival_on_grid},
\code{stiffness}, \code{nodes_dropped}, \code{m_report}, \code{objective} (the two jitted functions over $z$) and
\code{program_text}. Curves without a cohort: \code{engine.curves.curves} \src{engine/curves.py:curves}.

Notation. Time is in units of $L$ ($\tau=t/L$), $h=1/N$, and the dimensionless groups are those of
Eq.~\eqref{eq:pv-groups}. The step matrix, the Scharfetter--Gummel coefficients, the clip, $S=h\sum_jp_j$, the
stop rule ($S<G_{\rm tol}$ and $\tau>0.3$), the monotone truncation and the floor constants
$\log^{\ge}$ are those of Section~\ref{sec:pv-lik}
\src{engine/stepper.py:_coeffs, _step_q, truncate}. The sections below state what the engine adds.

\subsection{The time mesh and the march}\label{sec:fp-mesh}

\paragraph{Meshes.} A \code{Grid} carries either a uniform step ($\tau_n=n\Delta\tau$, $n_{\rm steps}=\lfloor
\tau_{\max}/\Delta\tau\rfloor$, as v0.3.3) or an increasing node list $0=\tau_0<\dots<\tau_{n}$
\src{engine/grid.py:Grid}. The graded start of a fit-time reference is
\begin{equation}
  \Delta_0=\Delta\tau_0,\quad\Delta_{j+1}=r\,\Delta_j\ \text{while}\ \Delta_j<\Delta\tau,\quad\text{then uniform steps }\Delta\tau\text{ to }\tau_{\max},
\end{equation}
\src{engine/grid.py:graded_start}; it is passed to the policy \code{"fixed"} as \code{nodes=}. The step
$\Delta_n=\tau_{n+1}-\tau_n$ is used for step $n$.

\paragraph{Segment edges are nodes (D-113).} Segment edges $e_s/L$ (finite, in $(0,\tau_{\rm end})$) are made
nodes by \src{engine/grid.py:insert_edges}: an edge on a node (to relative $10^{-12}$) changes nothing, so a
grid-aligned configuration keeps the uniform mesh like for like with \code{fp_piecewise}; otherwise the
\emph{nearest} node is moved onto the edge (the step count is unchanged and every step stays in
$[\Delta\tau/2,\,3\Delta\tau/2]$) and the step that starts at that node is a backward Euler restart. Two edges
claiming one node (a segment shorter than about one step) are refused. \code{fp_piecewise} instead switched
parameters at the node \emph{below} the edge while its clock ran from the exact edge (phase error of up to one
step; \code{tests/oracle/fp_piecewise_cb191f5}). Hazard-segment edges never touch the mesh: the hazard acts outside
the PDE (Section~\ref{sec:fp-obs}).

\paragraph{Variable-step BDF2.} With $w_n=\Delta_n/\Delta_{n-1}$ the step is
\begin{equation}
  (\alpha_nI-\Delta_nA_n)\,q^{n+1}=a_n\,p^n+b_n\,p^{n-1},\qquad
  \alpha_n=\frac{1+2w_n}{1+w_n},\quad a_n=1+w_n,\quad b_n=-\frac{w_n^2}{1+w_n},
  \label{eq:fp-bdf2}
\end{equation}
$p^{n+1}=\max(q^{n+1},0)$, and $(\alpha,a,b)=(1,1,0)$ (backward Euler) for $n<n_{\rm start}=20$, at $n=0$ and at
every restart \src{engine/stepper_gen.py:coefficients}. Note $a_n+b_n=\alpha_n$ (consistency), and on a uniform
mesh $w=1$ gives $(\tfrac32,2,-\tfrac12)$, Section~\ref{sec:pv-lik}. Scheme \code{"be"} is backward Euler throughout.

\paragraph{The general march.} \src{engine/stepper_gen.py:make_march, _step}. It marches any node list with
piecewise-constant parameters. The time axis is cut into \emph{spans} at the (finite) edges of the march's
segments; step $n$ uses the span containing the midpoint of $(\tau_n,\tau_{n+1})$ (ties to the earlier segment,
as \code{fp_piecewise}). For span $k$ the multipliers $\eta_k,\beta_k,\eps_k,\kappa_k$ give
$(\rho_\eta,\rho_\eps,r_j)_k$ on the common reference threshold $x_{\rm ref}=\max_k\xc^{(k)}$
\src{engine/evaluator.py:_theta_gen}. The production term is evaluated at the accumulated DPUs $Y$ (damage production
units; $A$, the \emph{DPU accumulation rate}, speeds up production only: removal, noise and the time step are not rescaled; a model that rescales all rates is time rescaling,
not A)
\begin{equation}
  \begin{gathered}
  Y(\tau)=Y_k+A_k(\tau-\tau_k)\ \ (\tau\in\text{span }k),\qquad Y_0=0,\ \ Y_{k+1}=Y_k+A_k(\tau_{k+1}-\tau_k),\\
  P_j=(\rho_\eta^{(k)}Y(\tau_{n+1})-r_j^{(k)})\,h/\rho_\eps^{(k)},
  \end{gathered}
\end{equation}
\src{engine/stepper_gen.py:_clock_starts}; with no $A$ segment $Y=\tau$ literally. A span whose threshold is
$\xc^{(k)}$ uses the first $M_k=\mathrm{round}(N\xc^{(k)}/x_{\rm ref})$ cells (the rest are masked out of the
solve); a raised threshold exposes empty cells, a lowered one is a \emph{lowering}, below. $M_k<2$ is refused
(D-114). The stop rule uses $\tau_{n+1}>0.3$ on the node, and \code{truncate_gen} returns the same quantities as
\code{truncate} on any mesh (density by the nonuniform \code{gradient}, one-sided at the stop)
\src{engine/stepper_gen.py:truncate_gen}. One span on a uniform mesh equals the uniform path to round-off.

\paragraph{The absorbed mass $F$.} Under numerics \code{"v1"} the march also carries $F\approx1-S$. Summing the
step over cells, the column sums of $\alpha I-\Delta A$ vanish except in the last cell, where the outflow is
$h\,c\,B(-P_{M-1})\,q_{M-1}$ with $c=\rho_\eps/h^2$. Because $a_n+b_n=\alpha_n$,
\begin{equation}
  F^{n+1}=\frac{a_n\,F^n+b_n\,F^{n-1}+\Delta_n\,\beta_n\,q^{n+1}_{M-1}}{\alpha_n}+h\sum_j\min\!\big(q^{n+1}_j,0\big),\qquad
  \beta_n=\frac{\rho_\eps}{h}\,B\!\big(-P_{M-1}\big),
  \label{eq:fp-F}
\end{equation}
\src{engine/stepper.py:_F_next, _bflux_coef; engine/stepper_gen.py:_step}: the scheme's own mass balance plus the
mass the clip adds back. Every term is a small positive quantity, so $F$ keeps full relative precision where
$S\approx1$ (where $1-S$ would not). $F_0=0$. The $S$ path is operation for operation the one without $F$.
\code{truncate_F} gives $F$ on the truncated grid (running maximum, clipped to $[0,1]$, held at $F_k$ after
the stop).

\paragraph{A lowering carries into BDF2's history (T079).} At the first step $n$ of a span with $M_k<M_{k-1}$,
the mass above the new threshold, $J=h\sum_{j\ge M_k}p^n_j$, is absorbed at once: $p^n_j\leftarrow0$ for $j\ge M_k$
and $F^n\leftarrow F^n+J$. The step is a backward Euler restart, so $F^{n-1}$ is not used in it. But the \emph{next}
step is BDF2 and uses the pre-step $F^{n}$ as its history. If that history excludes the jump, then Eq.~\eqref{eq:fp-F}
at step $n+1$ counts $J$ with weight $(a_{n+1}F^{n+1}+b_{n+1}F^n)/\alpha_{n+1}$, i.e.\ $\tfrac43J-\tfrac13\cdot0=\tfrac43J$
(for $w=1$) instead of $J$. The march therefore returns, with each new $F$, the post-lowering $F$ as the next step's
history \src{engine/stepper_gen.py:_step, return \code{pn, Fn, F}} (the same applies to $F^{n-1}\leftarrow F^{n-1}+J$,
which is moot because the step is backward Euler). Before this fix (review T064, T079) BDF2 counted the jump $4/3$
times.

\subsection{Gradients}\label{sec:fp-grad}

\paragraph{Uniform mesh: custom VJP.} For the uniform march the gradient is v0.3.3's hand-written adjoint of
$V=\sum_nS_n$-type functionals (Section~\ref{sec:pv-jax}; recurrence for $\lambda$, mask $\mathbf1[q>0]$, tail
folding) \src{engine/stepper.py:make_march, _adjoint_step}, with the same switch: all $q^{n+1}$ stored when
$(n_{\rm steps}+1)\,N\cdot8\le64$~MiB, otherwise a checkpoint $(p^n,p^{n-1})$ every
$K=\lceil\sqrt{n_{\rm steps}}\rceil$ steps and a re-march per block \src{engine/stepper.py:grad_mode}.
With $F$ carried (numerics \code{"v1"}) the VJP has a \emph{second adjoint} $\nu_m=\partial L/\partial F^m$
\src{engine/stepper.py:_adjoint_step_F}:
\begin{equation}\begin{gathered}
  \nu_{n+1}=\bar F_{n+1}+\frac{a_{n+1}}{\alpha_{n+1}}\,\nu_{n+2}+\frac{b_{n+2}}{\alpha_{n+2}}\,\nu_{n+3},\\
  \bar q=\bar p\odot\mathbf 1[q>0]+h\,\nu_{n+1}\,(1-\mathbf 1[q>0])+\frac{\Delta\,\beta_n\,\nu_{n+1}}{\alpha_n}\,e_{N-1},
\end{gathered}\end{equation}
where $\bar p=h\bar S_{n+1}+a_{n+1}\lambda_{n+1}+b_{n+2}\lambda_{n+2}$ is v0.3.3's and (uniform BDF2)
$a/\alpha=\tfrac43$, $b/\alpha=-\tfrac13$. The coefficient cotangent is pulled back through the recomputed band
\emph{and} $\beta_n$ (the cotangent $\nu\,\Delta\,q_{N-1}/\alpha$). The checkpoint stores only $(p,p_{-})$: $F$ is not
needed backward, only $\nu$. Without $F$ the code path is v0.3.3's, statement for statement.

\paragraph{General march: autodiff of a checkpointed scan.} The general march is differentiated by JAX, not by a
hand-written sweep: the gradient program marches the \emph{fixed length} $n_{\rm steps}$ in a \code{lax.scan} over
blocks of $K\approx\lceil\sqrt{n_{\rm steps}}\rceil$ steps, each block under \code{jax.checkpoint}, with a
\code{live} mask for the tail \src{engine/stepper_gen.py:make_march, mode "grad"}. Past the stop index the
survival is replaced by its held value in \code{truncate_gen}, so the fixed-length scan and the early-stopping
march compute the same function, and the gradient is exact for it by construction.

\paragraph{Values.} The value programs always use the early-stopping \code{lax.while_loop} (the value path is not
reverse-differentiable); entries after the stop index repeat $S_k$ (and $F_k$).

\subsection{The observation model}\label{sec:fp-obs}

\src{engine/observation.py:Assembly}. It has two stages so that a heterogeneity mixture can sit between them.

\paragraph{Stage 1: points.} From one march, \code{points} returns the survival quantities at the observation
ages, by linear interpolation in $\tau$ ($S=1$, $F=0$ to the left of the grid, the last value to the right):
$S_i=S(\tau_i)$, the density $p_i$ (exact times; $0$ beyond the stop, as v0.3.3), $S(\tau_i^+)$ and, under \code{"v1"}
for a death with resolution, $F(\tau_i^+)$, $F(\tau_s)$, $S(\tau_s)$ ($\tau_s=\max(t_i,e_i)/L$, below), and
$S(e_i)$, $S(t_{\rm lo})$, $S(t_{\rm hi})$ where entry or window are set. All are \emph{linear} in $S$ and $F$, and
nothing nonlinear is applied yet.

\paragraph{Stage 2: assemble.} Applies the hazard, the death rule, entry, window, the floor
$\log^{\ge}(x)=\max(\log\max(x,e^{-700}),-700)$ \src{engine/observation.py:_logp}, and returns the per-row
terms $\ell_i$; the total is one sum in input order, $-\infty$ if not finite.

\paragraph{External hazard.} A constant rate $h_e$, or a piecewise-constant one by hazard segment (rates $r_k$ on
spans $[a_k,a_k+w_k)$), with the cumulative hazard integrated exactly,
\begin{equation}
  \begin{gathered}H(x)=\sum_kr_k\,\mathrm{clip}(x-a_k,0,w_k),\qquad S_{\rm obs}(x)=S(x)\,e^{-H(x)},\\ f_{\rm obs}(x)=\big(f(x)+h(x)S(x)\big)e^{-H(x)},\end{gathered}
\end{equation}
\src{engine/observation.py:_H, _E, _rate, _dH}; for a constant rate $H=h_et$, SRK's \code{fp_loglik}. A death
beyond the march keeps the hazard-only mass (register D-115, Section~\ref{sec:fp-limits}).

\paragraph{Per-row terms (common structure).} Let $d_i$ be the death indicator.
\begin{equation}
  \ell_i=\Big[\,d_i\,\log^{\ge}\!\big(m_i\big)+(1-d_i)\log^{\ge}\!\big(S_{\rm obs}(t_i)\big)\Big]-\log^{\ge}\!\big(S_{\rm obs}(e_i)\big)\ \ (\text{entry}),
  \label{eq:fp-term}
\end{equation}
with $m_i$ the death mass (below). For exact times $m_i=f_{\rm obs}(t_i)$ (density in physical time, $p_i/L$). For a
death recorded with resolution $\delta_i$ (one value, or one per row) it is the interval mass over
$[t_s,\,t_i+\delta_i)$, with $t_s=t_i$ and in v0.3.3
\begin{equation}
  m_i^{033}=\max\!\big(S_{\rm obs}(t_s)-S(t_i+\delta_i)\,e^{-H(t_i+\delta_i)},\,0\big).
  \label{eq:fp-m033}
\end{equation}

\paragraph{Entry (left truncation).} Per-row entry ages $e_i$ condition each row on survival to $e_i$:
$-\log^{\ge}S_{\rm obs}(e_i)$ in Eq.~\eqref{eq:fp-term}, in both numerics (no legacy formula exists; research R8).
A death whose recording interval contains its entry ($e_i>t_i$, $e_i<t_i+\delta_i$) is scored from the entry:
$t_s=\max(t_i,e_i)$ in Eq.~\eqref{eq:fp-m033} and $\Delta$ below uses $t_i+\delta_i-t_s$. A cohort has either entry or a
window, not both (\code{cohort.py:_entry}).

\paragraph{Exact-time deaths by density.} $m_i=p_i/L$ with the hazard composition $(p_i/L+h\,S_i)e^{-H(t_i)}$. Under
\code{"v1"} $p_i$ itself is rebuilt from $F$ where the $S$ difference is round-off (D-106): at node $j$, with the
same stencil as the $S$ density (central, one-sided at the ends), $p_j^F=\nabla_\tau F_j$ is used instead of
$-\nabla_\tau S_j$ when $F(\tau_{j+1})<S(\tau_{j-1})$ (the difference whose operands are smaller has the smaller
ulp); where $S$'s is kept the result is v0.3.3's bit for bit \src{engine/stepper.py:density_F; engine/stepper_gen.py:density_F_gen}.
For a heterogeneous cohort the choice is made per node before mixing.

\paragraph{The fit window $[t_{\rm lo},t_{\rm hi}]$.} A common window is the special case of entry
(spec FR-012). Rows with $t_i<t_{\rm lo}$ are excluded (term $0$, counted). Every at-risk row is conditioned on
$S_{\rm obs}(t_{\rm lo})$. Rows that reach past $t_{\rm hi}$ are censored at $t_{\rm hi}$:
$\ell_i=\log^{\ge}$-style $\log\max(S_{\rm obs}(t_{\rm hi}),10^{-300})$, floored at $-700$. The rules differ:
\begin{center}\small
\begin{tabular}{@{}p{3.2cm}p{5.4cm}p{5.4cm}@{}}\toprule
 & \code{"v033"} \src{SRK/fp_likelihood.py:_windowed_ll} & \code{"v1"} (D-107, D-108)\\\midrule
censored at $t_{\rm hi}$ & $t_i>t_{\rm hi}$ & $t_i>t_{\rm hi}$, or a death with $t_i+\delta_i>t_{\rm hi}$ (D-107)\\
death with $t_i\le t_{\rm hi}<t_i+\delta_i$ & keeps its whole interval & censored at $t_{\rm hi}$\\
conditioning on $t_{\rm lo}$ & $\log\mathrm{clip}(S_{\rm obs}(t_{\rm lo}),10^{-12},1)$ & $\log^{\ge}S_{\rm obs}(t_{\rm lo})$ (D-108)\\\bottomrule
\end{tabular}
\end{center}
\src{engine/observation.py:Assembly.__init__, assemble}.

\paragraph{Heterogeneity.} One parameter $P\in\{\eta,\beta,\eps,\kappa,\xc\}$ is $P_k=P(1+v\,z_k)$ over the
probabilists' Gauss--Hermite nodes $(z_k,w_k)$, $k\le M$, $\sum w_k$ normalised over the nodes with $P_k>0$
(dropped nodes are marched at scale 1 with weight 0) \src{engine/evaluator.py:_node_scales, _points}. The
engine marches all nodes (\code{vmap}) and \emph{averages the stage-1 points} before anything nonlinear:
\begin{equation}
  \bar S(\cdot)=\sum_kw_kS_k(\cdot),\quad\bar F=\sum_kw_kF_k,\quad\bar p=\sum_kw_kp_k,
  \quad\text{and stage 2 is applied once to }(\bar S,\bar F,\bar p).
\end{equation}
Conditioning therefore divides the mixture: for entry $\ell_i$ contains $-\log^{\ge}\bar S_{\rm obs}(e_i)$, not
a mixture of conditional terms; likewise $S_{\rm obs}(t_{\rm lo})$ for the window. The D-004 switch
$F(t+\delta)<S(t_s)$ is evaluated on $\bar F,\bar S$, i.e.\ after mixing. The population mean is that of
SRK's \code{fp_loglik_het_physical}; $v$ may be free (\code{het_v}); \code{m_report} gives each early-death
term's change between $M$ and $M+4$ nodes. D-116: a node with $1+vz_k\le0$ is dropped and the weights renormalised,
so the likelihood jumps in $v$ where a node crosses zero; the count is reported and no gradient in $v$ is given
there. D-117: $P(1+vZ)$ itself puts probability $\Phi(-1/v)>0$ on non-positive parameter values (the model half
of D-116). Both are ported as they are.

\paragraph{D-004: the death mass from $F$ under \texttt{"v1"}.} For a death with resolution, the v0.3.3 mass
\eqref{eq:fp-m033} is a difference of two numbers near $1$ and underflows to $0$ (or is round-off) when $S\approx1$:
the term is $-700$ or a rounding-dependent value near $-31\ldots-37$. Under \code{"v1"}, after any mixing,
\begin{equation}
  \begin{gathered}\text{use\_F}=\big[\,F(t_i+\delta_i)<S(t_s)\,\big],\qquad
  m_i^F=\max\!\big(F(t_i+\delta_i)-F(t_s),\,0\big),\\
  m_i^{\rm v1}=\begin{cases}m_i^F,&\text{use\_F}\\ m_i^{033},&\text{otherwise,}\end{cases}\end{gathered}
\end{equation}
because $S(t_s)-S(t_i+\delta_i)=F(t_i+\delta_i)-F(t_s)$ in exact arithmetic and the difference with the smaller
operands carries the smaller error. Where \code{use_F} is false the expression is v0.3.3's own, bit for bit.
With an external hazard, factoring $e^{-H(t_s)}$ out of Eq.~\eqref{eq:fp-m033},
\begin{equation}
  m_i^F=e^{-H(t_s)}\Big[\,\big(F(t_i+\delta_i)-F(t_s)\big)+S(t_i+\delta_i)\,\big(-\mathrm{expm1}(-\Delta H)\big)\Big],\quad
  \Delta H=H(t_i+\delta_i)-H(t_s),
\end{equation}
which has no cancellation either (the switch itself is on the hazard-free $S$ and $F$)
\src{engine/observation.py:assemble}. Besides the cancellation, the $F$ route removes the $S$ route's stiffness
error: $S$'s difference is a tridiagonal-solve residual away from the true mass, growing like $u\rho_\eps\Delta\tau N^2$
(Section~\ref{sec:fp-guard}). It does \emph{not} remove the grid's discretisation error: a rescued early-death term can
still be several nats from its converged value at the production grid (Section~\ref{sec:fp-floors}, D-105).

\subsection{Grid policies}\label{sec:fp-grid}

Every result records the policy's name and version and the grid it produced (\code{Grid.summary}). No environment
variable is read \src{engine/grid.py:GridPolicy}.
\begin{itemize}[nosep]
  \item \textbf{\code{"v033"}} (version 1): v0.3.3's production grid, by calling
    \code{compat.v033.estimator.solver_grid} on the cohort's arrays \src{engine/grid.py:v033_grid}. That is $L$ from the
    Kaplan--Meier mean, $(\tau_{\max},\Delta\tau)$ from Eq.~\eqref{eq:pv-autogrid}, the dt ceiling
    Eq.~\eqref{eq:pv-ceiling}, $N=900$, $G_{\rm tol}=10^{-11}$, $n_{\rm start}=20$. A per-row $\delta$ has no v0.3.3
    rule; the ceiling then uses $\min_i\delta_i$.
  \item \textbf{\code{"fixed"}}: explicit $N$, $\tau_{\max}$, $L$, $G_{\rm tol}$ and exactly one of a uniform
    \code{dt} or an increasing node list \code{nodes} (e.g.\ \code{graded_start}).
  \item \textbf{\code{"model"}}: for a parameter set without a cohort (\code{curves}), $L$ is the model's own mean
    lifetime, as SRK's \code{fp_curves.auto_L} \src{engine/curves.py:model_L, model_grid_policy}. A coarse probe march
    ($N=300$, $\Delta\tau=4\times10^{-3}$, $\tau_{\max}=20$, $G_{\rm tol}=10^{-6}$) is run with time unit $L$ (start $L=1$ or a
    given \code{t_ref}); if the probe has absorbed ($S_{\rm end}<10^{-3}$, finite positive mean) the result is
    \begin{equation}
      L\leftarrow L\cdot\int_0^{\tau_k}S_{\rm probe}\,\dd\tau\quad(\text{trapezoid, to the probe's stop index }k),
    \end{equation}
    otherwise $L\leftarrow4L$ and repeat (at most 12 tries, then refusal). The curves are then computed on the standard grid
    $N=900$, $\Delta\tau=10^{-3}$, $\tau_{\max}=20$, $G_{\rm tol}=10^{-11}$ with this $L$.
\end{itemize}
A graded mesh without segments uses the general march with one span \src{engine/evaluator.py:_setup_segments}.

\subsection{The CI guard (constitution VII)}\label{sec:fp-guard}

The default numerics may differ from \code{"v033"} only where a register entry's machine-checkable condition says so
\src{engine/guard.py:compare}. For one cohort and one parameter point both numerics are evaluated term by term
and every term is classified as \emph{covered} by an entry, or \emph{uncovered} (and then required to be equal in bits).

\paragraph{Predicates.} A fixed vocabulary joined by \code{&} in the register's \code{terms} field:
\code{death}, \code{censored}, \code{interval}, \code{exact}, \code{use_F}, \code{hazard}, \code{window_close}
(a death with $t\le t_{\rm hi}<t+\delta$), \code{lo_floor} ($S_{\rm obs}(t_{\rm lo})<10^{-12}$, for every at-risk row)
\src{engine/guard.py:_mask}. D-004: \code{death & interval & use_F}; D-106: \code{death & exact & use_F}; D-107:
\code{window_close}; D-108: \code{lo_floor}. D-107 and D-108 have bound ``any'' (the entry defines the new value).

\paragraph{Value bound, in mass space.} For covered rows of D-004 and D-106,
\begin{equation}
  \big|m^{\rm v1}-m^{033}\big|\le c\,u\,K+f\,u,\qquad K=\rho_\eps\,\Delta\tau\,N^2,\quad c=1.8996,\quad f=348.2,
  \label{eq:fp-guard}
\end{equation}
where $u=2^{-52}$ ($f$ in ulps), $K$ is the solve's largest diffusion number over heterogeneity nodes and age spans
(\code{Evaluator.stiffness}), and the mass is the death mass, the density times the largest step for exact times, and the
survival for a censored row \src{engine/guard.py:masses}. The constants are frozen from the guard pilot
\src{tests/acceptance/results/guard_pilot_002.json}: $c=4\times0.4749$ (margin $4\times$ the worst
$|\Delta m|/(uK)$ at the stiff points, $K\ge10^4$; Denmark\_F\_1900 hetro truth at $N=450$, $\Delta\tau=10^{-3}$) and
$f=4\times87.1$ ulps (the worst excess over $0.4749\,uK$ at ordinary points). The pilot ran $39$ ordinary and $11$ stiff
points. The form follows the mass: the $S$ route's error is the tridiagonal solve's residual, which grows with $K$.
\paragraph{Floor.} $m^{033}=0$ for a term on the $-700$ floor (and likewise $m^{\rm v1}$): a floored v0.3.3 term gets no
free pass; v1 may rescue it only to a mass the $S$ route could not resolve.

\paragraph{Term/mass consistency.} For covered death rows (excluding rows an ``any'' entry defines, and rows outside the
window), the term must be its own mass minus the same conditioning in both numerics:
\begin{equation}
  \big|\,(\ell^{\rm v1}-\log^{\ge}m^{\rm v1})-(\ell^{033}-\log^{\ge}m^{033})\,\big|\le10^{-9},
\end{equation}
so a change to the term that its mass does not show is not that entry's.

\paragraph{\texttt{use\_F} recomputed.} For interval deaths the engine's \code{use_F} is checked against the same rule
recomputed from v0.3.3's own $S$: $(1-S_{033}(t+\delta))<S_{033}(t_s)$. They may differ only within $\mathrm{BAND}=10^{-9}$
of the switch.

\paragraph{Gradients.} With $w$ a 0/1 weight vector (\code{Evaluator.weighted_grad}):
\begin{enumerate}[nosep]
  \item \emph{uncovered}: the weighted gradient of the uncovered terms agrees to $10^{-10}$ relative to
    $\max|g^{033}|$ (not bits: XLA fuses the two programs differently), in addition to the \emph{terms} being equal in bits;
  \item \emph{covered, resolved by both} ($|\ell^{\rm v1}-\ell^{033}|\le10^{-8}$, neither on the floor): agree to $10^{-6}$ relative;
  \item \emph{rescued} (covered death rows, not floored under v1, and either $|\Delta\ell|>10^{-8}$ or floored under v033): v1's
    weighted gradient against central differences of v1 in $z$, with steps $10^{-4},10^{-5},10^{-6},10^{-7}$ (best accepted), each
    with its own round-off bound $50\,u\,\max(|\ell|,1)/h_{\rm step}+10^{-6}\max|g|$.
\end{enumerate}
\paragraph{Planted changes.} The guard must fail on changes no entry covers (constitution V); four are planted by
patching \code{Assembly.assemble} for v1 \src{tools/acceptance_002.py}: $+10^{-12}$ on every censored term;
a factor $(1+2.2\times10^{-16})$ on every term; $+10^{-3}$ on every death term; and a \code{stop_gradient} on the $F$ route
(values unchanged, rescued gradient wrong; caught by the covered and rescued gradient checks). Each must fail and name its
terms.

\subsection{FLOORED\_DEATHS}\label{sec:fp-floors}

A sixth truth-free check of the fit-cohort calibrator \src{floors.py:report, message}. At $\hat\theta$ the deaths are
evaluated on the fit's own grid with numerics \code{"v033"} (this fit) and \code{"v1"} (corrected).
\paragraph{Candidates.} A death is a candidate if its v033 term is on the $-700$ floor, or $|\ell^{\rm v1}-\ell^{033}|>1$~nat,
or $m^{\rm v1}<10^{-12}$ (robust to how v0.3.3 happened to round: one ulp takes the ITP example's added age-0 death from $-700$ to $-30.5$,
D-105).
\paragraph{Refinement.} A candidate inside the march ($\tau_i<\tau_k$ on the fit's grid) is re-evaluated under \code{"v1"} on a cohort of the
candidates only, with grid policy \code{"fixed"}: $N$ fixed, $\Delta\tau_k=\Delta\tau/4^k$, $\tau_{\rm end}=\max_i(t_i+\delta_i)/L+4\Delta\tau_k$,
$k=1,2,\dots$ until every term moves by less than $1$~nat between two refinements, at most $6$ refinements and
$4\times10^9$ node-steps in all. The step, not $N$, sets an early death's term (at the ITP $\hat\theta$ the age-0 death is
$-36.8$ at the fit's grid and $-114.0$ at $\Delta\tau/4^6$, while $N=900,1800,3600$ agree to $0.2$ nat). A candidate whose refinement did not converge
is reported with \code{converged=False}.
\paragraph{Categories.} Each candidate gets one category from its refined value:
\begin{enumerate}[nosep]
  \item \textbf{horizon}: $\tau_i\ge\tau_k$, beyond the march's early stop on the fit's grid: the grid, not the data, floored it;
  \item \textbf{impossible}: refined mass $<10^{-12}$ (or corrected mass $<10^{-12}$ with the refined term within $1$~nat of the corrected one): the model at $\hat\theta$ leaves no room for it;
  \item \textbf{grid}: otherwise, refined term differs from the corrected by more than $1$~nat: the fit's time step is too coarse for this death;
  \item \textbf{cancellation}: otherwise: v0.3.3's arithmetic misscored it and the corrected value is resolved.
\end{enumerate}
\paragraph{Example.} The ITP example (\code{examples/results/itp_c2016_het3_control_female.json}) plus an added age-0 death:
$-30.5$ in this fit, $-36.8$ corrected, $-114.0$ refined; category \emph{impossible}. The corrected-numerics
gradient at $\hat\theta$ is kept in the report. Nothing here changes the fit, its claims or the other checks.

\subsection{Verification: what is measured where}\label{sec:fp-verif}

Measured values are quoted only from committed result files under \code{tests/acceptance/results/}. A placeholder
\todo{from results/\ldots} marks a figure that is not yet in a committed file.

\begin{center}\small
\begin{longtable}{@{}p{1.55cm}p{4.6cm}p{4.2cm}p{4.5cm}@{}}
\toprule
Phase & What is measured & Against & Result\\\midrule\endhead
T0 & compiled program text; value and gradient bytes, numerics \code{"v033"}, all features off, side by side in one process
  & spec 001's port (\code{compat/v033/fp_jax.py}); 40 cohorts $\times$ 6 points
  & \code{t0_002.json}: programs equal $40/40$, bytes equal $240/240$, status pass\\
guard & the CI guard (Section~\ref{sec:fp-guard}) on its frozen panel, then four planted changes
  & \code{"v033"}, same grid, same process
  & \code{guard_002.json}: $78/78$ evaluations pass; $4/4$ plants fail\\
guard-pilot & $c$ and $f$ of Eq.~\eqref{eq:fp-guard} from $\max|\Delta m|/(uK)$ over $50$ points (ordinary and stiff)
  & \code{"v033"}
  & \code{guard_pilot_002.json}: $c=1.8996$, $f=348.2$ ulps ($4\times$ margin)\\
T1 & totals $\le10^{-5}$ nats per evaluation, resolved terms $\le10^{-8}$ nats, $\max|\Delta S|\le10^{-8}$, gradient $\le10^{-6}$ relative to autodiff of a fixed-length march
  & SRK's frozen NumPy code (\code{fp_loglik}, \code{fp_loglik_het_physical}); \code{fp_piecewise} for segments
  & \code{t1_002.json}: $474/480$ pass. Totals within $6.4\times10^{-10}$ nats except one cohort; resolved terms $\le2.0\times10^{-10}$; $\max|\Delta S|\le4.7\times10^{-10}$; gradient within $9.2\times10^{-7}$ of the fixed-length autodiff on all $350$ uniform-mesh evaluations. Fails: the $n=93{,}893$ Swedish cohort at $\beta,\eps\times1.5$ (totals $1.3\times10^{-5}$--$1.0\times10^{-4}$; cause not isolated: at the same $K=4.76\times10^5$ its truth, $\eta\times1.5$ and $\eta\times10$ agree to $\le2.1\times10^{-6}$, and all its terms fall below the low-mass cut, so it is checked by its total only), and central differences at two piecewise-smooth points; not applicable at $59$ (cause not diagnosed). Segments' gradient is checked independently on $111$ of $130$ evaluations; the other $19$ rest on the self-check\\
mpref & numerics \code{"v1"} terms, low-mass rows and the round-off panel, within $10^{-5}$ nats; $\pm3$-ulp scan of $\theta$ moves the total by $\le10^{-5}$ nats
  & the same discrete scheme in mpmath, 60 digits (\code{tests/reference/mp_scheme.py})
  & \code{mpref_002.json}: F-route low-mass rows of $108$ evaluations within $7.4\times10^{-6}$ nats (v033 up to $209$). Round-off panel: v1 matches at every fitted or stored $\hat\theta$ (Sweden stored: $-4.9\times10^{-8}$ vs v033's $+83.25$); 5 of 17 fail (float64 conditioning): at three stiff human truths ($K\lesssim8\times10^6$) both numerics are $10^{-4}$--$3\times10^{-4}$ off, and on the $n=93{,}893$ Swedish cohort v1's $\pm3$-ulp spread is $1.7\times10^{-5}$ at $\eta\times1.5$ and it is $2.4\times10^{-5}$ off at $\eps\times1.5$\\
early & residual of each early-death term ($S_{\rm obs}>1-10^{-3}$) at the v033 grid, v1 and v033
  & a converged graded reference (start $\Delta\tau_0=10^{-10}$, ratio $1.005$, then $\Delta\tau=\delta/(10L)$) at $N=1800,3600,7200$, extrapolated in $N$
  & \code{early_002.json}: $1313/1322$ references converged. Of the other $9$, $4$ with masses $<10^{-30}$ move by up to $3.5$ nats; $5$ with masses $10^{-21}$--$2\times10^{-14}$ move by $1.2\times10^{-4}$--$2.2\times10^{-3}$ nats, just above the $10^{-4}$ tolerance. v1 residual sums: Denmark M~1890 $-12.3$ (max $0.07$), pure-k0.1-p22 $+3.19$ (max $0.32$), Sweden F~1910 stored $+190$ (max $134$), ITP $+$ age-0 $+75.4$ nats\\
anchor & (a) $S$ and early interval masses for $\eta=\beta=0$ at three $\rho_\eps$, refining $N$ and the step together; (b) self-convergence of the total log-likelihood to $N=7200$ (reported, not judged)
  & the pure-diffusion closed-form series (\code{tests/reference/diffusion.py}); the scheme's reflecting wall sits at $-h/2$ from the start cell, so $S$ converges at order $O(h)$
  & \code{anchor_002.json}: pass. $S$ observed order $0.999$ at all three $\rho_\eps$; early-term errors fall monotonically; the stiff Denmark truth is pre-asymptotic at $N=900$ (differences $0.15, 0.034, 0.048, 0.031$ nats)\\
speed & CPU per evaluation, 1 core, $\ge5$ replicates, 95\% CI of the ratio
  & v0.3.3's port (same program: $\le1.03$ upper CI); \code{"v1"} ($\le1.20$); features on vs the NumPy oracle ($\le1.0$)
  & \code{speed_002.json}: pass. Off configuration $0.9996$ $[0.996, 1.002]$ (value), $0.9995$ $[0.997, 1.003]$ (value $+$ gradient); v1 $1.025$ $[1.020, 1.031]$ and $1.033$ $[1.026, 1.042]$; features against NumPy $0.15$--$0.37$; per step $11.6\,\mu$s (value) and $30.0\,\mu$s (value $+$ gradient) against the cost model's $11$ and $29$--$40$\\
floors & FLOORED\_DEATHS firing set and categories
  & the set computed in advance from per-term values at the stored $\hat\theta$
  & \code{floors_002.json}: the quick tier and Sweden match the committed prediction. The ITP $+$ age-0 CLI fit does not: the age-0 death moves $\hat\theta$ ($\eta\ 3.6\times10^{-5}\to6.5\times10^{-6}$; BASIN warnings) and the death at 53 days becomes impossible too. Post hoc (\code{floors_posthoc_002.json}): the independent predictor at that $\hat\theta$ gives exactly the fit's rows\\
score & mean gradient at the true $\theta$ over $R$ simulated cohorts, per feature (pure control, hazard, window, entry, heterogeneity; segments and the clock moved to spec 007)
  & SRK's Euler/Brownian-bridge simulator (\code{SR\_hetro}), read-only, at $\Delta t/8$ (measured: at $\Delta t$ the pure control was biased by $\sim5$ SE); zero within its standard error at a multiplicity-corrected threshold
  & \code{score_002.json}, \code{score_window_002.json}: pure control $|z|\le0.58$ ($R=200$); hazard $|z|\le2.76$; entry $|z|\le0.6$ ($R=40$), its alternative detected at $z=-123$; window ($t_{\rm hi}$ at the median) $|z|\le1.06$, SRK's window close detected at $|z|\le11.1$. Heterogeneity in $\xc$ ($v=0.2$, $M=41$, as $M=11$ biased the $\eps$ score by 14--29 SE, \code{mixture_M.json}), \code{score_het_002.json}: $|z|\le0.69$ ($R=200$, grid certified) but node-by-node conditioning is not detected ($|z|\le2.58$; the pilot needs $R=481$--$3893$): not informative at this $R$, so the mixture's conditioning has no truth-anchored check here (spec 006)\\
\bottomrule
\end{longtable}
\end{center}
Byte identity holds only in the pinned environment on the host that runs the test. The identities (xc reduction, full-life
segment equals chronic value, $M=1$ equals homogeneous, window $[0,\infty)$ equals the full likelihood, entry at 0 equals no entry, terms sum to the
total) are unit tests, not acceptance phases.

\subsection{Known limitations}\label{sec:fp-limits}

From spec 002's \emph{Not Covered}:
\begin{itemize}[nosep]
  \item $N=900$ is pre-asymptotic at $n\sim10^5$: the grid error is not removed by D-004 (it removes the cancellation and the stiffness error of $S$, not the
    discretisation error), and no $N$ rule is a policy here.
  \item D-115: with an external hazard, a death beyond the march scores a hazard-only mass $S_k(e^{-h t}-e^{-h(t+\delta)})$ in both numerics (SRK's semantics). $S_k$ jumps with
    the discrete stop index, so the likelihood is jagged in $\theta$ near the horizon. Fixing it belongs with hazard fitting (spec 006).
  \item D-116 and D-117: the heterogeneity mixture drops non-positive nodes, so the likelihood jumps in $v$ where a node crosses zero (a quadrature
    artefact; no gradient in $v$ is given there); the definition $P(1+vZ)$ itself puts probability on non-positive values. The shared definition is spec 004's; only
    \code{normal_frac_sd} exists here, for one parameter at a time.
  \item The heterogeneity mixture's conditioning at entry or at a window's start has no truth-anchored check: the score test does not detect node-by-node conditioning at $R=200$ (the pilot needs $R=481$--$3893$). It is checked against its definition only \src{tests/unit/engine/test_mixture.py}; spec 006 carries the check.
  \item Float64 conditioning is untouched by D-004: at $\rho_\eta\ge10^6$, and at very stiff points ($K=\rho_\eps\Delta\tau N^2\gtrsim5\times10^5$ on $n\sim10^5$; the mpref and T1 rows above give the measured values).
  \item The time grid's own error on early-death terms is the largest error this spec measured, and D-004 does not remove it: v1 on the fit's grid against a converged reference is $+190$ nats on Sweden F~1910 at its stored $\hat\theta$ (at most $134$ on one term), $+75$ on the ITP example plus an age-0 death and $-12.3$ on Denmark M~1890 (\code{early_002.json}). A graded default mesh is spec 003's decision.
  \item Also inherited: the censor-tie bias (D-006) and the p15 grid barrier; both are grid or data properties that no policy here removes.
\end{itemize}

% Section for specs/003-calibrator (T046). Paths are relative to src/sraging/ unless they start with tests/ or tools/.
\section{The successor calibrator (\texttt{profile\_jax2}) and the claims layer}\label{sec:pj}

\subsection{Overview}\label{sec:pj-overview}

\code{profile_jax2} and \code{profile_jax2_h} 1.0.0 are the tester's successor to v0.3.3, \code{mle_profile_jax2}
and \code{mle_profile_jax2_h} at \code{sr_calibration_tester_jax2} commit \code{3bc5775}, ported verbatim into
\code{compat/jax2} \src{compat/jax2/VERBATIM.toml}. \code{profile_v033} stays the default; the switch is spec 003b's.
On top of them sits a claims layer \src{claims/}: a claim on each combination of constitution X (v1.1.0), new
shapes and labels for weakly determined quantities, and $L_{\rm SR}$.

\paragraph{What jax2 changes against v0.3.3.} Each difference is a register entry (D-201 to D-215) with its impact
source and a regression test \src{tests/unit/test_asis_defects.py}:
\begin{itemize}[nosep]
  \item fix B: an early death's mass from the absorbed mass $F$ (D-201; the same formula as D-004 without a
    hazard), and exact-time densities from $F$ (D-202);
  \item the basin re-fit: a profile node that beats $\hat\theta$ by more than $\chi^2_1(0.95)/2$ restarts the MLE
    (D-203); the Nelder--Mead fallback is v0.3.3's (D-204);
  \item the censor-tie relabel, which fires only when every censored record shares one age, the maximum (D-205; so
    D-006's bias stays live elsewhere, and CENSOR\_TIE still warns there);
  \item self-consistency widening: each claim becomes its union with the evaluated points within half the threshold
    of the best (D-206);
  \item the N rule ($N=900\cdot2^k$, cap 3600; D-207), a graded early mesh (D-208), a floored $\Delta t$ ceiling
    (D-209), and an exception policy fixed to Phase C's \code{SRB_STRICT=1} (D-210);
  \item \code{_h}: a constant external hazard fitted as a nuisance, never claimed (D-211).
\end{itemize}
Where jax2's arithmetic differs from the spec 002 engine (the hazard bracket, $-700$ past the march's stop, the mesh
start and the window floor) the difference is recorded, not reconciled (D-212 to D-214).

\paragraph{What 1.0.0's claims are, and are not.} The benchmark's ratification of 1.0.0 (2026-10-08) lists items a
port must state; the port changes none of them, and the README states them to users (Naveh, 2026-10-08; numbers,
sources and sha256 in \src{specs/003-calibrator/review/evidence/jax2_reference_notes.json}):
\begin{itemize}[nosep]
  \item the 95\,\% intervals under-cover with the model right: on valid pure-arm cohorts the closed $\eta$ intervals
    covered the truth 0.892 of the time (415/465) and $\eps$ 0.909 (431/474), against v0.3.3's 0.884 and 0.890;
  \item their ends run narrow: on a selected set of 17 cohorts every end 1.0.1 moved was too narrow under 1.0.0, by
    up to 0.11 dex;
  \item $\kappa$'s bounds are not calibrated: one-sided $\kappa$ bounds covered the truth 0.402 of the time (190/473;
    v0.3.3 0.409), nominal 0.95;
  \item a human $\eta$ or $\eps$ upper end set at $\beta$'s box edge ($10^3$) is box-conditional and flips open or
    closed across machines (27 such records under \code{profile_jax2}, 19 under \code{_h}); claims differ between
    the Mac and the cluster on 17 of 17 cohorts compared, and 3 of 4 human records change shape;
  \item the censor-tie relabel (D-205) changes how the data are read; it is what lifts $\eta$'s closed coverage on
    the censor arm from 0.634 to 0.912, and the fit's \code{mechanisms} line counts the rows it relabelled;
  \item \code{profile_jax2} is misspecified on external deaths and on a fit window: use \code{_h} for external
    deaths; on the window arm its closed $\eta$ coverage is 0.533 (v0.3.3 0.528). The benchmark's \code{_h} models a
    window; sraging refuses windows;
  \item first-bin deaths are recorded at age 0, not $\delta/2$ (AGE0\_DEATHS checks it; the benchmark's library
    records them so from v5).
\end{itemize}

\subsection{The port and its proof}\label{sec:pj-port}

\code{tools/port_jax2.py} copies source text out of the frozen snapshot \src{tests/oracle/jax2_3bc5775} by AST
position and declares every edit; \code{tests/unit/test_verbatim.py} applies the declared edits to the snapshot's
AST and requires equality with the port, and fails on planted one-token changes. The edits: import rewrites to
\code{compat.v033}'s hash-equal SRK subsets and a new Makeham subset; \code{PJRT_NPROC} made explicit; the
\code{SRB_STRICT} read fixed to Phase C's value; the tester's deployment guards (sibling loader, version lock,
\code{srk_sha}) deleted, their job done by the verbatim test and the snapshot's \code{VERSION_LOCK.json} check.
The registered names' class constants are ported as data with their inheritance \src{compat/jax2/settings.py}.

\paragraph{The archive, without editing the port.} The estimator's objective and its memo are local to \code{fit}.
\code{calibrators/profile_jax2.py:capture_objectives} subclasses \code{_Objective} for the duration of a fit to
record each instance; the archive is read with the port's own \code{_sc_archive}. Each profile's final nodes are
captured the same way, by wrapping \code{_repair} (output unchanged); jax2 stores a node's free coordinates only,
so the capture rebuilds the full $z$ as the port's \code{_full_points} does, and since jax2 re-runs \code{_repair}
after any polish improvement the last capture is final. A slow test fits the same cohort with and
without the captures and requires identical results, claims and widening records (timing leaves excluded as the
tester's own identity comparison excludes them).

\subsection{Combination claims}\label{sec:pj-combos}

The combinations, with $\bar r=\beta\bigl(1-\kappa\ln(1+1/\kappa)\bigr)$ (which already contains $\beta$), are
$\eta/\eps$, $\bar r/\eta$, $\sqrt\eta/\eps$, $\beta/(\eta(\kappa+1))$, $\bar r$ and $\beta/(1+\kappa)$ \src{claims/combinations.py}.
Their order is frozen by measurement (FR-016): the lead group is the combinations invariant on the scaling ridge
$(\eta,\beta,\eps)\to\lambda(\eta,\beta,\eps)$ by construction whose median spread along the stored kfix
$\kappa$-slices is $\le0.05$ dex on the validity-valid panel; the rest follow by that spread
\src{tools/acceptance_003.py:phase_order}.

\paragraph{The walker} \src{claims/walker.py}. Each node maximises the port's log-likelihood subject to
$\log g(z)=c$ inside the original rectangular box: SLSQP with the box as bounds and one equality constraint, the
port's \code{nll_grad} for values and gradients, so every point joins the fit's archive. Four of the six
combinations make the box a value-dependent interval on one coordinate; $\bar r/\eta$ and $\beta/(\eta(\kappa+1))$
couple $\eta$ and $\kappa$, which no box expresses, hence the constrained solver. The walk copies jax2's defences
against false caps:
\begin{itemize}[nosep]
  \item the centre node is re-solved at $g(\hat\theta)$ from every start, so a stale $\hat\theta$ does not anchor it;
  \item its step rule (0.05 dex first, $\times3$ or $\times2$, at most 2 dex) and jump rejection (4 nats); at the
    minimum step the jumped node is kept, so the profile crosses at a cliff rather than reading open;
  \item stops at deviance 10, at the box image, after 24 nodes, and, on the lower side of a combination containing
    $\beta$, once a node's $\beta$ sits at its floor: the combination is then consistent with no removal, and going
    lower only forces $\eta$ up or $\kappa$ to its cap, so the claim is $[0,\text{cap}]$; and, past the threshold,
    once its last 3 nodes span less than 0.01 dex (``stalled''; Naveh, 2026-10-08): the crossing is then bracketed
    at SC-004's resolution. The step rule can crawl when solves past the edge keep failing (each rejection divides
    the step by 3); on the ITP example the lower side of $\bar r/\eta$ ran to the 24-node limit (at \code{2306641}),
    and the stop cut the lead pair from 49,825 to 26,726 evaluations with the claims unchanged
    \src{examples/results/itp_c2016_het3_control_female_jax2.json}. jax2's own \code{_walk_side} crawls the
    same way (272 of 7,846 final-node profiles have $\ge8$ nodes within 0.01 dex) and is not changed upstream;
  \item mid-walk starts from the previous node, $\hat\theta$, and, for the combinations the scaling ridge moves
    ($\sqrt\eta/\eps$, $\bar r$, $\beta/(1+\kappa)$), the previous node carried along the ridge by the step;
  \item jax2's edge polish: the first node at or past the threshold is re-solved from every start (the seed, the
    $\kappa$ corners of $\hat\theta$ and of the node, the h starts, the neighbours, ridge starts at $\pm1$ dex, archive
    points near the node); if it falls inside, the polish moves outward (at most 12 nodes) and walks on once past the
    last node; the node just inside each edge is re-solved too (\code{profile_ci} interpolates between the two);
  \item monotonicity repair.
\end{itemize}
Walker step 1's archive check found 6 of 183 ends contradicted by points the fit and the walks had already evaluated
(at \code{6ceb692}); refits traced them to three mechanisms, and four defences were added (Naveh, 2026-10-06):
\begin{itemize}[nosep]
  \item for $\eta/\eps$ and $\sqrt\eta/\eps$, whose constraint leaves $\beta$ free, the centre and every polished node
    also start from $\beta$ spread over its box (6 values): three $\sqrt\eta/\eps$ lower ends were 1.1--1.2 dex too
    narrow, the inside set on a branch at $\beta\approx10$--$17$, $\kappa\approx160$--$300$, while $\hat\beta$ was
    $10^{-9}$--$10^{-8}$ (deviance 0.07--0.16 just 0.01 dex past the end);
  \item before the lower end of a combination containing $\beta$ is closed, the best fit with $\beta$ at its floor is
    found; inside the threshold, it joins the profile and the side reads open: an $\bar r$ walk had stopped at a solver
    cliff at 0.041 while $\beta$'s floor sat at deviance 1.7;
  \item each end's crossing is bisected to 0.01 dex: two ends had been interpolated 0.033 and 0.16 dex short;
  \item a final pass re-walks any side an archived point contradicts (beyond an end by more than 0.01 dex inside the
    threshold, or above the reference), from the contradicting points, while the archive grows (at most 3 rounds);
    the lead pair is reconciled before the other four walk, so \code{"all"} still leaves the default claims as they
    are. Every walk starts from the archive of every captured objective.
\end{itemize}
The archive check is then this pass's convergence check, and the brute force (walker 2) the independent test.
One deviance reference serves the polish, the repair and the interval: the profile's own maximum, at least the
fit's log-likelihood. The interval is the port's \code{profile_ci}.

\paragraph{Which combinations a fit computes.} By default the frozen lead group, $\eta/\eps$ and $\bar r/\eta$;
\code{combinations="all"} (\code{--combinations all}) adds the other four, and \code{False} none (Naveh,
2026-10-06). Walker step 1 measured the cost on the 20 walker cohorts, with the defences above and the stall stop
(\code{a9727f4}): all six walks a median $7.5\times$ the fit's own objective evaluations ($6.6\times$ its CPU, range
$2.5$--$13.1\times$), the lead group $2.3\times$ (range $0.61$--$4.25\times$; Naveh accepted it above the planned $2\times$,
2026-10-07); before the defences ($6.4\times$ and $1.8\times$) the walks were cheaper and wrong on 6 of 183 ends
\src{tests/acceptance/results/walker_step1_003.json}. The lead walks run first, and the lead claims are widened over
the archive as it stands when they end, so \code{"all"} adds four claims and never changes the default two; a later
walk's point beyond a lead end is reported (\code{combination_claims_contradicted}), as the walks' points are for the
per-parameter claims. All six are truth-scored before any ships.

\paragraph{Widening and the claim.} The claim is the port's one-quantity mapping (\code{h_claim_kwargs}) with
$g$'s box image as its box; it is widened by the port's own \code{sc_extents} and \code{sc_widen_claim} over the
whole archive with a $\log g$ column appended \src{claims/widening.py}. If a walk finds a point more than 1.92 nats
above $\hat\theta$, $\hat\theta$ and the per-parameter claims stay as the port gives them; the point joins the
archive and BASIN\_GAIN fires on the archive gap (FR-017).

\subsection{Shapes, labels and $L_{\rm SR}$}\label{sec:pj-shapes}

\src{claims/shapes.py}, \src{claims/lsr.py}.
\begin{itemize}[nosep]
  \item \textbf{open above: box-limited}: at the node or archive point that sets an upper endpoint, some
    coordinate (nuisances included) lies within 0.01 dex of its \emph{upper} box edge; scored apart from
    data-determined bounds. Lower edges never count: $\beta$, $\kappa$ and $h$ at their floors are the model's own
    limits (no removal, saturated removal, no external hazard), which a wider box cannot move;
  \item \textbf{negligible}: an upper bound on $\bar r/\eta$, in $L_{\rm SR}$ units, below its regime's floor
    $c/\sqrt{n_{\rm eff}}$ ($n_{\rm eff}\le2000$); the coefficients are the \S8b sweep's measured regime values
    converted by $f(\kappa)$ \src{tools/acceptance_003.py:phase_floor}. The regime is chosen by
    $\rho_\eps=\eps L_{\rm SR}$ alone ($\kappa$ is not identified when removal is negligible, and $c_{\bar r}$ barely
    depends on it: mid 1.535, $\kappa$-large 1.556): the smallest coefficient among the regimes nearest any $\rho_\eps$
    in the $\eps$ claim's range; refused below $\rho_\eps=0.1$, where no floor is measured, and never on a widened
    upper side. A label on an upper bound, never a shape of its own;
  \item \textbf{the $\kappa$ box on raw $\beta$} \src{labels.py}: every raw-$\beta$ upper end, of an upper bound or
    a closed claim, reads ``set by the kappa box ($\kappa_{\max}=300$), not by the data''. The $\kappa_{\max}=3000$
    study refitted 10 cohorts whose raw-$\beta$ cap sat at $\kappa_{\max}$ and 3 closed negative controls: every cap
    and every control's upper end moved by about 1 dex, the controls' lower ends by at most 0.0014 dex, the caps on
    $\bar r$ and $\bar r/\eta$ by at most 0.064 \src{tests/acceptance/results/kmax_003.json};
  \item \textbf{saturated}: a $\kappa$ cap below $0.364/\sqrt{n_{\rm eff}}$;
  \item \textbf{widened}: ``at least $X$; edge not found'', with the unwidened claim kept;
  \item $L_{\rm SR}=\int_0^\infty S\,dt$ at $\hat\theta$ from the engine's curves, refined to $10^{-4}$ relative; the
    box and each combination's box image are also reported in $L_{\rm SR}$ units (FR-021; the raw box's unit
    dependence is D-215).
\end{itemize}

\subsection{Checks for jax2 fits}\label{sec:pj-checks}

FLOORED\_DEATHS evaluates terms with the fit's own likelihood and grid (\code{floors.report_jax2}): candidates are
deaths with mass below $10^{-12}$ or a term on the floor, refined by $\Delta t/4^k$, and, as a second group with its
own refinement budget, every death whose fit-grid mass is below $10^{-6}$, since the category is defined on the
refined mass; categories horizon, impossible and grid (no ``cancellation'' under fix B); each listed death shows its
share of the log-likelihood, which is not its leverage: the leverage check is not run for jax2 fits until spec 003b. BASIN\_GAIN reads
the archive gap; CENSOR\_TIE is quiet where the relabel applied and says why it did not elsewhere.

\paragraph{A death the model calls impossible is fitted in, not floored.} Under fix B such a death has a finite term
that pulls the fit. Added to the ITP example, one age-0 death moves $\hat\eta$, $\hat\beta$ and $\hat\eps$ by 1.38,
1.64 and 1.58 dex until its mass is $2.3\times10^{-6}$, at a cost of 5.4 nats to the other rows; at that $\hat\theta$
FLOORED\_DEATHS has nothing impossible to list. At the $\hat\theta$ fitted without it, the death is impossible (mass
$4.6\times10^{-37}$, term $-83.7$ nats) \src{tests/acceptance/results/checks_003.json}. AGE0\_DEATHS, from the data
alone, is what fires; no check measures the pull (the leverage check, spec 003b).

\subsection{Verification}\label{sec:pj-verify}

Each phase runs on a clean committed tree in a worktree, priced first; its pass rule is a function of the committed
result file alone (\code{tests/acceptance/test_results_003.py}). Measured values are quoted only from those files.

\begin{center}\small
\begin{longtable}{@{}p{1.55cm}p{4.6cm}p{4.2cm}p{4.5cm}@{}}
\toprule
Phase & What is measured & Against & Result\\\midrule\endhead
pilot & the panels and their prices (bytes, fits, walker, kmax, coverage, speed)
  & ---
  & \code{pilot_003.json}, \code{panel_003.toml} (at \code{17de2ec}); nothing passes or fails\\
bytes & value, value $+$ gradient and grid bytes of both names, side by side in one process (SC-001)
  & the frozen jax2 code \src{tests/oracle/jax2_3bc5775}; spec 001's 40 cohorts $\times$ 6 points, C.~elegans points, the floored $\Delta t$ ceiling binding
  & \code{bytes_003.json}: pass; \code{jax2} $244/244$, \code{jax2_h} $488/488$ (two $h$ points)\\
fits & $\hat\theta$, claims, log-likelihood and diagnostics of whole fits, and the mechanisms firing (SC-002)
  & the frozen code on the quick tier and the named mechanism cohorts
  & \code{fits_003.json}: $\hat\theta$ bytes, results, claims and widening identical on $66/66$ (33 cohorts $\times$ both names), mechanisms $6/6$; the tester's stored Mac fits equal on $18/18$ with zero diffs (\code{fits_stored_003.json}, R3d6/L5's NaN strings read as floats)\\
adapter & the tester's scoring of the oracle's quick-tier records (SC-003)
  & the tester's own functions
  & \code{adapter_003.json}: pass; the port's ClaimSets valid and equal to the frozen method's, claim for claim, on the quick tier under both names; the tester's scores equal on every parameter and shape\\
checks & AGE0\_DEATHS and FLOORED\_DEATHS on ITP $+$ age-0; CENSOR\_TIE on staggered censoring (SC-009, amended)
  & the specified firing
  & \code{checks_003.json}: pass. The age-0 death is impossible at the $\hat\theta$ fitted without it (mass $4.6\times10^{-37}$, $-83.7$ nats, 4.1\,\% of the log-likelihood) and absorbed by the fit with it (1.38--1.64 dex); CENSOR\_TIE fires on the recorded construction (no library cohort qualifies, 0 of 1,990)\\
order & spread of each combination along the stored kfix $\kappa$-slices, validity-valid panel (FR-016)
  & the pre-registered hypothesis order
  & \code{order_003.json}: lead $\eta/\eps$, $\bar r/\eta$ (0.003, 0.026 dex), then $\sqrt\eta/\eps$, $\bar r$, $\beta/(\eta(\kappa+1))$, $\beta/(1+\kappa)$; frozen in \code{claims/order.json}\\
floor & negligibility coefficients on $\bar r/(\eta L_{\rm SR})$ by regime, from the stored \S8b sweep (cited by md5)
  & ---
  & \code{neg_floor_003.json}: $c_{\bar r}$ sharp 0.790, mid 1.532, $\kappa$-large 1.555, diffuse 7.54 (a lower bound: censored)\\
walker 1 & cost of the walks; the free archive check (no archived point $>0.01$ dex beyond an unwidened end at deviance $<3.84$)
  & the fit's own evaluations and archive; 20 cohorts
  & \code{walker_step1_003.json} (at \code{a9727f4}): all six $7.54\times$ the fit's evaluations ($6.63\times$ CPU), the lead group $2.30\times$ (accepted, 2026-10-07); archive check $182/182$. The first run (\code{6ceb692}: $6.4\times$, $1.8\times$, $177/183$) led to the defences above\\
walker 2 & each stated end against a brute-force constrained maximisation (SC-004: $\le0.01$ dex narrower, $\le0.05$ dex wider)
  & Sobol starts in the $\pm1.5$ dex cube and over the whole $\beta$ and $\kappa$ box, the scaling ridge, the $\kappa$ corners
  & \code{walker_003.json} (at \code{01799ed}): pass, $182/182$ ends and $182/182$ archive checks; the inward point capped at the profile's maximum on 27 narrow ends (Naveh, 2026-10-08); 92\,\% of the starts converged. The thinnest margins are on one cohort whose profiles are flat at the threshold (outward deviance 3.843--3.855, inward 3.72--3.84)\\
kmax & raw-$\beta$ caps refitted at $\kappa_{\max}=3000$ (SC-007, FR-025)
  & the same cohorts at $\kappa_{\max}=300$; 3 closed negative controls
  & \code{kmax_003.json}: box-limited, $10/10$ caps moved 0.998--1.013 dex; controls' upper ends 0.999 dex; $\bar r$ and $\bar r/\eta$ caps $\le0.064$ dex\\
coverage & paired, cluster-robust coverage gap of each combination against the same fits' $\eta$ and $\eps$ claims (SC-005)
  & the truth, medium tier
  & \code{coverage_003.json} (at \code{d903979}; 215 fits on the cluster's AMD EPYC nodes, scored here): \textbf{fail}, recorded and attributed (Naveh, 2026-10-10; \code{decisions_003.toml [coverage_result]}). On censor, $\eta/\eps$ loses to $\eta$ ($z=-2.70$, 35 cells, 14 clusters) and to $\eps$ ($z=-2.74$, 38, 13): closed coverage $36/44$ (Wilson 0.68--0.90), every miss above the upper end. It is $14/15$ on the tester's valid cohorts and falls with its CV-bias statistic: $14/15$, $11/13$, $6/7$, $5/9$ for $2n(\Delta\ln{\rm CV})^2<0.1$, $0.1$--$1$, $1$--$10$, $\ge10$ (\code{coverage_cv_bias_003.json}). 23 of the 36 gaps are underpowered (mostly one-sided claims); 11 pass, both lead combinations on pure among them. On hazard (\code{profile_jax2_h}), $\eta/\eps$ nearly loses: to $\eta$ $z=-2.13$ on 29 cells, one under the floor of 30, so underpowered; to $\eps$ $z=-1.92$, the gap's 95\,\% interval reaching only $+0.002$. Its paired gap is negative on all three arms (pure $z=-1.02$, $-0.43$). Mac against cluster on 15 cohorts: no hit flips (\code{coverage_xhost_003.json})\\
speed & port against the frozen code, 5 cohorts $\times$ 3 replicates (SC-010)
  & the frozen code, alternating order; one process; 1-minute load $\le5$ at the start, $\le6$ after
  & \code{speed_003.json} (at \code{38afa67}): pass; CPU ratio 0.9995, 95\,\% CI 0.996--1.003; cohorts 0.993--1.006 at 142--313 CPU-s per fit; loads 2.33 at the start, 3.23--4.31 after. The bar was 1.0 until the run; this Mac's load read 1.5--7.1 with its CPU about 90\,\% idle, so it is spec 002's (Naveh, 2026-10-10; \code{decisions_003.toml [speed_load]})\\
unchanged & compat/v033, the engine and the oracles untouched; spec 001's claims re-run (SC-011)
  & \code{main}; spec 001's oracle
  & \code{claims-001_003.json} (at \code{d903979}): pass. No untouched path changed against \code{main}; 22 files of the fit path's import closure did, so spec 001's claims were re-run: values, gradients, arrays, $\hat\theta$, claims and warnings identical on $40/40$, roles reconfirmed $10/10$, the tester's tree unchanged\\
\bottomrule
\end{longtable}
\end{center}

\paragraph{The coverage fail.} The one loss, $\eta/\eps$ on censored cohorts, is attributed to a defect of the
panel; the measurements support this but do not prove it. The tester's validity audit, fixed before any coverage was
looked at, found that its generator over-spreads lifespans: at the truth the generated lifespan CV exceeds the model's
(median $\Delta\ln{\rm CV}=0.086$, up to 0.17, on sharp cohorts). Data more spread out than the model pull
$\hat\eta/\hat\eps$ down, which is the side of every miss, and the misses grow with the bias. The direction is the
stronger evidence: on flagged cohorts the $\eta/\eps$ point lies below the truth in 60, 79 and 80\,\% of fits
(pure, censor, hazard; median $\log_{10}$ point/truth $-0.009$, $-0.015$, $-0.014$), on valid ones in 38, 47 and
60\,\% ($+0.009$, $+0.002$, $-0.004$). Coverage alone separates them weakly: on censor valid $14/15$ against flagged
$22/29$ (Fisher $p=0.23$), over the three arms $44/46$ against $77/89$ ($p=0.14$); $14/15$ has Wilson interval
0.70--0.99. The walker's precision is not excluded: 9 of the 14 $\eta/\eps$ misses, 6 of the 8 on censor, lie within
0.01\,dex of the end, inside SC-004's tolerance, and walker step 2 checked no censored cohort. On its 20 $\eta/\eps$
upper ends (pure, hazard, real) the deviance 0.01\,dex beyond the end was 4.8--18.0, well above 3.84, so the walker
was not short there; the next spec checks the 8 censor misses' ends by brute force before any re-test (Naveh,
2026-10-10). All figures: \code{coverage_cv_bias_003.json}. $\eta/\eps$'s intervals are the narrowest here (median 0.25 dex on censor, $\eta$'s
0.87), so the mismatch shows there while the wider $\eta$ and $\eps$ intervals still cover; the paired gate does not
cancel it, as \code{[coverage_validity]} had assumed. The order stays as measured (Naveh, 2026-10-10). The re-test
needs a panel without the bias, which no library version through v5 provides; it is the next spec's.

% Section for specs/003f-leverage-check (T026; constitution X and XII). Paths are relative to src/sraging/ unless they
% start with tests/ or tools/.
\section{The leverage check on the jax2 calibrators (\texttt{sraging.leverage})}\label{sec:lev}

\subsection{Purpose}\label{sec:lev-purpose}

Constitution X requires that every fit ask whether its estimate rests on one individual. Spec 003f implements the
check for \code{profile_jax2} and \code{profile_jax2_h}; \code{profile_v033}'s check is spec 005's. It runs inside the
calibrator's \code{fit}, after the fit and after the floor report, on the fit's own likelihood and grid
\src{calibrators/profile_jax2.py:_make_fit}. It never changes $\hat\theta$, the claims or the likelihood: its result goes
to \code{diagnostics["leverage"]}, which the reference comparison does not read \src{result.py:SRAGING_DIAGNOSTICS}.

\subsection{The free statistics}\label{sec:lev-free}

With the per-individual terms $\ell_i(z)$ of the port's own likelihood at $z=\log\theta$ (the value march with
\code{with_terms=True}, on the prepared data and grid of \code{profile_jax2.fit_inputs}):
\begin{itemize}[nosep]
  \item gap/sd on the death times, flag when $g>2$ and $n_D\ge 20$ \src{leverage.py:gap_sd};
  \item the worst individual's share in the absolute form, $|\ell_w|/\sum_i|\ell_i|$ with $w=\arg\min_i\ell_i$
    \src{leverage.py:share}.
\end{itemize}
The share does not pick the individual: on the ITP cohort with one added age-0 death, that death carries 0.66\,\% of the
log-likelihood at its own $\hat\theta$, and the worst term belongs to another individual (pilot P1,
\code{tests/acceptance/results/p1_003f.json}).

\subsection{Influence}\label{sec:lev-influence}

Let $g_i=\partial\ell_i/\partial z$ (forward-mode differentiation of the terms) and $H$ the Hessian of the total
(central differences of the port's analytic gradient, step $10^{-4}$ in $z$; forward-over-forward agrees to
$\le 10^{-7}$ relative on P1). Coordinates within $10^{-6}$ (in $\log$) of a box bound are fixed. On the free coordinates,
$I=-H$; its eigen-directions with eigenvalue at or below $10^{-6}$ times the largest, including non-positive ones, are
cut. The cut sits above the Hessian's accuracy ($9\times10^{-8}$ relative); the $\kappa$ ridge makes $I$ near-singular,
and on ITP one eigenvalue is $-1.7\times10^{-3}$. The negative eigenvalues, the distance of each coordinate to its box
bound and the gradient at $\hat\theta$ are recorded. If no direction survives the cut, the result is
\code{not_assessable}. The one-step shift of dropping $i$ is
\begin{equation}
  \Delta z_i \approx -I^{+} g_i ,
\end{equation}
since $\sum_k g_k=0$ at $\hat\theta$. Its effect on a quantity $q$ is $\nabla q\cdot\Delta z_i$, for
$q\in\{\log\eta,\ \log\eps,\ \log(\eta/\eps),\ \log(\bar r/\eta)\}$ with $\bar r=\beta(1-\kappa\ln(1+1/\kappa))$
\src{leverage.py:influence, shifts}. The individual dropped is the one with the largest shift on the rule's quantities
(all four); a tie keeps the lowest row and is counted \src{leverage.py:ranking}. The one-step shift ranks; it does not measure: on ITP it predicts 0.028\,dex on $\eta/\eps$
against the refit's 0.20.

\subsection{The drop-one refit}\label{sec:lev-refit}

The cohort without the dropped individual is refitted twice.
\begin{itemize}[nosep]
  \item The calibrator's own MLE stage: the two steps of \code{compat.jax2.claims.run_fit}, with its exact keywords,
    plus \code{stages=("mle",)} \src{leverage.py:refit_mle}. A unit test compares the two calls. The MLE stage
    re-seeds and starts twice, the second start with $\eps\times4$.
  \item The estimator's own L-BFGS-B (\code{estimator._lbfgs}, with the MLE stage's options and fallbacks), started at
    the reported $\hat\theta$ on the same reduced objective (the warm start) \src{leverage.py:refit_warm}.
    \code{_lbfgs} never returns a point worse than its start.
\end{itemize}
Both points, and $\hat\theta$, are evaluated on the reduced cohort's objective by the value march, never from an
optimiser's stored value \src{leverage.py:reduced_loglik}. The MLE stage stands unless the warm start beats it by more
than $10^{-3}$ nats, the estimator's own optimiser noise; this keeps the calibrator's path wherever it is as good. The
point that stands is compared with the reported $\hat\theta$, the one the user gets, on every fit, including those whose
basin re-fit moved $\hat\theta$:
\begin{equation}
  R=\exp\max_{q}\bigl|q(\hat\theta_{(-j)})-q(\hat\theta)\bigr|,\qquad
  \Delta_j=\ell^{(-j)}(\hat\theta_{(-j)})-\ell^{(-j)}(\hat\theta)\ \text{(total nats)},
\end{equation}
with $q$ over $\log(\eta/\eps)$, $\log(\bar r/\eta)$, $\log\eta$ and $\log\eps$ \src{leverage.py:_drop}.

\subsection{The rule (Naveh, 2026-10-18; revised 2026-10-19; constitution X)}\label{sec:lev-rule}

\begin{itemize}[nosep]
  \item \code{failed} when $R$ or $\Delta_j$ is not finite, or $\Delta_j<-0.1$ nats: the refit stopped short of
    $\hat\theta$ on its own objective, so its move is not a measurement;
  \item \code{flag} when $R>1.5$ and $\Delta_j>1.92$ nats, the claims' one-parameter 95\,\% threshold (the joint
    four-parameter one is 4.74);
  \item \code{flat_valley_slide} when $R>1.5$ at a smaller gain;
  \item \code{no_flag} otherwise \src{leverage.py:label}, \src{leverage_rule.py:RULE};
  \item \code{not_assessed} (spec 003g, Naveh 2026-10-22; per individual since 2026-10-28, deviation DV-3) in place of
    \code{no_flag} or \code{flat_valley_slide} when the optimiser could not move from $\hat\theta$ (its warm start on the
    full data stopped where it began, $g_0=0$) at a free-gradient norm above $5\times10^{-5}$ per individual
    (\code{leverage_rule.py:G_CRIT_PER_IND}): no refit without one row can be compared with that point. The evidence
    does not say whether $\hat\theta$ is non-stationary or the objective is kinked there. The block records L-BFGS-B's
    termination (scipy's integer status, iterations, message) and the gradient per individual
    \src{leverage.py:not_assessed}, \src{leverage.py:warm_termination}. Measured per individual: $5.9\times10^{-6}$ at
    ITP+age0 and $1.9\times10^{-6}$ at spec 003g's fit 1 (assessed); $4.0\times10^{-4}$ to $2.6\times10^{-2}$ at the seven
    human hazard pilot cohorts and $2.8\times10^{-2}$ at fit 2 (not assessed). The registered criterion, a summed norm
    above 1, grew with $n$ and would have read the same pathology as \code{no_flag} at $n=1{,}000$. scipy's status is 2
    (ABNORMAL) on all eight pilot fits, the stationary one included, so the termination does not discriminate
    (\code{tests/acceptance/results/lev_pilot_003g.json}).
\end{itemize}
A flag, a failed refit, a fit not assessable, or a check that did not run raises the warning \code{LEVERAGE}
\src{warnings.py:leverage_warnings}.

The rule of 2026-10-18 read the lead pair ($\eta/\eps$, $\bar r/\eta$) with $10^{-3}$ nats per individual, and compared
the MLE stage with the MLE stage where the basin re-fit had fired. The sr-naveh review and a re-score of pilot P3's
stored drops changed it (\code{specs/003f-leverage-check/review/}):
\begin{itemize}[nosep]
  \item Nine MLE-stage refits ended 0.1--38 nats below $\hat\theta$, all on hazard cohorts with $n\ge25{,}000$, and had
    been labelled \code{no_flag}.
  \item On the 12 cohorts with $n\le10{,}000$, $\bar r/\eta$ crossed $1.5\times$ on 12 of 24 drops at $\le0.17$ nats in
    total: removal negligible on both sides, or a slide along a flat direction. Raw $\eta$ and $\eps$ never exceeded
    $1.40\times$ there.
  \item The library's negligibility floor cannot repair $\bar r/\eta$: it is not computable below $\rho_\eps=0.1$, which
    covers 10 of those 12.
  \item A per-individual guard is $10^{-3}n$ nats in total: 0.3 nats at $n=304$, and costless slides become flags below
    $n\approx170$.
\end{itemize}

\subsection{Evidence (the rule of 2026-10-19, all phases re-run at the final code)}\label{sec:lev-evidence}

\begin{itemize}[nosep]
  \item \textbf{ITP with an added age-0 death} (\code{itp_003f.json}).
    \begin{itemize}[nosep]
      \item The check drops that death and flags: $\eta$ moves $24\times$ and $\eps$ $38\times$, $\eta/\eps$ $1.59\times$,
        and the refit gains 5.42 nats.
      \item The MLE stage stood, and equals the stored fit without the death bit for bit; the warm start reached the
        same point.
      \item $\hat\theta$ and the log-likelihood are bit-equal with and without the check.
      \item The check cost 74 CPU-s, 4.9\,\% of the fit.
    \end{itemize}
  \item \textbf{Pilot P3}, 26 library v6 cohorts with $n$ from 1,000 to 100,000 (\code{p3_003f.json},
    \code{panel_003f.json}).
    \begin{itemize}[nosep]
      \item The median-influence drops flagged 0 of 26 (Wilson 95\,\% upper bound 0.13). Only the 12 cohorts with
        $n\le10{,}000$ are an informative null: 0 of 12, upper bound 0.24.
      \item No refit failed. The warm start stood on 11 of 52 drops. On 7 of them, on five human hazard cohorts with
        $n\ge25{,}000$, the MLE stage had ended 0.72--12.3 nats below $\hat\theta$ and would have failed; there every
        ratio was at most $1.0005\times$. On the other 4 the warm start was better by 0.002--0.14 nats.
      \item Raw $\eta$ and $\eps$ never passed $1.5\times$. The 9 slides are all $\bar r/\eta$ moves on cohorts with
        $n\le10{,}000$, at gains of at most 0.17 nats.
      \item The overhead was a median 8.7\,\% and a p90 of 19.6\,\% of the fit.
    \end{itemize}
  \item \textbf{The README example} (\code{readme_example_003f.json}). The check drops the death at day 53 in the CSV
    (recorded day 54). $\eta$ and $\eps$ move about $2\times$ at a gain of 0.51 nats: a slide, inside both intervals.
\end{itemize}

\subsection{Limits}\label{sec:lev-limits}

\begin{itemize}[nosep]
  \item The 1.5 threshold is a convention, measured but not calibrated. The positive evidence is one cohort ($N=1$).
    The held-out null, the small-$n$ null ($n\approx100$--300) and a planted panel are spec 003b's re-acceptance
    (AC11).
  \item Influence is a first-order ranking: on ITP it predicts 0.028\,dex on $\eta/\eps$ against the refit's 0.20.
  \item Drop-one cannot see a cluster: dropping one of $k$ similar deaths moves nothing while the others remain
    (masking). \code{AGE0_DEATHS} fires on 10 of P3's 14 human cohorts; the README example's three early deaths act
    together.
  \item $\hat\theta$ is re-evaluated at $\log\hat\theta$, not at the fit's exact $z$. On box-corner human cohorts the
    likelihood can move about 1 nat under 1--4\,ulp. That figure is the jax2 session's (sesync note of 2026-10-19,
    from its C1 record); it is NOT RE-MEASURED here. $\Delta_j$ there may carry that term.
  \item $\Delta_j$ also carries $\hat\theta$'s own position against the box.
    \begin{itemize}[nosep]
      \item On all seven human hazard cohorts of P3, fitted without a hazard term, $\beta$ sits 0.000--0.009 (in
        $\log$) below its upper bound with a free-gradient norm of 10--640. The optimiser could not move from
        $\hat\theta$ there (DV-3's wording: a kink fits the evidence as well as a non-stationary point), and the
        one-step shift's $\sum_k g_k=0$ fails.
      \item Three drops on Sweden\_F\_1910 gained 1.87--3.09 nats with every ratio at most $1.013\times$; two of them
        exceed 1.92 nats, one by the MLE stage.
      \item On those fits the ratio, not the gain, decides. The gain of a warm start on the full data, beside
        $\Delta_j$, is spec 003b's measurement.
    \end{itemize}
  \item gap/sd is $0$ on cohorts whose top death ages tie on the grid (14 of P3's 26).
\end{itemize}

% Section for specs/003b-jax2-101-port (T040; research R16; constitution XII). Paths are relative to src/sraging/ unless
% they start with tests/, tools/ or specs/. Every figure comes from a 003b result file, named where it is used.
\section{jax2 1.0.1: what changed against 1.0.0 (\texttt{compat.jax2\_101})}\label{sec:j101}

\tplPurpose
Spec 003b ports \code{mle_profile_jax2} 1.0.1, frozen by the jax2 session at tag \code{mle_profile_jax2-1.0.1} (commit
6884fc8; handoff \code{reports/freeze101/HANDOFF.txt}, sha256 \code{a424d48f...}). The names \code{profile_jax2} and
\code{profile_jax2_h} now mean 1.0.1; \code{profile_jax2@1.0.0} selects the port of Section~\ref{sec:pj} unchanged
\src{registry.py:get}. The default calibrator stays \code{profile_v033}: whether 1.0.1 replaces it is spec 003g's
comparison with the benchmark's library-v6 reference. This section describes only what 1.0.1 changes; everything else is
Section~\ref{sec:pj}.

\tplMath
\textbf{The scoped time step (\code{DT_RULE}).} With $\mathrm{CV}$ the Kaplan--Meier coefficient of variation of the
prepared data, $\delta$ the recording resolution and $L$ the Kaplan--Meier mean, the binning-corrected CV is
\begin{equation}
  \mathrm{CV}_c=\sqrt{\max\!\bigl(0,\ \mathrm{CV}^2-(\delta/L)^2/12\bigr)}\qquad(\mathrm{CV}_c=\mathrm{CV}\text{ for exact times}),
\end{equation}
Sheppard's correction for recording on a $\delta$ grid \src{compat/jax2_101/estimator.py:_cv_c}. Where
$0<\mathrm{CV}$ and $\mathrm{CV}_c<0.05$ (strict), the march's step is capped,
$\Delta t\leftarrow\min(\Delta t,\ \mathrm{CV}/80)$, with the uncorrected CV in the cap
\src{compat/jax2_101/estimator.py:_grid}. Outside that scope the grid is 1.0.0's. The cap removes the defect 1.0.0 and
v0.3.3 share on sharp data: at $\mathrm{CV}\lesssim0.025$ and $\Delta t=0.002$, BDF2 overshoot and the positivity clip set
$S$ to 0 past the mean, and every later death scores the $-700$ floor (handoff \S6 issue 13; Section~\ref{sec:pj}).

\textbf{Three changes to the profile stage}, each on in both registered classes and off in 1.0.0
\src{compat/jax2_101/settings.py}:
\begin{itemize}[nosep]
  \item \emph{inner polish} (\code{INNER_POLISH_J}): the inner re-optimisation \code{_polish_inner}, which 1.0.0 ran only
    with the hazard, now runs on \code{profile_jax2} too \src{compat/jax2_101/estimator.py:_polish_inner};
  \item \emph{bracket bisection} (\code{BRACKET_BISECT}): a profile bracket that crosses the threshold over more than
    0.3\,dex is bisected, at most 6 times per side, so an end is not set by a coarse node
    \src{compat/jax2_101/estimator.py:_bisect_brackets};
  \item \emph{extent $z$} (\code{EXTENT_Z}): each self-consistency extent stores its point's exact log-parameters
    (\code{z_qmin}, \code{z_qmax}), a bit-exact key into the fit's memo; no claim changes.
\end{itemize}
\code{BOX_FLAG} exists and is off. The likelihood itself is unchanged: of the four runtime files, \code{fp_jax.py} and
\code{observation_model.py} are byte-identical to 1.0.0's; \code{estimator.py} and the adapter changed
(\code{tests/oracle/jax2_6884fc8/MANIFEST.json}).

\textbf{\code{diagnostics["estimate"]}} (schema 1) is built by the adapter from \code{run_fit}'s result alone, with no
likelihood call \src{compat/jax2_101/claims.py:estimate_block}: the reported point $\hat\theta$, its stored
log-likelihood, and the best point the fit evaluated (its log-likelihood, the gap, and its $\theta$ and $z$ where the
record identifies them). Two caveats travel with it. The stored value is not guaranteed to be the likelihood at
$\hat\theta$ (scipy's L-BFGS-B keeps the last trial's value after a failed line search; [T-LBFGS-STORED-VALUE]), and
$\hat\theta$ has no $z$, so a re-evaluation at $\log\hat\theta$ on box-corner human cohorts can differ by about one nat
(handoff \S6 issues 11, 14, 17c). The 1.0.1 FitResult also records \code{diagnostics["switches"]}, the class constants,
as the frozen ClaimSet does \src{calibrators/profile_jax2.py:method_diagnostics}.

\tplImpl
The port is generated, not edited: \code{tools/port_jax2_101.py} reads the frozen files from the snapshot
\code{tests/oracle/jax2_6884fc8/} (built from git objects at the tag) and writes \code{compat/jax2_101/}, keeping
\code{fp_jax}, \code{observation_model} and \code{estimator} whole, and the adapter's functions the fit needs, with its
guard deletions and sibling-import replacements counted and asserted \src{compat/jax2_101/VERBATIM.toml}. The 1.0.0
port is never touched. Both ports are handles (\code{JaxPort}); every function that reads a port takes it, and the
leverage check runs on the fit's own port: its MLE-stage refit passes the port's own \code{run_fit} keywords (1.0.1 adds
\code{inner_polish}, \code{bracket_bisect} and \code{extent_z}) plus \code{stages=("mle",)} \src{leverage.py:refit_mle}.
The six $\kappa$-fixed names 1.0.1 ships are ported but not registered: a fit asking for one is refused with its reason
\src{registry.py:UNREGISTERED}.

The check also records $g_0$, the warm start's gain at $\hat\theta$ on the full data \src{leverage.py:g0}. No label reads
it (Naveh, 2026-10-20, after the pilot below).

\tplLimits
As in Section~\ref{sec:pj}, plus the limits jax2 lists for 1.0.1 (handoff \S6; numbers copied with their lines into
\code{specs/003b-jax2-101-port/review/evidence/jax2_101_limits.json}): human claims depend on the CPU model (issue 1);
\code{profile_jax2_h}'s grid was 1.2--13.6 nats off at the truth (14.9 against a second reference) on the 3 of the 56
registered sharp hazard cohorts where it was measured, all on library v5 data, none on v6; 69 more v6 cohorts are
inferred, and the two of them measured were 0.11--0.60 nats off ([T-H-HAZARD-SHARP], issue 8; every \code{_h} 1.0.1 fit
prints this as a note, without saying whether the cohort is in the stratum); interval coverage inside \code{DT_RULE}'s
scope is expected but untested (issue 6); stale reported points (issue 11); \code{profile_jax2} is misspecified on hazard and window data (issue 12);
\code{best_point.theta} is empty where the gap is at most 0.01 nats (issue 14); and \code{DT_RULE}'s own limits: the
grid error at the profile edges on sharp cohorts of about $10^4$ is not measured, the scope's edge cases are untested,
and nothing is claimed for $\mathrm{CV}_c\ge0.05$ (issue 17). The public fit still refuses windows, entry ages and
per-row $\delta$ (spec 006).

\tplWhere
\code{compat/jax2_101/} (generated); \code{calibrators/profile_jax2.py} (the handles, \code{fit_result},
\code{method_diagnostics}, \code{version_warnings}, \code{notes_for}); \code{leverage.py}; \code{registry.py}
(several versions per name).

\tplEvidence
All side by side on this Mac (MC06268844), in \code{tests/acceptance/results/}:
\begin{itemize}[nosep]
  \item \code{census_003b.json}: the port's grid equals jax2's census bit for bit on 2,798 of 2,798 library-v6 rows;
    290 J and 52 H rows are in \code{DT_RULE}'s scope.
  \item \code{identity_off_003b.json}: with 1.0.0's ten keywords, the 1.0.1 port's result, claims and self-consistency
    block equal the 1.0.0 port's on 11 of 11 of jax2's identity lines.
  \item \code{parity_003b.json}: at its class constants the port equals the frozen 1.0.1 code on 7 of 7 pairs: 5 at the
    \code{run_fit} level and as the benchmark's ClaimSet (\code{estimate} and \code{switches} included), 2 at the MLE
    stage; 3 inside the scope, one of them a full fit; the mechanism pairs (inner polish, bracket bisection, the hazard
    path, the basin re-fit) ran on the library-v5 arrays where jax2's records showed them firing, and each fired in the
    port's run. The full in-scope fit (\code{censor-k0.01-p05-c1-s1}, $n=10^4$, no combination claims) took 19,654
    CPU-seconds, against 1,850 for a full fit outside the scope at the same $n$ (\code{hazard-k0.01-p04-c1-s1}).
  \item \code{bytes_003b.json}: value and gradient bytes equal where an age-0 death's first-interval mass is
    $2.4\times10^{-21}$ and $1.8\times10^{-35}$. v0.3.3's $S(0)-S(\delta)$ scores the $-700$ floor at both; the two
    $\theta$ were chosen by a pre-registered scan for exactly that, and across the scan v0.3.3 scored the floor on only
    34 of 127 and 56 of 164 such $\theta$, and round-off noise ($-29$ to $-37$ nats) on the rest, where the true term is
    $-47$ or lower (\code{tests/acceptance/panel_003b.toml} [parity.bytes]).
  \item \code{v100_unchanged_003b.json}: \code{profile_jax2@1.0.0} on ITP with an added age-0 death reproduces
    1.0.0's stored $\hat\theta$, log-likelihood and leverage block.
  \item \code{lev_measure_003b.json} (the pilot), \code{itp_003b.json}, \code{ac11_003b.json}: the leverage check on 1.0.1
    flags the planted age-0 death on ITP (R 38.2, 5.42 nats; its refit equals 1.0.1's fit without the death); against
    the ceilings pre-registered in \code{tests/acceptance/panel_003b.toml}, it flags 0 of 12 held-out cohorts, 0 of 47
    cohorts of 150 and 0 of 18 generated twins, and finds a planted age-0 death in 5, 6 and 6 of 6 cohorts at $n$ = 100,
    300 and 1,000 and a planted late death in 5, 6 and 4 of 6. These cohorts sit where removal is negligible
    ($\rho_\beta/\rho_\eta<0.1$ on 56 of the 59 library cohorts, 0.055--0.060 at every planted point): there
    $\bar r/\eta$ is not determined and its ratio slides freely, so the nulls test the 1.92-nat gain alone, and 9 of the
    15 late detections moved only $\bar r/\eta$ (\code{specs/003b-jax2-101-port/review/evidence/ac11_reading.json}).
    Cohorts where removal matters, fits inside \code{DT_RULE}'s scope, and \code{profile_jax2_h}'s check, are not
    tested.
\end{itemize}

\tplIssues
\begin{itemize}[nosep]
  \item $g_0$ was 0 on all 8 pilot fits: the estimator's L-BFGS-B stops at its first line search at $\hat\theta$ (8 of
    8, \code{review/evidence/pilot_lnsrch.json}); on the 7 human hazard cohorts $\hat\theta$ is not stationary
    (gradient norms 10--642), so $g_0$ cannot fail there and a \code{no_flag} means the check could not see
    ([T-LEV-G0-LNSRCH]).
  \item Deferred to a later spec (Naveh, 2026-10-21): an $\bar r/\eta$ floor in the drop's comparison, floored deaths
    ranked first, a $g_0$ that can fail, a removal-dominated AC panel, the \code{_h} note naming the fit's own SR CV, and
    a printed \code{DT_RULE} line.
  \item The fixes jax2 defers to its next version: the stored value at the returned point, recording the best point and
    $\hat\theta$'s $z$ always ([T-LBFGS-STORED-VALUE]); a hazard-aware grid for \code{_h} ([T-H-HAZARD-SHARP]).
  \item The default switch, the v6 scoring and the $\kappa$-fixed names are later specs (003g and after).
\end{itemize}

\tplRefs
\code{specs/003b-jax2-101-port/} (spec, plan, research R1--R16); jax2's handoff (\S\S1--7) and defect register REV 6
(\code{defects/DEFECT_REGISTER.txt}, sha256 \code{20ba3bc9...}); Section~\ref{sec:pj}; Section~\ref{sec:lev}.

% Section for specs/003d-cohort-generator (T049, research R14). Paths are relative to src/sraging/ unless they start
% with tests/ or tools/.
\section{The cohort generator (\texttt{sraging.generator})}\label{sec:gen}

\subsection{Overview}\label{sec:gen-overview}

\code{sraging.generator} draws a synthetic cohort from the SR model as the lab benchmark's library specifies one. It
replaces SRK's \code{benchmark.cohorts.materialize} (SRK \code{cb191f5}), whose backward-Euler grid biased the
coefficient of variation (CV) of sharp cohorts. The version is \code{"sraging-gen-1"} (git tag \code{sraging-gen-1})
\src{generator.py:GENERATOR_VERSION}. For a fixed (spec, seed, version) the output bytes never change; any change to
them is a new version. The call \src{generator.py:CohortSpec, materialize} is
\begin{verbatim}
  materialize(CohortSpec(name=, rho_eta=, rho_beta=, rho_eps=, kappa=,
                         n=, L_phys=, seed=, delta=, censor_frac=0,
                         external_hazard=None, het_param=None, het_var=0,
                         time_range=None))
\end{verbatim}
 The fields are the benchmark's own and are keyword-only. Anything else
is refused by name: the benchmark's organisational fields, a positional call, $\delta=$ \code{None}, heterogeneity in
anything but $\xc$, an off-grid window, and out-of-range values \src{generator.py:_validate}.
\code{materialize} runs, in order:
\begin{enumerate}[nosep]
  \item the model's absorbed-mass curve (one, or $K$ node curves under heterogeneity; Sections~\ref{sec:gen-curve}
    and \ref{sec:gen-het});
  \item the uniforms (Section~\ref{sec:gen-draws});
  \item the death times by inverse CDF, then the external hazard as a competing risk;
  \item the administrative censoring time $C$ and the recording on the $\delta$ grid (Section~\ref{sec:gen-obs});
  \item \code{time_range}, then the grid error and the caveats (Section~\ref{sec:gen-err}).
\end{enumerate}
It returns \code{death_times}, \code{events} (1 death, 0 censored), \code{censor_time},
\code{realised_censored_fraction}, \code{grid_error}, \code{caveats} and \code{provenance}
\src{generator.py:GeneratedCohort}.

\subsection{The curve}\label{sec:gen-curve}

Time is $\tau=t/L_{\rm phys}$, and $(\rho_\eta,\rho_\beta,\rho_\eps,\kappa)$ are the dimensionless groups of
Eq.~\eqref{eq:pv-groups} at $\xc=1$: the benchmark's $\tau$-mean $=1$ convention, SRK's. The curve is the spec 002
engine's uniform value march (Section~\ref{sec:fp-mesh}), numerics \code{"v1"}: \code{nstart} backward-Euler steps,
then BDF2. Its coefficients are $(\rho_\eta,\rho_\eps,\rho_\beta\,y_f/(y_f+\kappa))$ at the faces $y_f=(j+\frac12)/N$.
With \code{carry_F} it returns both $S(\tau_n)$ and the absorbed mass $F(\tau_n)$
\src{generator.py:_march, _solve_curve}. The march stops at the first $\tau_n>0.3$ with $S<G_{\rm tol}$. If none
exists before $\tau_{\max}$, the horizon does not absorb, and the call is refused rather than truncating the draws.
The CDF that is sampled is $F$, piecewise linear between the grid nodes. $\max_n|S+F-1|$ is recorded with every
cohort.

\paragraph{The grid.} One grid per version: $N=6000$, $\Delta\tau=2\times10^{-4}$, $\tau_{\max}=30$,
$G_{\rm tol}=10^{-12}$, \code{nstart} $=20$, and $K=161$ heterogeneity nodes \src{generator.py:GRID}. It was fixed
by a refinement sweep over every library configuration, before any acceptance run
\src{tools/acceptance_003d.py:sweep}. A configuration is what fixes a cohort's law
\src{generator.py:configuration_key}:
\begin{equation}
  (\rho_\eta,\rho_\beta,\rho_\eps,\kappa,\ \delta/L_{\rm phys},\ \text{censor\_frac},\ h L_{\rm phys},\
  \sigma_{\rm het}).
\end{equation}
The reference grid doubles $N$ and $K$ and halves $\Delta\tau$ \src{generator.py:reference_grid}. Each
configuration's binned law (Section~\ref{sec:gen-err}) is compared with the reference's, under the generator's own
$C$:
\begin{equation}
  \lambda=2\,n\,\mathrm{KL}_{\ge5}\!\left(P_{\rm gen}\,\Vert\,P_{\rm ref}\right),\qquad
  \Delta\ln\mathrm{CV}=\ln\frac{\mathrm{CV}(P_{\rm gen})}{\mathrm{CV}(P_{\rm ref})}.
\end{equation}
Here $\mathrm{KL}_{\ge5}$ merges cells from the left until $n\,P_{\rm ref}\ge5$, the G-test's binning
\src{generator.py:kl_binned}. The bars (spec 003d SC-001) are:
\begin{itemize}[nosep]
  \item homogeneous: $\lambda\le0.05$ and $\tfrac43|\Delta\ln\mathrm{CV}|\le\tfrac12\sqrt{0.1/(2n)}$, at the
    configuration's largest library $n$;
  \item heterogeneous (Naveh, 2026-10-14): $\lambda\le0.2$ at $\min(n,20{,}000)$, and the benchmark's own rule
    $|\Delta\ln\mathrm{CV}|\le\sqrt{0.1/(2n)}$ at the real $n$.
\end{itemize}
The sweep's rungs and every configuration's numbers are in \code{tests/acceptance/results/sweep_003d.json}.

\subsection{The draws}\label{sec:gen-draws}

Each cohort has one PCG64 generator, \code{numpy.random.default_rng(seed)}, used only through \code{.random(n)}
\src{generator.py:_uniforms}. It draws the death uniforms first, then the hazard uniforms, then the heterogeneity
uniforms, each only when its feature is on; the order is recorded in \code{provenance}. The death uniforms are SRK's
first draw. So a homogeneous cohort without a hazard is an individual-level twin of its SRK cohort in the SR death
time.

\paragraph{Inverse CDF.} For uniform $u$ and an individual's CDF $F_i$ (piecewise linear on the grid), the death
time is $\tau=\tau_{k-1}+\frac{u-F_i(\tau_{k-1})}{F_i(\tau_k)-F_i(\tau_{k-1})}(\tau_k-\tau_{k-1})$. Here $k$ is the
smallest node with $F_i(\tau_k)\ge u$, found by bisection \src{generator.py:_inverse}. A $u$ above the CDF's last
value is placed at the last node and counted (\code{n_beyond_cdf}). The death age is $t=\tau L_{\rm phys}$.

\paragraph{External hazard.} A competing exponential at rate $h=e^{-\text{external\_hazard}}$ per age unit:
$t\leftarrow\min\bigl(t,\,-\log(1-u_h)/h\bigr)$ \src{generator.py:_hazard_times}.

\subsection{Heterogeneity in $\xc$}\label{sec:gen-het}

$\xc\to\xc s$ with $s=1+\sigma z$, where $\sigma$ is \code{het_var} and $z$ is a standard normal truncated to
$s>0$. In dimensionless terms a node at $s$ is $(\rho_\eta/s,\ \rho_\beta/s,\ \rho_\eps/s^2,\ \kappa/s)$
\src{generator.py:_node_point}. An individual's $z$ is continuous, from its own uniform:
$z=\Phi^{-1}\bigl(a+u_s(1-a)\bigr)$ with $a=\Phi(-1/\sigma)$ \src{generator.py:_het_z}.

There are $K$ node curves, uniform in $z$ on
$[z_a,z_b]=[\max(\Phi^{-1}(10^{-9}),\,-(1-s_{\min})/\sigma),\ \Phi^{-1}(1-10^{-9})]$ with $s_{\min}=0.05$
\src{generator.py:_het_nodes}. The individual's CDF is the linear interpolation in $z$ between the two bracketing
node CDFs, $F_i=(1-w)F_j+wF_{j+1}$; outside $[z_a,z_b]$ it is the end node's, counted as \code{n_outside_het_nodes}
\src{generator.py:_bracket}. The cohort's law is then exactly the hat-weighted mixture
$\sum_jW_jF_j$, with
\begin{equation}
  W_j=\frac{1}{1-\Phi(-1/\sigma)}\int\mathrm{hat}_j(z)\,\phi(z)\,\dd z,
\end{equation}
and the truncated mass beyond either end assigned to the end node. The convergence check and the G-test compute
this mixture. SRK used 40 equal-weight midpoint nodes, whose CV is 0.7--1.9\% lower (the hand-off message,
\code{docs/planning/handoff_2026-10-14_generator_to_tester.md}).

\subsection{Censoring, recording and the window}\label{sec:gen-obs}

The observed survival is $S_{\rm obs}(t)=\bigl(1-\sum_jW_jF_j(t/L_{\rm phys})\bigr)\,e^{-ht}$
\src{generator.py:_S_obs}.

\paragraph{Administrative censoring (R7).} $t^*$ solves $S_{\rm obs}(t^*)=$ \code{censor_frac} by bisection. $C$
is whichever of the grid points $\lfloor t^*/\delta\rfloor\delta$ and $(\lfloor t^*/\delta\rfloor+1)\delta$ (excluding
0) has $S_{\rm obs}$ closer to \code{censor_frac}; a tie goes to the earlier one
\src{generator.py:_censor_time}. The call is refused if $S_{\rm obs}(C)<10^{-12}$. \code{censor_frac} is thus a
target for the model, not for the sample: on a sharp cohort with a coarse $\delta$, no grid point may come near it,
and the cohort then says so (Section~\ref{sec:gen-err}). Nearest in probability replaced nearest in time on
2026-10-14 (Naveh), when the latter missed by up to 0.33.

\paragraph{Recording.} With censoring, $t\ge C$ gives $(C,0)$, and otherwise
$(\min(\lfloor t/\delta\rfloor,\,C/\delta-1)\,\delta,\ 1)$. Without censoring every row is
$(\lfloor t/\delta\rfloor\delta,\ 1)$ \src{generator.py:_record}. There is no clamp: a death in the first bin is
recorded at 0, and no death is recorded at $C$. Rows are sorted by (age, event).

\paragraph{The window.} \code{time_range} $=(lo,hi)$ must lie on the $\delta$ grid, or it is refused. Ages below
$lo$ are dropped; ages above $hi$ become $(hi,0)$; a death recorded at $hi$ stays a death. A window that keeps no
one is refused \src{generator.py:_apply_time_range}. This is the benchmark's \code{apply_window}, closed at both
ends.

\subsection{The binned law, the grid error and the caveats}\label{sec:gen-err}

The exact law of the recorded outcome, before the window, is \src{generator.py:binned_distribution}:
\begin{equation}
  P_k=S_{\rm obs}(k\delta)-S_{\rm obs}((k+1)\delta)\ \ (k<C/\delta),\qquad P_C=S_{\rm obs}(C).
\end{equation}
Without censoring, the tail beyond the march's end joins the last cell. Every cohort carries its configuration's
$\lambda$ (at its own $n$) and $\Delta\ln\mathrm{CV}$ in \code{grid_error}. They come from the packaged point table
when the configuration is in it, and otherwise are computed against the reference grid at the call
\src{generator.py:_grid_error_for}. \code{caveats} lists, with numbers:
\begin{itemize}[nosep]
  \item $\lambda$ or $|\Delta\ln\mathrm{CV}|$ above its bar at the cohort's own $n$;
  \item a heterogeneous cohort above $n=20{,}000$;
  \item a realised censored fraction more than 0.1 from \code{censor_frac}.
\end{itemize}
An empty tuple means none. Each configuration's converged and binned CV, $C$ and censored mass are in
\code{tests/acceptance/results/point_table_full_003d.json}.

\subsection{The backward-Euler control}\label{sec:gen-be}

SRK \code{cb191f5}'s \code{materialize} is kept as an acceptance tool, never as a generator version
\src{tools/v5_reproduce_003d.py:curve, materialize_srk}. It is defined by:
\begin{itemize}[nosep]
  \item a backward-Euler march on $N=1500$, $\Delta\tau=5\times10^{-4}$;
  \item inverse CDF from $S$ interpolated linearly;
  \item censoring at the sample's $(1-f)$ quantile;
  \item flooring, then a clamp to $\delta/2$, which the benchmark's \code{correct_first_bin} undoes.
\end{itemize}
It rebuilds the benchmark's library v5 homogeneous cohorts from sraging alone
(\code{tests/acceptance/results/v5_reproduce_003d.json}), and v2's 55 censored coverage cohorts byte for byte.
It is the positive control of the CV-fidelity check and the G-test, both of which it must fail. Its $n\to\infty$
law gives the fit-free prediction \src{tools/acceptance_003d.py:be_law}. Under that law, the deaths in
$[\lfloor C\rfloor_\delta, C]$ share the censored records' age.

\subsection{Validation evidence}\label{sec:gen-evidence}

Each check, run on a clean committed tree, and where its result is recorded:
\begin{center}\small
\begin{tabular}{@{}>{\raggedright\arraybackslash}p{0.5\linewidth}>{\raggedright\arraybackslash}p{0.46\linewidth}@{}}
\toprule
check & result file, in \code{tests/acceptance/results/} \\
\midrule
grid refinement, every library configuration (SC-001) & \code{sweep_003d.json}, \code{het_nodes_all_003d.json} \\
CV of recorded ages against the binned reference (SC-002) & \code{fidelity_003d.json} \\
G-test of every library cohort, with a null control (SC-003) & \code{gtest_003d.json} \\
recording and \code{time_range} against \code{apply_window} (SC-004, SC-005) & \code{gtest_003d.json}, \code{window_003d.json} \\
determinism on the Mac and on WEXAC AMD EPYC (SC-006) & \code{determinism_003d.json}, \code{determinism_epyc_003d.json} \\
cost (SC-008) and the per-configuration table (FR-015) & \code{cost_003d.json}, \code{point_table_full_003d.json} \\
an independent FP code at the truth & \code{calibrator_dt_003d.json} \\
\bottomrule
\end{tabular}
\end{center}

\paragraph{An independent code.} \code{tools/calibrator_dt_003d.py} evaluates \code{profile_jax2}'s own objective
(jax2's \code{fp_jax}, independent of the spec 002 engine) at the truth. It does this on the expected-count cohorts
of 40 configurations, and compares with the generator's binned law. On a grid refined to $(2N,\Delta t/16)$ the two
agree to $1.15\times10^{-6}$ nats per individual or better. On \code{profile_jax2} 1.0.0's production grid
($\Delta t=0.002$ at $\delta=1$, $L\approx100$) its likelihood at the truth is up to 0.214 nats per individual too
low at sharp, small-$\kappa$ points. That is a defect of the calibrator's time step, not of the generator: its own
\code{DT_RULE} grid ($\Delta t\le\mathrm{CV}/80$, built and off) brings it to $1.8\times10^{-5}$ or better
(\code{calibrator_dt_003d.json}).

\subsection{Assumptions and validity limits}\label{sec:gen-limits}

\begin{itemize}[nosep]
  \item \textbf{One model.} The generator and the calibrators solve the same Fokker--Planck model, so an error in its
    equations would be common to both. Only the discretisations are checked independently.
  \item \textbf{Heterogeneity.} Only $\xc$ at $\sigma=0.2$ is verified. Heterogeneous cohorts are accepted at
    $n\le20{,}000$; above that they say so.
  \item \textbf{Validated sizes.} The grid is validated at each configuration's library $n$. A configuration
    outside the packaged table is checked at the call (a heterogeneous one costs about 1.3 CPU-h, inferred) and not
    at the version level.
  \item \textbf{Process state.} Importing the module sets JAX to float64 and \code{PJRT_NPROC=1} for the whole
    process \src{generator.py:_runtime.configure}.
  \item \textbf{Not provided.} SRK's planted defects, $\delta=$ \code{None}, and heterogeneity in $\eta$, $\beta$,
    $\eps$ or $\kappa$.
\end{itemize}

% Section for specs/003g-default-switch (T054; research R15; constitution V, XI, XII). Paths are relative to src/sraging/
% unless they start with tests/, tools/ or specs/. Numbers come from 003g's result files, named where they are used.
\section{The default's gate (\texttt{sraging.default\_gate}; spec 003g)}\label{sec:gate}

\tplPurpose
A calibrator becomes sraging's default only after it passes where a registered defect of it, or of the calibrator it
replaces, fires (constitution V), against rules fixed before the records it scores were read, every later change
recorded as a dated deviation (XI). Spec 003g applies this
to the switch from \code{profile_v033} (v0.3.3) to \code{profile_jax2} (jax2 1.0.1). The evidence is jax2's library-v6
records, which it delivers in three windows (W1: the reference set; W2: \code{_h}; W3: the sharp and nladder groups);
sraging fits two cohorts of its own (Section~\ref{sec:gate-cross}).

\tplMath
\textbf{The panel.} A frozen list (\code{tests/acceptance/panel_003g_ids.tsv}), built read-only from six sources each
pinned by sha256 (\code{tools/panel_003g.py}): the defect register's v6 lists, the v6 manifest (each cohort's array
checked against its content digest), the reference set, the validity audit, the basin labels and jax2's grid census.
Its strata: the 53 \code{[T-JAX-BASIN]} cohorts mapped onto v6 (E.\,coli has no v6 cohort: a recorded deviation; the 13
hazard ones are \code{_h}'s, reported); the register's lists in J's domain (no window, no hazard arm), 507 cohorts, 269
inside the reference set and 238 outside; and three strata of the predecessor's other defects, by a data rule: a death
in the first recording bin (324 cohorts), the censor-tie rule of jax2's own detector (0 on v6: an empty stratum, never
a pass), and $\delta/L<10^{-3}$ (16).

\textbf{The witness.} For a cohort with truth $\theta^\ast$ and fitted $\hat\theta$,
\begin{equation}
  W=\ell(\theta^\ast)-\ell(\hat\theta),
\end{equation}
with $\ell$ the candidate's own log-likelihood on its own grid for that cohort, by the value-and-gradient program, as
jax2's G2 witness takes it (\code{GATEV6/PREREG_V6.txt}:517); on a human cohort jax2's Richardson-extrapolated
$W_{\mathrm{REF}}$. $W$ is defined only on informative cohorts (no heterogeneity): the candidate's likelihood cannot
score a heterogeneous truth. The candidate's $W$ is read from jax2's witness files, each pinned by sha256 and checked
against the record it witnesses (\code{tools/witness_003g.py}).

\textbf{The rules} (pre-registered; Naveh's OK 2026-10-22, with two corrections made before any v6 record was read;
PR-4's followed a dry run on library-v5 records)
\src{default_gate.py}:
\begin{enumerate}[nosep]
  \item[PR-1] a stratum fails if any valid informative cohort has $W>1.92$; it is incomplete if any lacks a fresh
    witness; the exact one-sided 95\,\% bound on the failure rate is reported;
  \item[PR-2] inside the reference set, per arm and quantity, the tester's own paired comparison of the candidate
    against v0.3.3 (\code{build_comparison}): a stratum fails if any calibration outcome (two-sided or one-sided)
    reads ``loses'' ($z<-2$);
  \item[PR-3] outside it (no v0.3.3 record), the candidate's coverage hits against its own on the same arm and
    quantity inside: the one-sided 95\,\% Newcombe hybrid-score lower bound of $p_{\mathrm{out}}-p_{\mathrm{in}}$ must
    exceed $-0.10$. (The first approved margin, 0.05, was corrected after its power was computed. A power table built
    from W1's closed counts and the panel's $N$ alone, committed before W3 (\code{pr3_power_003g.json}), shows that
    $N\le3$ cannot pass, and that equal coverage passes at the counts the inside closed share predicts: on SRK-1,
    154 of 161 against 617 of 644 gives $-0.038$ and 45 of 46 against 194 of 199 gives $-0.068$. Beside every PR-3 FAIL
    the gate prints that row and, after W3, the probability of a count as low under exactly equal coverage,
    $P(K\le k_{\mathrm{out}}\mid n_{\mathrm{out}},p_{\mathrm{in}})$.)
  \item[PR-4] the candidate's median closed-interval width (dex) on $\eta$ and $\eps$ below the data-ignoring null's by
    more than the tester's 0.05\,dex tolerance; its closed share is reported. (The tester's information axis also
    fails a method whose closed share is below the reference's, and the null is closed on every cohort.)
  \item[PR-5] v0.3.3 is shown failing on each of its defects' strata when at least one cohort shows the defect:
    $W_{33}>1.92$ (the basin and round-off strata), its stored CPU time above the candidate's (the $\delta/L$ ceiling);
  \item[PR-6] the switch is blocked by any failed or incomplete PR-1 to PR-4 verdict, any PR-5 ``not shown'', or a
    required verdict missing; each blocked stratum goes to Naveh, who decides it (fix first, or a stated limit).
\end{enumerate}
\textbf{Deviations after W1} (the sr-naveh interim review, Naveh 2026-10-28; the panel's \code{[deviations]}, the
registered \code{[rules]} unchanged): DV-1, the $\delta/L$ ceiling is a library gap (v0.3.3's cap $\delta/(3L)$ costs
more than 1.0.1's floor only below $\delta/L\approx4.4\times10^{-4}$; v6's 16 cohorts sit at $\approx6\times10^{-4}$), so
its PR-5 reads ``reported''; DV-2, a PR-2 verdict whose every axis is ``not measured'' reads ``reported'', not PASS (7
strata); DV-3, \code{not_assessed} per individual (Section~\ref{sec:lev-rule}). At the close (the sr-naveh review,
Naveh 2026-10-29): DV-4, P-ROUNDOFF's PR-5 reads ``reported'', because every cohort shown there is human with the basin's
signature and none of the 112 non-human ones is shown, so the round-off defect itself is not shown on v6; DV-5, a PR-2
PASS whose measured axes all have fewer than 3 paired cells reads ``reported'' (below 3, srbench's cluster-robust $z$
cannot reach $-2$, enumerated over every cell difference and cluster partition). Naveh's decision on each blocking
entry is the panel's \code{[[decisions]]}, copied onto the gate's entries. The rule code is frozen at W1's pin
except for named changes, and a result is stale when a commit since it ran is not a reviewed one
(\code{tests/acceptance/test_results_003g.py}).

\textbf{v0.3.3's point} (no $\theta$ is stored; \code{tools/v033_points.py}): the four claim points; else $\eta$,
$\beta$, $\eps$ from the claims and $\kappa$ by a 1-D maximisation of v0.3.3's own likelihood over $\log\kappa$ on the
record's grid, kept iff it reproduces the stored log-likelihood within 0.5\,nats (the tester's basin-scope recipe).
$W_{33}$ is then taken in the candidate's likelihood, never in v0.3.3's: at its own wrong-basin points v0.3.3's value
is host-dependent, by more than 1\,nat on 5 of the 32 human basin cohorts and up to 3,346\,nats
(\code{tests/acceptance/results/v033_points_verify_003g.json}). Where no v0.3.3 record exists (the three p15 cohorts), its
MLE stage is replayed from outside the frozen port, call for call (\code{tools/v033_mle.py}; equal to a real fit bit
for bit, \code{v033_mle_verify_003g.json}).

\tplImpl
\code{tools/acceptance_003g.py}: one phase per step, each on a clean committed tree; each window's runs, records and
witness files are pinned (\code{tests/acceptance/pins_003g.json}) and committed before anything is scored; v0.3.3's v6
records are sealed until jax2's gate run, and only the predecessor phase reads them.

\subsection{sraging's own fits}\label{sec:gate-cross}
Two fits tie jax2's records to the port a user runs. Fit 1 is the default settings on an in-scope sharp cohort
(\code{censor-k0.01-p05-c1-s1}); its claim shapes and coverage hits must agree with its W3 record, and the interval
ends are reported against the 1.0.0 host spread (spec 003e: 0 shape flips on 140 cohorts, 6 of 745 ends beyond
0.03\,dex). Fit 2 was registered as a human basin cohort whose basin re-fit fired; in W1 the re-fit fired on no basin
cohort (only on 22 heterogeneous and 13 realfull records), so it is a valid human realfull cohort
(\code{realfull-real-Sweden_F_1910_homo-c1-s1}) whose re-fit fired on the cluster. On this Mac it did not fire: the
MLE stage went straight to $\beta$ at the box. Fit 2 is therefore a host-agreement check, and the port's re-fit path
has not run on this Mac. Its shapes and hits agree with the record; 2 of its 6 interval ends differ by more than
0.03\,dex (0.04 on $\beta$'s lower end, 0.11 on $\kappa$'s upper end), and the log-likelihoods by $-0.41$\,nats, which is
$1.7\times10^{-5}$ per individual at the $n=23{,}448$ fitted; that the two hosts sit at different points of the ridge is inferred,
not measured. Fit 2's stored leverage reason predates DV-3's wording; its label is current.

\tplEvidence
From \code{tests/acceptance/results/}: \code{gate_003g.json} (``blocked; decided (11 of 11)'', all 11 entries
outside the reference set),
\code{score_w1_003g.json}, \code{predecessor_003g.json}, \code{score_w3_003g.json}, \code{pr3_power_003g.json}.
\begin{itemize}[nosep]
  \item Inside: PR-1 passes on every scored stratum and arm (largest $W=0.287$); PR-2 loses nowhere, and the strata
    that measured nothing (DV-2) or could not show a loss (DV-5) read ``reported''; PR-4 passes everywhere. 1.0.1 often declines to cap $\eta$ or $\eps$ from above (on the
    censored LB-2 cohorts $\eta$ is a lower bound only on 30 of 66), which PR-4's widths do not see.
  \item PR-5: v0.3.3 is shown failing on 29 of 32 basin cohorts and on 3 of 3 p15 cohorts. On the round-off stratum
    the 19 of 178 shown are all human with the basin's signature (15 of them basin cohorts), and none of the 112
    non-human ones is shown (largest $W_{33}=-0.034$): DV-4. Of the 3 basin cohorts not shown, one was recovered with
    $W_{33}=-0.06$; the other two were not recovered, one a real failure (recon $-27$), one a host call (recon $-0.545$
    against the $-0.5$ cut). The human $W_{33}$ is $W_4$, not $W_{\mathrm{REF}}$ (a deviation from research R6: all 29
    are above 10\,nats, and $W_4-W_{\mathrm{REF}}$ is about 1\,nat, so nothing flips). The $\delta/L$ ceiling is DV-1;
    the censor tie is empty.
  \item Outside: PR-1 passes (largest $W=-0.01$ on 56 censored and 171 witnessed pure sharp cohorts), except that one
    valid cohort, \code{pure-k0.0316-p15-c0-s1}, has no record (incomplete). PR-4 passes. PR-3 passes on $\eta$ and $\eps$ for
    the pure sharp cohorts and on $\eps$ for the censored ones; it fails where the power table says it cannot pass
    (the p15, LB-1x and SRK-2 strata, $N\le3$, every claim covering; $\beta$ on 2 and 4 closed claims), and on
    censored $\eta$, 43 of 45 against 194 of 199, one claim short; a count that low has probability 0.31 under
    exactly equal coverage.
  \item What PR-3 and PR-4 do not see, reported in \code{score_w3_003g.json} (\code{limits}) and printed on every
    default fit in DT\_RULE's regime (\code{warnings.SHARP_COHORT_LIMIT}): $\eta$ has no upper end on 49 of the 227
    witnessed sharp cohorts, every one of those lower bounds covering, and on each of them $\beta$ has no claim or only
    the $\kappa$ box's end; median closed widths are 1.50--1.67\,dex ($\eta$) and 2.42--2.47\,dex ($\eps$); raw $\beta$
    is closed on 6 of 227, and its upper end is always the $\kappa$ box's.
  \item Naveh (2026-10-28) made every blocked entry a stated limit and switched the default
    (\code{registry.py:DEFAULT}). Fit 1 agrees with its record (shapes and hits equal, ends within $10^{-4}$\,dex); it
    took 59,867 CPU-s, over its price of 8.6--15\,h.
  \item On realfull cohorts $W$ cannot fail: 1.0.1 models neither their fit window nor their external hazard. Of the 66
    valid ones inside, 39 have $W<-10$ (from $-0.09$ to $-14{,}259$; per individual of the library cohort, which can exceed the $n$ fitted by about 6\,\%, a median of $-0.010$, from
    $-1.3\times10^{-5}$ to $-1.58$), and 1.0.1's closed intervals covered 26 of 61 ($\eta$) and 20 of 57 ($\eps$)
    (\code{score_w1_003g.json}, \code{realfull_all}).
  \item Human records: \code{BOX_EDGE} fires on 85 of 396 ($\beta$ 45, $\kappa$ 40); closed $\eta$ widths are
    bimodal (quartiles 0.11, 0.62, 2.28\,dex) (\code{human_box_edge}).
\end{itemize}

\tplRefs
\code{specs/003g-default-switch/} (spec with its Clarifications, research R1--R15); jax2's \code{GATEV6/PREREG_V6.txt};
Section~\ref{sec:j101}; Section~\ref{sec:lev}.

% Section for specs/004-samplers (T132; FR-020; research R1-R8, R26, R28). Paths are relative to src/sraging/ unless
% they start with tests/ or tools/. Every measured number carries a \verify marker against
% tests/acceptance/results/*_004.json (macros.tex).
\section{Heterogeneity: the truncated normal on a checked Gauss rule}\label{sec:het}

\subsection{Overview}\label{sec:het-overview}

Spec 004 gives \code{sraging} one heterogeneity definition at a known cost. A heterogeneous parameter $P$ (one of
$\eta,\beta,\eps,\kappa,\xc$) is $P\,s$ for each individual, with the scale $s$ drawn from a truncated normal: gen-1's
law, so the FP engine, the cohort generator (Section~\ref{sec:gen}) and the Langevin engine
(Section~\ref{sec:lgv}) share one model \src{engine/inputs.py:Heterogeneity}. The FP likelihood of a cohort is the
mixture over $s$, evaluated on a Gauss rule for the truncated measure whose node count $M$ is chosen per cohort by a
checked ladder. This replaces spec 002's \code{normal_frac_sd}: $P(1+vZ)$ on $M$ probabilists' Gauss-Hermite nodes
with the non-positive nodes dropped and $M$ fixed by the caller and never checked (registers D-116, D-117). That
definition stays selectable by name, byte for byte; the change of default is register D-401.

\subsection{The definition}\label{sec:het-def}

For a spread $\sigma\in(0,1]$ (\code{sd}),
\begin{equation}\label{eq:het-tn}
  s = 1+\sigma Z,\qquad Z\sim\mathcal N(0,1)\ \text{truncated to}\ s>0,\ \text{i.e.}\ Z>-1/\sigma,
\end{equation}
with density $\varphi(z)/\Phi(1/\sigma)$ on $z>-1/\sigma$. The parameter is $P\,s$; at $\sigma=0.2$ the truncated
mass $\Phi(-5)$ is $2.9\times10^{-7}$, at $\sigma=0.5$ it is $\Phi(-2)=2.3\%$. The object
\code{Heterogeneity(parameter=, sd=, M="auto", check=True)} names one parameter; a spread given as \code{v} without
naming \code{definition="normal_frac_sd"} is refused (OD-8), so \code{v} never silently changes meaning. $M$ is
\code{"auto"} (the checked ladder, Section~\ref{sec:het-ladder}) or a fixed integer in $[1,192]$, checked against the
smallest rung $\ge\lceil1.5M\rceil$ unless \code{check=False}, which is never clean
\src{engine/inputs.py:Heterogeneity}. Several parameters are
a tuple, which only the Langevin engine accepts (Section~\ref{sec:lgv-multi}).

\subsection{The Gauss rule and its derivative in $\sigma$}\label{sec:het-rule}

\code{truncated_normal_rule(sd, M)} is Gaussian quadrature for the measure~\eqref{eq:het-tn} in $z$
\src{engine/quadrature.py:truncated_normal_rule, _tn_nodes}:
\begin{enumerate}[nosep]
  \item Discretise the measure on 4,000 Gauss-Legendre points over $[\max(-1/\sigma,-8),\,8]$ (\code{BASE_POINTS},
    \code{Z_HI}); the base grid is computed once per process.
  \item Run the monic discretised Stieltjes procedure on it: the recurrence coefficients $(a_k,b_k)$ of the orthogonal
    polynomials of the discretised measure, $k<M$.
  \item Golub-Welsch: the eigenvalues of the $M\times M$ Jacobi matrix are the nodes $z_k$; the squared first
    components of its eigenvectors are the weights $w_k$, which sum to 1.
\end{enumerate}
The scales are $s_k=1+\sigma z_k>0$ for every $k$: no node is ever dropped, so nodes and weights move smoothly with
$\sigma$ (spec 002's dropped node was the jump of D-116). The rule is written in \code{jax.numpy} in float64, one
jitted program per $M$, so the $\sigma$-derivatives $\partial z_k/\partial\sigma$ and $\partial w_k/\partial\sigma$
come from \code{jax.jvp} on the same implementation as the values (\code{dz_dsd}, \code{dw_dsd}). Rules are memoised on
$(\sigma,M)$, returned read-only, and carry the sha256 of their $z\Vert w$ bytes, so an evaluation records the nodes it
used. \code{composite_rule(sd, panels, order=8)} is the cross-check: Gauss-Legendre panels of order 8 on the same
measure \src{engine/quadrature.py:composite_rule}.

\paragraph{The mixture and its gradient.} With $P_k$ the node-$k$ survival quantities at the observation ages (the
model at parameters scaled by $s_k$), the cohort's quantities are $\sum_k w_kP_k$ and the log-likelihood is the
engine's assembly on them \src{engine/mixture.py:_fold, _ll}. The gradient takes two passes
\src{engine/mixture.py:_pass2}: pass 1 computes and mixes the points; the parent runs \code{jax.vjp} of the assembly at
the mixture, giving the cotangent $c$; pass 2 runs each node's VJP with $w_kc$ and sums the parameter parts in node
order. The $\sigma$ term is
\begin{equation}\label{eq:het-dsd}
  \frac{\partial\ell}{\partial\sigma}=\sum_k\Big[c\cdot P_k\,\frac{\partial w_k}{\partial\sigma}
  +w_k\Big(c\cdot\frac{\partial P_k}{\partial s}\Big)\Big(z_k+\sigma\frac{\partial z_k}{\partial\sigma}\Big)\Big],
\end{equation}
since $\partial s_k/\partial\sigma=z_k+\sigma\,\partial z_k/\partial\sigma$. The value never depends on whether a
gradient was asked for.

\subsection{The ladder, its check, and the curve check}\label{sec:het-ladder}

\paragraph{The ladder.} The rungs are $192/1.5^j$ rounded: 5, 8, 11, 17, 25, 38, 57, 85, 128, 192
\src{engine/mixture.py:RUNGS}. No node count above 192 is ever evaluated (\code{CAP}). For \code{M="auto"} the walk
\src{engine/mixture.py:_walk} starts at a guided rung, evaluates $\ell(M)$ and $\ell$ at the next rung, and stops at
the first rung with
\begin{equation}\label{eq:het-check}
  |\ell(M)-\ell(M_{\rm next})|<0.05\ \text{nats},
\end{equation}
returning $\ell(M)$ with the verdict \code{converged}. If the top pair (128, 192) still fails, it returns $\ell(192)$
with \code{unconverged_at_cap} and the warning \code{HET_UNCONVERGED}: never clean, with no fallback. A fixed $M$ is
checked against the smallest rung $\ge\lceil1.5M\rceil$ \src{engine/mixture.py:check_rung, _fixed}.

\paragraph{The start.} The guide \src{engine/mixture.py:guide} measures how much the mixture spreads the
first-passage time beyond the homogeneous curve: from six node marches (the homogeneous one and a 5-node rule),
\begin{equation}
  r=\frac{\sqrt{\mathrm{CV}_{\rm mix}^2-\mathrm{CV}_c^2}}{\mathrm{CV}_c},\qquad
  \text{start}=\min\{\text{rung}\ \ge 10\,r\,g(n)\},\quad g(n)=\max(1,(n/10^4)^{0.3}).
\end{equation}
Where $r$ is undefined (a zero first moment, or $\mathrm{CV}_c=0$ at a degenerate point) the ladder starts at the
lowest rung and the check still chooses $M$. The start saves rungs; it never decides the value.

\paragraph{The curve check (R26).} A reference curve has no data and so no likelihood to check. Its node count is
chosen instead by the first rung $M$ with $\sup_t|S_M(t)-S_{M_{\rm next}}(t)|$ below a stated bar: $10^{-4}$ for the
exact population of \code{population(..., n=None)} \src{population.py:fp_curves}, and ${\rm crit}_{95}/10$ at the
sample's own $n$ (${\rm crit}_{95}=1.358/\sqrt n$, over the sample's raw times) for the Langevin's error statement
\src{langevin/sample.py:Reference.curve_check}. The same rule, applied per dimension, sizes the tensor reference
(Section~\ref{sec:het-tensor}) and the damage density (Section~\ref{sec:lgv-paths}).

\subsection{Adaptive $M$ in a series}\label{sec:het-adaptive}

Inside a maximisation at fixed $\sigma$ the node count can be lower while the optimiser explores
\src{engine/series.py:maximise_fixed_sd}: (A) L-BFGS-B in log parameters at $M_{\rm low}$, two rungs below the start
rung (floor 5); (B) the checked ladder at $\hat\theta_A$ gives $M_c$; (C) L-BFGS-B at $M_c$ from $\hat\theta_A$;
(D) the check at $\hat\theta$, raising $M_c$ one rung and repeating (C) while it fails. The final value is always a
checked one; \code{evaluate(params, phase="final")} runs the ladder afresh from the guide, so no value depends on the
call history. A profile point is the same maximisation with one more parameter fixed.

\subsection{The parallel mixture and its fixed-order sum}\label{sec:het-pool}

A mixture's nodes can run over $W$ worker processes \src{engine/pool.py:NodePool}. Each worker builds the evaluator
once and computes the nodes it is sent; the parent folds the returned $P_k$ in ascending node index, in float64, on the
host \src{engine/pool.py:in_order}. The sum's order never depends on which worker finished first, so the value and
gradient bytes are those of $W=1$ for every $W$ and $M$, including $M$ not divisible by $W$. Node programs are compiled
once per process and their compile CPU is recorded apart \src{engine/mixture.py:Programs}.

\subsection{The cost model}\label{sec:het-cost}

Each rung or phase states its CPU before it runs \src{engine/mixture.py:_Ledger, time_hom; engine/cost.py}:
\begin{equation}
  \text{stated CPU}=M\times t_{\rm hom}\times f_{\rm mode},
\end{equation}
with $t_{\rm hom}$ the median of three timed homogeneous values through the same per-node programs after a warm-up,
measured in the same process, and $f_{\rm value}$, $f_{\rm grad}$ constants frozen from pilot P0. The statement keeps
the start rung, the worst case up to the cap, and after the work the measured CPU (parent and workers) and wall time;
compilation is stated and measured apart. P0 measured
$f_{\rm value}=\verify{tests/acceptance/results/p0_het_cost_004.json}{f_value}{.3f}{0.870}$ and
$f_{\rm grad}=\verify{tests/acceptance/results/p0_het_cost_004.json}{f_grad}{.3f}{3.132}$ homogeneous values per
node.

\subsection{The tensor-product reference}\label{sec:het-tensor}

FP's likelihood and curves take one heterogeneous parameter. For a Langevin sample with $d=2$ or 3 spreads, the error
statement's reference is the tensor product of the per-dimension rules \src{engine/tensor_reference.py:TensorReference}:
\begin{equation}
  S(t)=\sum_{k_1,\dots,k_d}\Big(\prod_j w_{j,k_j}\Big)\,S\big(t;\theta\odot s_{(k_1,\dots,k_d)}\big),
\end{equation}
run by the same per-node program with a five-factor scale vector (one factor per parameter, 1.0 for the others: exact)
and folded lexicographically over the parameter-sorted node indices \src{engine/tensor_reference.py:fold}. The node
counts are chosen per dimension by R26's check: dimension $j$ climbs to its next rung only while
$\sup|S_{(\dots,M_j,\dots)}-S_{(\dots,\mathrm{next}(M_j),\dots)}|$ exceeds the bar, and only the failing dimensions
climb \src{engine/tensor_reference.py:TensorReference._climb}. The product is capped at 20,000 nodes
(\code{BUDGET}); above it, or at $d>3$ (\code{D_MAX}), the sample's statement reads ``no FP reference'' and is not
clean. One member at $\sigma=0$ is one node at $s=1$, $w=1$, and reduces to the single-parameter curve bit for bit.
The composite cross-check (\code{composite_rule} panels per dimension) is HR-11's (Section~\ref{sec:lgv-multi}).

\subsection{Validation evidence}\label{sec:het-evidence}

Each rule below is a pure function in \code{sampler_rules.py} judged on a frozen panel (\code{panel_004.toml}), with a
planted case beside its twin (contracts/register-and-panel.md §3).
\begin{itemize}[nosep]
  \item \textbf{HR-1, the check:}
    \verify{tests/acceptance/results/hr1_004.json}{verdicts[0].status}{}{pass}; no false pass
    (\verify{tests/acceptance/results/hr1_004.json}{verdicts[0].statistic.n_false_pass}{d}{0}) among
    \verify{tests/acceptance/results/hr1_004.json}{verdicts[0].statistic.n_checked_valid}{d}{159} checked cells,
    Clopper-Pearson upper bound
    \verify{tests/acceptance/results/hr1_004.json}{verdicts[0].statistic.cp_upper_95}{.3g}{0.0187}.
    The result's overall status is
    \verify{tests/acceptance/results/hr1_004.json}{status}{}{fail} through FR-009
    (\verify{tests/acceptance/results/hr1_004.json}{verdicts[1].status}{}{fail}): gen-1's 161-node law sits more
    than 0.1 nats from the converged mixture on 4 cohorts above $n=20{,}000$ (register D-412).
  \item \textbf{HR-2, $\partial\ell/\partial\sigma$ against central differences:}
    \verify{tests/acceptance/results/hr2_004.json}{status}{}{pass}; the largest difference
    \verify{tests/acceptance/results/hr2_004.json}{verdicts[0].statistic.max_abs_diff}{.3g}{0.101} against the
    tolerance \verify{tests/acceptance/results/hr2_004.json}{tol}{.3g}{0.601}.
  \item \textbf{HR-3, adaptive against full $M$:}
    \verify{tests/acceptance/results/hr3_004.json}{status}{}{fail}: the largest
    $|\Delta\ell|$ is \verify{tests/acceptance/results/hr3_004.json}{verdicts[0].statistic.max_abs_dll}{.3g}{7.35e+04}
    nats, and the CPU saving is not claimed (upper 95\,\% bound of the mean log ratio
    \verify{tests/acceptance/results/hr3_004.json}{verdicts[0].statistic.saving.upper_95_log_ratio}{.3g}{0.0548}).
  \item \textbf{HR-4, the parallel bytes:} \verify{tests/acceptance/results/hr4_004.json}{status}{}{fail}: every
    (W, M) value and gradient is byte-equal to $W=1$, but the remainder plant (dropping $M\bmod W$ nodes) left the bytes
    unchanged on some sized cells, which the rule reads as a failure. The wall-time speedup at
    $W=\verify{tests/acceptance/results/hr4_004.json}{verdicts[0].statistic.speedup.W}{d}{5}$ is claimed (upper bound
    of the mean ratio \verify{tests/acceptance/results/hr4_004.json}{verdicts[0].statistic.speedup.upper_95_ratio}{.3g}{0.518}).
  \item \textbf{HR-5, the cost statement:} \verify{tests/acceptance/results/hr5_004.json}{status}{}{pass}; a share
    \verify{tests/acceptance/results/hr5_004.json}{verdicts[0].statistic.share_inside}{.4f}{0.9938} of cells within
    $[0.5,2]$ of the stated CPU and a median ratio
    \verify{tests/acceptance/results/hr5_004.json}{verdicts[0].statistic.median_ratio}{.3f}{1.097}.
  \item \textbf{D-401, the change of default:} the largest per-cohort $|\Delta\ell|$ between the two definitions at
    fixed $\theta$ is \verify{tests/acceptance/results/d401_004.json}{max_abs_dloglik}{,.0f}{196,269} nats, nearly all
    of it $M=7$'s error on \code{normal_frac_sd}, not the definition (divergences/register.toml, D-401).
\end{itemize}

\subsection{Known issues}\label{sec:het-issues}

\begin{itemize}[nosep]
  \item FR-009 fails as written (D-412): gen-1's heterogeneous law is not within 0.1 nats of the converged mixture on 4
    v6 cohorts above gen-1's \code{HET_N_CAP}; the limit is stated for the tester. It also covers analyses that pool
    replicate cohorts past $n=20{,}000$ (Naveh 2026-11-05): at S5's configuration gen-1's law is
    \verify{specs/004-samplers/review/evidence/sr_naveh_004/s5_binned_kl.json}{kl_per_individual}{.2e}{4.70e-06} nats per individual from the converged
    mixture.
  \item HR-3 and HR-4 read \code{fail} (above); their cells are listed in their results.
  \item No calibrator consumes engine heterogeneity until spec 006. HR-2 shows that the gradient in $\sigma$ is
    correct; whether $\sigma$ can be told apart from $\eps$ is spec 006's to measure, not a claim made here.
\end{itemize}

% Section for specs/004-samplers (T133; FR-020; research R13-R18, R20, R21, R24-R30). Paths are relative to
% src/sraging/ unless they start with tests/ or tools/. SRK citations are to the read-only oracle snapshot
% tests/oracle/v033/SRtools_kappa/SRK/. Every measured number carries a \verify marker (macros.tex).
\section{The Langevin engine (\texttt{sraging.langevin})}\label{sec:lgv}

\subsection{Overview}\label{sec:lgv-overview}

\code{sraging.langevin.simulate} draws a seeded sample of $n$ first-passage times from the SR model's stochastic
differential equation, records it through gen-1's observation transform, and attaches an error statement at the
sample's own $n$ \src{langevin/sample.py:simulate, run}. It replaces SRK's Langevin samplers (\code{getSrHetro},
\code{getISR}), re-implemented from SRK's code and compared with it by distribution, never edited (FR-010). numba is an
optional extra that the FP path never imports (FR-014). Every option is named; SRK's (\code{nsteps},
\code{time_step_multiplier}, \code{step_size}, \code{parallel}, \code{n_jobs}, \code{seed}, \ldots) are refused by
name.

\subsection{The kernels and the two walls}\label{sec:lgv-kernels}

Each individual $i$, with its own parameters, follows
\begin{equation}\label{eq:lgv-sde}
  \dd X=\Big(\eta\,t-\frac{\beta X}{X+\kappa}\Big)\dd t+\sqrt{2\eps}\,\dd W,\qquad X(0)=0,
\end{equation}
killed when $X\ge\xc$ \src{langevin/kernels.py:run_chunk}. One step of length $h$ (the step $dt$, or $dt/2$ on the
coupled rung) from age $kh$ is
\begin{equation}\label{eq:lgv-step}
  \tilde x=x+h\Big(\eta\,(kh)-\frac{\beta x}{x+\kappa}\Big)+\sqrt{2\eps h}\,\xi,\qquad
  x'=\begin{cases}\max(\tilde x,0)&\text{clamp (default)}\\ |\tilde x|&\text{reflect}\end{cases}
\end{equation}
\src{langevin/kernels.py:step}. The age $kh$ is computed from $k$, never accumulated, and there is no outer/inner
split. The kill is $x'\ge\xc$ at the step's end; with the bridge kernel (the default) also the Brownian-bridge crossing
probability
\begin{equation}
  p=\exp\!\Big(-\frac{2(\xc-x)(\xc-x')}{2\eps h}\Big)
\end{equation}
when both ends are below $\xc$ and $x>0$, evaluated (and one uniform drawn from the bridge stream) only while
$2(\xc-x)(\xc-x')<30\cdot2\eps h$, SRK's $e^{-30}$ cutoff, after the wall \src{langevin/kernels.py:bridge_p}. The death
time is the end of the kill step. A competing external death time $t_{\rm ext}$, drawn outside the kernel, stops the
path at $t_{\rm ext}$ if it comes first; a survivor at $t_{\max}$ is censored there.

\paragraph{The walls.} The default \code{wall="clamp"} is SRK's $\max(0,\cdot)$ (\src{SR_hetro.py:383, 420, 481}), the
projected Euler scheme: a step that crosses 0 lands on 0 exactly, so young-age damage carries a lump of exact zeros
that the FP density does not have, and the scheme is weak order $\frac12$ at the wall. \code{wall="reflect"} is the
symmetrized Euler scheme $|\cdot|$ (Bossy, Gobet and Talay 2004), weak order 1 under its conditions. The wall is
recorded in every sample's provenance. Keeping the clamp as the default is register D-410 (ported as is; OD-12).
Which wall needs the finer step depends on the configuration. In P1's ladders (one key per rung, $n=25{,}000$, KS
against FP, crit \verify{tests/acceptance/results/p1_langevin_004.json}{cells[8].ladder[2].crit_95}{.4f}{0.0086}), reflect was further from FP than the clamp at coarse
steps for E.~coli and the human (E.~coli homogeneous at $m=4$: clamp \verify{tests/acceptance/results/p1_langevin_004.json}{cells[8].ladder[2].ks_D}{.4f}{0.0065}, reflect \verify{tests/acceptance/results/p1_langevin_004.json}{cells[9].ladder[2].ks_D}{.4f}{0.0345}), about equal for Jack Russell (homogeneous, $m=16$:
\verify{tests/acceptance/results/p1_langevin_004.json}{cells[4].ladder[4].ks_D}{.4f}{0.0399} and \verify{tests/acceptance/results/p1_langevin_004.json}{cells[5].ladder[4].ks_D}{.4f}{0.0388}), and closer for mice\_F
(homogeneous, $m=2$: \verify{tests/acceptance/results/p1_langevin_004.json}{cells[0].ladder[1].ks_D}{.4f}{0.0109} and \verify{tests/acceptance/results/p1_langevin_004.json}{cells[1].ladder[1].ks_D}{.4f}{0.0045})
and the README's $\theta$. The advice is therefore to halve $dt$ until the sample reads clean; reflect's advantage is
near the wall at a converged step (E.~coli's homogeneous reflect met its bar at
$m=\verify{tests/acceptance/results/p1_pooled_004.json}{frozen.ecoli|hom|reflect.m}{d}{133}$, its clamp at no step tried).

\subsection{The fixes and their register entries}\label{sec:lgv-fixes}

Each fix is a sampler entry in \code{divergences/register.toml}, measured by phase \code{plants} against SRK's own code
through an oracle worker, with a planted case that can fail on SRK beside its twin (research R14; the result
\code{tests/acceptance/results/plants_004.json},
\verify{tests/acceptance/results/plants_004.json}{status}{}{pass}). The impact is $\sup|\Delta S|$ at
$n=25{,}000$ unless stated.

\begin{longtable}{@{}p{1.2cm}p{5.6cm}p{4.4cm}>{\raggedright\arraybackslash}p{3.0cm}@{}}
\toprule
entry & SRK (legacy) & sraging & $\sup|\Delta S|$\\
\midrule\endhead
D-402 & B1: the bridge kernel froze the drift's age at $t_j$ for all $m$ substeps (\src{SR_hetro.py:478}) & one step,
  age $kh$ from $k$ & \verify{tests/acceptance/results/plants_004.json}{plants[0].sup_abs_dS}{.3g}{0.0222}\\
D-403 & B2: kernels step $t_{\rm end}/((N-1)m)$, the grid helpers mean $t_{\rm end}/(Nm)$ (\src{stepsize.py:52},
  \src{SR_hetro.py:108}) & one $dt$, realised bit for bit &
  \verify{tests/acceptance/results/plants_004.json}{plants[1].sup_abs_dS}{.3g}{0.013}\\
D-404 & B4: $n\bmod n_{\rm jobs}$ individuals dropped and padded as censored (\src{SR_hetro.py:434-435}) & every one
  of $n$ drawn & \verify{tests/acceptance/results/plants_004.json}{plants[2].sup_abs_dS}{.3g}{0.013} at
  $n=\verify{tests/acceptance/results/plants_004.json}{plants[2].n}{,}{1,003}$\\
D-405 & B11: $\xc\,\mathcal N(1,\sigma)$ unfiltered (\src{SR_hetro.py:471}) & the truncated normal
  (Section~\ref{sec:het-def}) & \verify{tests/acceptance/results/plants_004.json}{plants[3].sup_abs_dS}{.3g}{0.024}
  at $\sigma=0.5$\\
D-406 & numba's generator unreachable by the seed; workers reseeded from the OS & per-individual keys
  (Section~\ref{sec:lgv-rng}) & reproducibility, not $\Delta S$\\
D-407 & a first-order kill, probability $h\,\delta$ per substep (\src{SR_hetro.py:462}) & an exact competing
  exponential & \verify{tests/acceptance/results/plants_004.json}{plants[5].sup_abs_dS}{.3g}{0.0174}\\
D-408 & deaths recorded at $(j+1)\Delta t$ on the outer grid (\src{SR_hetro.py:511-513}) & $\lfloor t/\delta\rfloor\delta$
  (Section~\ref{sec:lgv-obs}) & \verify{tests/acceptance/results/plants_004.json}{plants[6].sup_abs_dS}{.3g}{0.0242}\\
D-409 & \code{getSrHetro} heterogeneous in $\xc$ (0.2) by default (\src{SR_hetro.py:145, 148}) & homogeneous unless
  named & \verify{tests/acceptance/results/plants_004.json}{plants[7].sup_abs_dS}{.3g}{0.0611}\\
D-411 & \code{getISR} stored 0 for the dead, at the outer step within $dt/2$ of the age & NaN iff dead, ages on the
  $dt$ grid (Section~\ref{sec:lgv-paths}) & \verify{tests/acceptance/results/plants_004.json}{plants[9].sup_abs_dS}{.3g}{0.00691}
  ($\sup|\Delta F|$ of $X$)\\
\bottomrule
\end{longtable}
B1's plant is like for like (gate C3): SRK at $(\Delta t,m=16)$ against SRK at $(\Delta t/16,m=1)$, both on the
$\Delta t$ grid, so only the frozen age differs; HR-9 reads
\verify{tests/acceptance/results/plants_004.json}{hr9.status}{}{pass}. Each entry names its regression test, which
fails on SRK's behaviour and passes on the fix.

\subsection{The random streams and the seed}\label{sec:lgv-rng}

A sample depends only on $(\text{key},i)$ \src{langevin/rng.py}. Individual $i$'s stream is xoshiro256**, seeded by
four SplitMix64 outputs from $\mathrm{hash64}(\text{key},i,\text{label})$; the key has 63 bits. The label names the
stream: the path's normals (empty label), the bridge's uniforms (\code{"bridge"}), each heterogeneous parameter
(\code{"het:<parameter>"}) and the external hazard (\code{"hazard"}), so heterogeneity and the hazard never shift the
path stream. Uniforms are 53-bit; normals come in Marsaglia polar pairs. Two implementations, numba's
(\code{next_u64}, \code{uniform}, \code{normal_pair}) and a pure NumPy reference (\code{RefStream}), agree bit for bit
on the first 1,000 draws. What follows: a sample is the same for any worker count and chunk size, and the first $n'$
individuals of a draw of $n>n'$ are the draw of $n'$. HR-8 checks the bytes at $W=1$, 7 and 15 with
$n=\verify{tests/acceptance/results/hr8_004.json}{config.n}{,}{25,013}$, not divisible by any of them:
\verify{tests/acceptance/results/hr8_004.json}{status}{}{pass}.

\subsection{gen-1's observation transform}\label{sec:lgv-obs}

The parameters of individual $i$ are $P\,s_i$ with $s_i=1+\sigma z_i$ from gen-1's own inverse CDF of the truncated
normal (\code{generator._het_z}), so no draw is $\le0$; the external death time is gen-1's \code{_hazard_times} with
rate $h_{\rm ext}$ \src{langevin/sample.py:draw_parameters}. The raw times are the kernel's, after the competing hazard.
The recorded sample is gen-1's own functions, called through the module, so the cohort generator and the Langevin share
one transform (C8) \src{langevin/sample.py:observe}:
\begin{enumerate}[nosep]
  \item the censoring time $C$ from \code{_censor_time} with the FP observed survival, when \code{censor_frac} $>0$;
  \item the recording \code{_record}: a death at $t<C$ is recorded at $\lfloor t/\delta\rfloor\delta$ (first bin at 0),
    $t\ge C$ censored at $C$;
  \item gen-1's sort, then the window \code{_apply_time_range}.
\end{enumerate}
The default $t_{\max}$ is the age where the FP observed survival falls below $0.01/n$; a survivor there is censored at
$t_{\max}$ and counted, with \code{LANGEVIN_SURVIVORS_AT_TMAX}.

\subsection{The error statement}\label{sec:lgv-error}

Every sample carries an \code{ErrorStatement} at its own $n$ \src{langevin/sample.py:ErrorStatement}:
\begin{itemize}[nosep]
  \item \textbf{KS} of the raw times (before recording, so binning, censoring and windows are never read as model
    disagreement) against the FP observed survival, its node count by R26's curve check
    (Section~\ref{sec:het-ladder}); $D$, ${\rm crit}_{95}=1.358/\sqrt n$, the $p$-value and the verdict
    ``consistent at $n$'' or ``inconsistent at $n$''.
  \item \textbf{The step bias} by a coupled half-step rung: the same streams at $dt/2$, the $dt$ path stepping with
    $(\xi_a+\xi_b)/\sqrt2$ and the $dt/2$ rung with $\xi_a$ then $\xi_b$, so both follow one Brownian path. It reports
    $D=\sup|\hat S_{dt}-\hat S_{dt/2}|$ with a paired bootstrap SE.
\end{itemize}
For a scheme of weak order $p$, $D=b(dt)-b(dt/2)=b(dt)(1-2^{-p})$, so the bias of the user's own sample is
\begin{equation}\label{eq:lgv-bdt}
  \hat b(dt)=\frac{D}{1-2^{-\hat p}},
\end{equation}
with $\hat p=\frac12$ for every configuration ($\hat b=3.41\,D$; the library reads no test result, Naveh
2026-11-05, T142 C4). $\hat b$ is an
estimate from one halving, not a bound: on the README's $\theta$ (the ITP fit, days) under the clamp at $dt=0.5$ and
$n=25{,}000$ it reads
\verify{specs/004-samplers/review/evidence/sr_naveh_004/clamp_vs_reflect.json}{rows[0].bias_at_dt}{.4f}{0.0132} while the sample's own KS distance
from FP is \verify{specs/004-samplers/review/evidence/sr_naveh_004/clamp_vs_reflect.json}{rows[0].ks_D}{.4f}{0.0242}. \code{error="ks"} skips the rung:
the step bias reads NOT MEASURED and the sample is never clean.

\paragraph{The step bias against the crit.} A sample whose $\hat b(dt)$ exceeds its own ${\rm crit}_{95}$ carries
\code{LANGEVIN_STEP_BIAS_ABOVE_CRIT} and is not clean \src{langevin/sample.py:run} (FR-012, Naveh 2026-11-05, T141):
at small $n$ the KS has no power to see the step. On the same $\theta$ under the clamp at $dt=4$ and $n=2{,}000$, the
KS passes ($D=\verify{specs/004-samplers/review/evidence/sr_naveh_004/clean_vs_bias.json}{rows[2].ks_D}{.4f}{0.0290}$ against
${\rm crit}_{95}=\verify{specs/004-samplers/review/evidence/sr_naveh_004/clean_vs_bias.json}{rows[2].crit_95}{.4f}{0.0304}$) with
$\hat b=\verify{specs/004-samplers/review/evidence/sr_naveh_004/clean_vs_bias.json}{rows[2].bias_at_dt}{.4f}{0.1093}$, and the sample is not clean.
The warning is never given under \code{error="ks"}, where the bias is NOT MEASURED. A clean sample therefore certifies
the KS at $n$, an FP reference, no survivor at $t_{\max}$ and a stated bias within the crit; it does not certify a
converged step, since $\hat b$ is not a bound.

\paragraph{The judged rung (gate C1).} The acceptance phases judge the $dt/2$ rung, whose bias is
\begin{equation}\label{eq:lgv-judged}
  \hat b(dt/2)=\frac{D}{2^{\hat p}-1},
\end{equation}
and a row is assessable iff $\hat b(dt/2)\le{\rm crit}/10$. Because P1's noise floor exceeded ${\rm crit}/30$, each
cell's step was frozen once from a measurement at pooled $n$ (\code{p1-pooled}: 18 keys at $n=25{,}000$) and not
re-measured per row \src{tools/acceptance_004.py:frozen_steps}; for example mice\_F's homogeneous clamp step is
$m=\verify{tests/acceptance/results/p1_pooled_004.json}{frozen.mice_F|hom|clamp.m}{d}{170}$
(amendment A1, \code{langevin_004.toml}).

\subsection{Paths and the FP damage density}\label{sec:lgv-paths}

With \code{record_ages} the kernel stores $X_{ij}$, individual $i$'s damage at the end of step
$\mathrm{round}(a_j/dt)$, NaN iff $i$ is dead by $a_j$, read with no draw, so the times are the same with and without
paths \src{langevin/kernels.py:run_chunk; langevin/sample.py:Paths}. The ages must lie on the $dt$ grid
(\code{grid_dt} offers a step that divides them), and the block is capped by \code{max_path_bytes}.

\paragraph{The reference.} \code{damage_density} runs the engine's own march with every step stored, only to the
oldest age, and returns the conditional CDF of damage among the living \src{engine/density.py:damage_density}:
\begin{equation}
  F(x\mid\text{alive},a)=\frac{\sum_kw_kG_k(x/x_{c,k},a)}{\sum_kw_kS_k(a)},
\end{equation}
with $G_k$ node $k$'s living mass below $y$ on its grid, interpolated linearly in $\tau=a/L$, and $x_{c,k}=\xc s_k$
under $\xc$ heterogeneity. The external hazard does not enter: it kills independently of $X$. Its node count comes from
R26's check at $10^{-4}$; its grid error $\sup_x|F_{(N,dt)}-F_{(2N,dt/2)}|$ is recorded beside it.

\paragraph{The censored statistic.} Under the clamp, the lump at $X=0$ is not in the FP density. The test therefore
censors on the left \src{sampler_rules.py:hr10_censoring}: with $p_0$ the share of the living at exactly 0,
\begin{equation}
  u_c=\max(0.01,\,3p_0),\qquad x_c=\max\big(F^{-1}(u_c),\,5\sqrt{2\eps_{\max}dt}\big),
\end{equation}
and $D_c$ is the KS distance above $x_c$. Under reflect the full $D$ is used. The ages are Bonferroni-corrected, and
the coupled rung's $\sup|F_{dt}-F_{dt/2}|$ on the same range is reported beside each. HR-10 reads
\verify{tests/acceptance/results/paths_004.json}{verdicts[0].status}{}{pass} on the clamp
(\verify{tests/acceptance/results/paths_004.json}{verdicts[0].n_valid}{d}{11} tests; E.~coli has no clamp
configuration, amendments A0 and A2) and \verify{tests/acceptance/results/paths_004.json}{verdicts[1].status}{}{pass}
on reflect (\verify{tests/acceptance/results/paths_004.json}{verdicts[1].n_valid}{d}{16} tests).

\subsection{Multi-parameter heterogeneity}\label{sec:lgv-multi}

The Langevin accepts a tuple of \code{Heterogeneity} objects: independent truncated normals, one per parameter,
canonicalised by parameter name \src{langevin/sample.py:HetTuple}. Each parameter draws from its own stream, so a member
at $\sigma=0$ (\code{Heterogeneity.zero}) is $s=1$ exactly with no draw, and a tuple whose other spreads are 0 is the
single-parameter draw bit for bit (the exact reduction). FP's likelihood and curves refuse a tuple. The reference for a
tuple's error statement is the tensor product (Section~\ref{sec:het-tensor}) at $d=2$ or 3 within 20,000 nodes;
otherwise the statement reads ``no FP reference''. A tuple's paths are recorded, not checked (register D-413).

HR-11 reads \verify{tests/acceptance/results/multi_004.json}{status}{}{pass} on its pre-declared re-run at fresh keys (amendment A9): every
one of its six cells is below crit, the largest
$D=\verify{tests/acceptance/results/multi_004.json}{cells[4].D}{.3g}{0.00709}$ (xc+$\eta$+$\epsilon$) against
${\rm crit}=\verify{tests/acceptance/results/multi_004.json}{cells[4].crit}{.4g}{0.01047}$, and the tensor reference built at
$10^{-4}$ agrees with the composite rule to \verify{tests/acceptance/results/multi_004.json}{composite.sup_dS}{.2g}{2.6e-05}. The reductions
are exact, and the above-reference statement (a 4-parameter cell, past $d\le3$: amendment A7) reads ``no FP
reference'' and is not clean. The first run failed one cell, xc+$\eta$ at
$D=\verify{specs/004-samplers/review/evidence/multi_004_first_252ad3a.json}{cells[0].D}{.4g}{0.01112}$, diagnosed as chance (step bias and the reference ruled out;
four fresh draws pass), and compared the composite with the cells' coarser reference (bar crit/10); both runs are
kept.

\subsection{The engine switch}\label{sec:lgv-switch}

\code{population(params, ages, heterogeneity=, engine=, n=, ...)} chooses the engine \src{population.py:population}:
\begin{center}
\begin{tabular}{@{}lp{4.4cm}p{6.6cm}@{}}
\toprule
engine & $n=$ None & $n=$ int\\
\midrule
\code{fp} & the FP curves & gen-1 inside its domain, else refused naming why\\
\code{langevin} & refused: ``pass $n$'' & \code{simulate}\\
\code{auto} & the FP curves (a tuple: ``pass $n$'') & gen-1 inside its domain, else the Langevin\\
\bottomrule
\end{tabular}
\end{center}
gen-1's domain is $\xc$ or no heterogeneity, $n\le$ \code{HET_N_CAP} $=20{,}000$ when heterogeneous, and a valid
\code{CohortSpec} \src{population.py:gen1_spec}. Outside it \code{auto} draws with the Langevin and the warning
\code{ENGINE_AUTO_LANGEVIN} names the reason \src{population.py:auto_engine, P3_RULE}. The boundary is pilot P3's
committed rule (C10): gen-1's domain, not the cost. P3 timed both engines on the human preset with $\xc$ sd 0.2,
the Langevin at the preset's frozen heterogeneous clamp step with its rung, in one process, one timing each: gen-1's
heterogeneous draw took
\verify{tests/acceptance/results/p3_cost_004.json}{boundary.rows[3].cpu_gen1_s}{.0f}{690} CPU-s at
$n=2{,}000$ against the Langevin's
\verify{tests/acceptance/results/p3_cost_004.json}{boundary.rows[3].cpu_langevin_s}{.0f}{127}, and
\verify{tests/acceptance/results/p3_cost_004.json}{boundary.rows[4].cpu_gen1_s}{,.0f}{691} against \verify{tests/acceptance/results/p3_cost_004.json}{boundary.rows[4].cpu_langevin_s}{,.0f}{1,254} at $n=20{,}000$. gen-1 is exact up to its stated grid error, so \code{auto} keeps it wherever it can draw. A sample's $S$ is the Kaplan-Meier estimate at
the ages with Greenwood's 95\,\% bands \src{population.py:kaplan_meier}.

\paragraph{JAX.} Pilot P4 priced a JAX kernel at
\verify{tests/acceptance/results/p4_jax_004.json}{ratio_geomean}{.3g}{5.81} times numba's CPU (geometric mean over
\verify{tests/acceptance/results/p4_jax_004.json}{n_pairs}{d}{20} pairs, threshold
\verify{tests/acceptance/results/p4_jax_004.json}{threshold}{}{1.2}), so the kernels stay in numba.

\subsection{Known issues}\label{sec:lgv-issues}

\begin{itemize}[nosep]
  \item The default clamp is not shown to agree with FP for E.~coli at any affordable step: E.~coli has no clamp
    configuration in S4's or HR-10's families (amendment A0, \code{langevin_004.toml [limits]}).
  \item Paths of a tuple with two or more spreads are not checked (D-413): no tensor damage density in spec 004.
  \item Reflect's coverage: S4's \verify{specs/004-samplers/review/evidence/sr_naveh_004/s4_aggregate.json}{walls.reflect.rows}{d}{8} reflect rows are each preset's
    homogeneous row and its $\xc$ sd 0.2 row; no $\eta$, $\beta$, $\eps$ or $\kappa$ row and no HR-11 cell ran on
    reflect.
\end{itemize}

% Section for specs/005-data-plots-cli (T056; FR-010 to FR-012; research R2, R3; register D-501). Paths are relative to
% src/sraging/. Replaces lifelines' KaplanMeierFitter as SRK Dataset uses it, and SRK's Dataset.getMedianLifetime /
% getSteepness (named, not edited). Measured numbers carry \verify markers against km_005.json or the development
% checks of the sr-naveh review (review/evidence/devchecks_005.json).
\section{The Kaplan--Meier estimator and summary statistics (\texttt{sraging.survival})}\label{sec:surv}

\subsection{Overview}\label{sec:surv-overview}

\code{sraging.survival.km} builds one Kaplan--Meier estimator $\hat S$ with a named 95\,\% band from a
\code{Cohort} or an \code{(ages, events)} pair \src{survival.py:km, _km}; \code{summary} reads the named statistics
from it \src{survival.py:summary}. It replaces the KM that SRK's \code{Dataset} takes from lifelines
(\code{KaplanMeierFitter}) and its \code{getMedianLifetime} and \code{getSteepness}. Every statistic carries its
definition as text, never implied. Left truncation (entry ages) and a fit window are refused by name. The linear
band of \code{population.kaplan\_meier} (spec 004) and the compat KMs are untouched (FR-012, D-007).

\subsection{The estimator}\label{sec:surv-km}

Let $u_1<\dots<u_J$ be the distinct death times, $d(u)$ the deaths at $u$ and $r(u)=\#\{i: t_i\ge u\}$ the number at
risk. A censoring at exactly $u$ is at risk at $u$, so \emph{deaths come first at a tie}:
\begin{equation}\label{eq:surv-km}
  \hat S(t)=\prod_{u_j\le t}\frac{r(u_j)-d(u_j)}{r(u_j)}
  =\exp\sum_{u_j\le t}\big[\log(r-d)-\log r\big],\qquad \hat S=1\ \text{before } u_1 .
\end{equation}
$r(u)$ is $n$ minus the number of ages below $u$, found with \code{np.unique} and \code{np.searchsorted}
\src{survival.py:_core}; the cost is $O(n\log n)$ and the result does not depend on row order. The sum-of-logs form is
lifelines 0.30.3's arithmetic, so $\hat S$ is meant to equal lifelines' bit for bit (HR-3: $|\Delta\hat S|\le10^{-12}$,
the median identical, against lifelines side by side on the cages of the two records and on synthetic ties). A
planted defect, censorings before deaths at a tie, must fail HR-3. Because the arithmetic is lifelines', the
comparison locks it rather than checking it independently.
The recorded run (phase \code{km}) compared the mosquito and fly records' cage cohorts and \verify{tests/acceptance/results/km_005.json}{cells[0].cohorts}{d}{201} synthetic tie cohorts, with a largest
$|\Delta\hat S|$ of \verify{tests/acceptance/results/km_005.json}{cells[1].max_dS}{g}{0}, \verify{tests/acceptance/results/km_005.json}{cells[2].max_dS}{g}{0} and \verify{tests/acceptance/results/km_005.json}{cells[0].max_dS}{g}{0} (mosquito, fly, synthetic), a largest band difference of \verify{tests/acceptance/results/km_005.json}{cells[1].max_dband}{g}{0}, \verify{tests/acceptance/results/km_005.json}{cells[2].max_dband}{g}{0} and \verify{tests/acceptance/results/km_005.json}{cells[0].max_dband}{g}{0}, and \verify{tests/acceptance/results/km_005.json}{cells[1].median_mismatches}{d}{0}, \verify{tests/acceptance/results/km_005.json}{cells[2].median_mismatches}{d}{0} and \verify{tests/acceptance/results/km_005.json}{cells[0].median_mismatches}{d}{0} median mismatches. The planted tie defect
gives $\hat S(2)=\verify{tests/acceptance/results/km_005.json}{planted.planted}{.3f}{0.533}$ against the correct \verify{tests/acceptance/results/km_005.json}{planted.ours}{.3f}{0.600}, and fails as it must.
HR-3 has no reference for the other statistics. In a development check the restricted mean matched lifelines'
\code{restricted\_mean\_survival\_time} to relative \verify{specs/005-data-plots-cli/review/evidence/devchecks_005.json}{n5.rmst_rel}{.1e}{6.4e-16} and the median's CI matched on every cage of the two records
(\verify{specs/005-data-plots-cli/review/evidence/devchecks_005.json}{n5.ci_mismatch}{d}{0} mismatches); the quartile ages, the CV and the conditional statistics are compared with nothing.

\paragraph{The band.} With $V(t)=\sum_{u_j\le t} d/\big(r(r-d)\big)$ (a term with $r=d$ adds $0$) and $z$ the
$1-(1-\text{level})/2$ normal quantile, the $\log(-\log)$ Greenwood band is
\begin{equation}\label{eq:surv-band}
  \hat S_{\mp}(t)=\exp\!\Big(-\exp\!\Big[\log(-\log\hat S)\;\mp\;\frac{z\sqrt{V}}{\log\hat S}\Big]\Big),
\end{equation}
the lower limit takes $\mp=-$ and the upper $\mp=+$ \src{survival.py:_band}. As in lifelines,
the band is $[0,0]$ where $\hat S=0$ and $[1,1]$ before the first death (a NaN limit is set to $1$). HR-3 holds the
band to $10^{-10}$. The band's name \code{log(-log) Greenwood} is stored on the result.

\paragraph{Conditioning.} \code{condition\_on=}$t_0$ rebuilds the KM from the individuals observed beyond $t_0$, which
is $\hat S(t)/\hat S(t_0)$ for $t\ge t_0$; the result records $(t_0,\hat S(t_0))$ and every statistic's definition
text says it is conditional. It is refused when no one is observed beyond $t_0$, and on a KM already conditioned.

\subsection{The named statistics}\label{sec:surv-stats}

\paragraph{Median.} The first-crossing median is $\min\{t:\hat S(t)\le\tfrac12\}$ (lifelines' median), $\infty$ if
$\hat S$ never reaches it \src{survival.py:_first}. A crossing at an exact rational level follows the \emph{computed}
$\hat S$, which can sit one ulp either side, as lifelines' does. Its CI is the pair of ages where the band's lower and
upper limits first reach $\tfrac12$.

\paragraph{Quartile ages and steepness.} Two definitions of the age $t(q)$ at which $\hat S$ reaches $q$
\src{survival.py:quartile_age}:
\begin{align}
  \text{first-crossing (default):}\quad & t(q)=\min\{t:\hat S(t)\le q\},\\
  \text{nearest:}\quad & t(q)=\text{the first age of the KM timeline minimising }|\hat S(t)-q|,
\end{align}
the timeline being $\{$start$\}$ and every observed age (lifelines'), which is SRK \code{Dataset.getSteepness}'s
choice. The IQR steepness is
\begin{equation}\label{eq:surv-steep}
  \sigma_{\rm IQR}=\frac{t_{1/2}}{t(0.25)-t(0.75)}=\frac{-t_{1/2}}{t(0.75)-t(0.25)},
\end{equation}
computed in SRK's arithmetic \src{survival.py:steepness}; it is non-finite when the interquartile range is zero. The
summary names the quartile definition in every entry. Cleaning (Section~\ref{sec:clean}) asks for \code{nearest} by
name, since the stored records use it; the summary's default is first-crossing (clarification 2026-11-08).

\paragraph{Restricted mean and CV.} With horizon $\tau$ (default the last observed age) and start $t_0$ (0 unless
conditioned),
\begin{equation}
  m_1=t_0+\int_{t_0}^{\tau}\hat S\,\dd t,\qquad m_2=t_0^2+\int_{t_0}^{\tau}2t\,\hat S\,\dd t,\qquad
  \mathrm{CV}=\frac{\sqrt{\max(m_2-m_1^2,0)}}{m_1},
\end{equation}
integrated exactly on the step function \src{survival.py:_moments}. They are labelled \emph{restricted} when
$\hat S(\tau)>0$, in which case $m_1$ is a lower bound on the mean.

\subsection{Assumptions and limits}\label{sec:surv-limits}

Right censoring only, noninformative; entry ages and windows are refused. The band is pointwise, not simultaneous.
The mean and CV are restricted to the observed range. Every option is validated and refused by name
(\code{UnsupportedOptionError}), including a \code{level} outside $(0,1)$, an unknown \code{quartiles}, and a
\code{horizon} at or before the start.

% Section for specs/005-data-plots-cli (T056; FR-007 to FR-009; research R6; Naveh's Dc). Paths relative to src/sraging/.
% Checked by HR-8 (round trips, log-likelihood invariance, the factor table through SRK getTheta) and HR-7.
\section{Time units (\texttt{sraging.units})}\label{sec:units}

\subsection{Overview}\label{sec:units-overview}

\code{sraging.units} makes the time unit explicit: \code{convert} returns a parameter set or a cohort in another unit,
\code{with\_unit} labels an object whose unit is not yet named, and nothing converts unless \code{convert} is called
\src{units.py:convert, with_unit}. It replaces SRK's \code{getTheta(name, time\_unit=)} factor table and its
\code{rescaleTheta}, which apply the same table (\code{rescaleTheta} is in no frozen snapshot, HR-8).

\subsection{The unit table}\label{sec:units-table}

\code{UNITS} is the closed list of days per unit \src{units.py:UNITS}:
\begin{center}\begin{tabular}{lrrrrr}
\toprule
unit & hours & days & weeks & years & generations\\
\midrule
days per unit $D$ & $1/24$ & $1$ & $7$ & $365$ & $1/8$\\
\bottomrule
\end{tabular}\end{center}
A year is 365 days and a generation is $1/8$ day (Naveh's Dc; SRK's table has the same values, its $3/24$ being
$1/8$). The generation is SRK's fixed three-hour convention for its yeast and \textit{E.~coli} presets, not a property
of the user's organism. \code{convert} refuses any other unit by name, months included (\code{UnsupportedOptionError});
a \code{time\_unit} given to \code{read\_cohorts} or \code{Cohort.from\_arrays} is a label, checked only when
converting. \code{convert} also refuses
\code{time\_unit="input units"} as a source, the default label of an object that has not been told its unit: name it
first with \code{with\_unit(obj, unit)}, which refuses an object that already has a unit.

\subsection{How the parameters transform}\label{sec:units-theta}

Going from unit $a$ to unit $b$, write $s=D_b/D_a$ \src{units.py:factor}; a quantity in $a$-units, multiplied by $s$
(rate) is in $b$-units. A time $T$ in $a$ is $T/s$ in $b$. The model $\dd X=(\eta t-\beta X/(X+\kappa))\dd t
+\sqrt{2\eps}\,\dd W$, killed at $x_c$, has $\eta$ of dimension $[x]\,[t]^{-2}$, $\beta$ of dimension $[x]\,[t]^{-1}$ and $\eps$ of
dimension $[x]^2[t]^{-1}$; with $X$ itself unchanged,
\begin{equation}\label{eq:units-theta}
  \eta'=\eta\,s^{2},\qquad \beta'=\beta\,s,\qquad \eps'=\eps\,s,\qquad \hext'=\hext\,s,\qquad \kappa'=\kappa,\quad
  x_c'=x_c ,
\end{equation}
as built by the \code{SRParams} constructor with \code{time\_unit=to} \src{units.py:convert}. For example days to
weeks: $\eta\times49$, $\beta,\eps,\hext\times7$. The dimensionless groups $\reta,\rbeta,\reps$ at the grid length
$L$ are unchanged when $L$ is converted with the ages ($L'=L/s$), so the converted problem is the same problem. For a
$\delta$ cohort, whose terms are probabilities of intervals, the log-likelihood is invariant up to the numerical grid.
A continuous cohort's terms are densities: each death's term gains $\ln s$, the Jacobian of $t'=t/s$, so its
log-likelihood shifts by $n_{\rm deaths}\ln s$. $\hat\theta$ and likelihood ratios are unchanged, but log-likelihoods
and AIC must not be compared across units.

\paragraph{Presets' external hazard.} SRK's preset tables store $\mathrm{ExtH}=-\ln h$ in the preset's own unit,
SRK's \code{getTheta} 'auto' unit (years for humans, dogs and cats: every preset that has an ExtH), and
\code{getTheta} returns it unscaled in every unit. \code{presets.get(..., ExtH=True)} therefore returns
$\hext=e^{-\mathrm{ExtH}}/D_{\rm auto}$ per day and names the unit in \code{info.ext\_h\_unit}
\src{presets.py:get}; HR-7 compares $\hext$ converted back to that unit with $e^{-\mathrm{ExtH}}$.

\paragraph{Cohorts.} Ages, $\delta$ (scalar or per-individual), entry ages and the window are divided by $s$
($\times D_a/D_b$) and the cohort is rebuilt through \code{Cohort.from\_arrays}, so every validation re-runs
\src{units.py:convert}. Events are unchanged.

\subsection{Acceptance and limits}\label{sec:units-accept}

HR-8 pre-registers: round trips within relative $10^{-15}$; $\mathrm{convert}(\mathrm{get}(\text{name}),\text{to}=u)$
against SRK \code{getTheta(name, time\_unit=u)} from the frozen snapshot, for every preset and the five units, within
relative $10^{-15}$; and a log-likelihood tolerance set by the P3 pilot (about $10\times$ the measured change under
conversion, floor $10^{-12}$). HR-7 adds the preset table in days exactly and in the other units within $10^{-15}$.
The recorded run (phase \code{units}) on a \verify{tests/acceptance/results/units_005.json}{cohort.n}{d}{200}-individual cohort (\verify{tests/acceptance/results/units_005.json}{cohort.n_deaths}{d}{172} deaths) gave a worst $|\Delta\log L|$ of \verify{tests/acceptance/results/units_005.json}{worst_abs_dloglik}{.3g}{9.55e-12} under joint
conversion from days to weeks and to years, and the frozen tolerance is \verify{tests/acceptance/results/units_005.json}{loglik_tolerance}{.3g}{9.55e-11}, ten times P3's value. The run is the same
computation as P3, so this clause is a regression lock; the evidence is the size, not the pass. Round trips returned parameters to relative \verify{tests/acceptance/results/units_005.json}{round_trips.params_rel}{.2g}{1.4e-16} and the cohort to \verify{tests/acceptance/results/units_005.json}{round_trips.cohort_rel}{.2g}{2.7e-16}.
The factor table matched \code{getTheta} on \verify{tests/acceptance/results/units_005.json}{factors.cells}{d}{175} cells with a largest relative difference of \verify{tests/acceptance/results/units_005.json}{factors.max_rel}{.2g}{2.2e-16}, and the planted wrong factor is caught.
The conversion is arithmetic on stored floats, so repeated conversions accumulate rounding of order $10^{-16}$ per
step; it does not change $\kappa$ or $x_c$, which are in the units of $X$.

% Section for specs/005-data-plots-cli (T056; FR-013 to FR-017; research R4, R5; register D-501, D-503, D-504).
% Paths relative to src/sraging/. SRK/lab citations are to the legacy brief docs/planning/spec005_brief_legacy.md.
% Measured numbers carry \verify markers against the acceptance results (exit_005, noise_005) or the development
% checks of the sr-naveh review (review/evidence/devchecks_005.json).
\section{MAD cage cleaning and its noise baseline (\texttt{sraging.cleaning})}\label{sec:clean}

\subsection{Overview}\label{sec:clean-overview}

\code{clean} judges every cage against the other cages of its cohort and returns a report; \code{drop} removes the
flagged cages by naming that report \src{cleaning.py:clean, drop}. Cleaning \emph{flags by default} and \code{k} and
\code{seed} are required; removal is the separate call \code{drop}, which refuses a report not computed on the same
set (D-504). The rule is the lab's \code{cage\_cleaning.clean\_by\_mad} (mosquito reanalysis, fly redo), reimplemented
formula for formula; its score KM replaces SRK \code{Dataset}'s, and \code{noise\_baseline} replaces the fly redo's
\code{null\_drops} and \code{null\_by\_group} constructions.

\subsection{The rule}\label{sec:clean-rule}

\paragraph{Scores.} Cage $i$ of a cohort is scored by $(x_i,y_i)=$ (KM median, IQR steepness), from
Section~\ref{sec:surv}: median by first crossing, quartiles \code{nearest} (SRK \code{Dataset}'s, which the stored
records use), both named in the report \src{cleaning.py:cage_scores, definitions}.

\paragraph{Geometry.} On one axis with finite values $v$: centre $c=\mathrm{med}(v)$, scale
$\sigma=1.4826\,\mathrm{MAD}(v)=1.4826\,\mathrm{med}|v-c|$, and radius
\begin{equation}\label{eq:clean-radius}
  \rho=\max\big(k\sigma,\;k\,\varphi\,|c|\big),
\end{equation}
the floor $k\varphi|c|$ taking over when it is larger or when $\sigma$ is not finite; $\rho=\infty$ if it is still not
positive (every cage tied, no floor) \src{cleaning.py:axis_geometry}. The floors $\varphi=(\varphi_{\rm ml},
\varphi_{\rm st})$ are pooled over the cohorts with at least \code{min\_cages}$=10$ cages as
$\varphi=\mathrm{med}_{\rm cohorts}\,|\sigma/c|$ per axis, or given \src{cleaning.py:pooled_floors}; if no cohort
qualifies and none are given, cleaning is refused.

\paragraph{Distance and flag.} With more than two cages the geometry of cage $i$ uses the \emph{other} cages
(leave-one-out); with two or fewer, the whole group (\code{cage\_cleaning.py:232}). With $u_i=(x_i-c_x)/\rho_x$ and
$v_i=(y_i-c_y)/\rho_y$,
\begin{equation}\label{eq:clean-dist}
  D_i=\sqrt{u_i^2+v_i^2}\quad(\text{or }|u_i|\text{ for }\code{axes="median"}),\qquad
  \text{flagged}\iff D_i>1,
\end{equation}
an axis-aligned ellipse, not oriented and not the bootstrap ellipse of SRK's \code{clean\_DatasetCollection}
\src{cleaning.py:judge}. A non-finite score on an axis the rule uses is refused, naming
the cage and axis; the legacy code kept such a cage at distance 0 (D-503).

\subsection{The noise baseline}\label{sec:clean-noise}

Flags are expected even without a real outlier, so each cohort row carries the flags the same rule gives exchangeable
cages whose scores scatter under two named nulls, always both (FR-016). The scatter is Gaussian, independent on the two
axes and independent of the number of animals per cage; it is not the sampling noise of a KM median. For a cohort of $m$ cages, draw $R$ (default 4000) samples of $m$
scores per axis $1+s\,\xi$, $\xi\sim N(0,1)$, with the same $k$, floors and leave-one-out, and report the flag rate
per cage and the expected flags $=$ rate$\times m$ \src{cleaning.py:noise_baseline, _flags_batch}. The generator is
\code{numpy.random.default_rng} seeded from \code{(seed, m, null)} so each null is reproducible and independent.
\begin{itemize}[nosep]
  \item \textbf{floor null:} $s=\varphi$ on each axis.
  \item \textbf{own-spread null:} $s$ is the leave-one-out corrected spread \emph{pooled} over a group of cohorts:
    those sharing the value of the design column \code{spread\_by}, or every cohort of the set by default (Naveh,
    2026-11-09; the fly redo's \code{group\_spreads}) \src{cleaning.py:pooled_spread, clean}. With $e_i=x_i/\mathrm{med}(x_{-i})-1$ the leave-one-out
    relative deviations of a cohort of $m$ cages \src{cleaning.py:loo_relative_deviations},
    \begin{equation}\label{eq:clean-own}
      s=1.4826\,\mathrm{MAD}\big(\{e_i/r(m_c)\}_{i\in c,\;c\in\text{group}}\big),
    \end{equation}
    where $r(m)$ is the robust sd of $x_i-\mathrm{med}(x_{-i})$ over $m$ i.i.d.\ $N(0,1)$ cages (20\,000
    replicates). $r(m)>1$ ($r(3)\approx1.22$): a leave-one-out deviation carries the noise of the others' median too,
    and dividing by $r(m)$ removes that inflation \src{cleaning.py:loo_correction}. A single cohort cannot be its own
    group: with three cages two ties give $s=0$, a null with no scatter. A pooled $s$ that is $0$ or not finite is
    reported as not measurable, with no expected flags.
\end{itemize}
Neither null is the analytic $e^{-k^2/2}$, which understates the rate at small $m$ (research R5). The reading of how
many flags are noise depends on the null, which is why both are shown.

\subsection{Acceptance and stated limit}\label{sec:clean-accept}

\paragraph{HR-1 (the exit).} The code is compared with the \emph{stored} mosquito and fly cleaning records (read-only,
pinned by sha256) at every recorded $k$ (mosquito 2.5; fly 2.5 and 3.0) with floors pooled over the record's cleaned
cohorts: verdicts identical, median exact, distances and floors within $10^{-12}$ (relative for floors), steepness
within $10^{-12}$, with data parsed with round-trip floats. Four planted defects must each fail first: no
leave-one-out, no floors, $k$ swapped, first-crossing quartiles.
The recorded run (phase \code{exit}) reproduced both records: every verdict identical at each recorded $k$, the
largest distance difference \verify{tests/acceptance/results/exit_005.json}{records[0].per_k[0].max_abs_distance}{g}{0} (mosquito), \verify{tests/acceptance/results/exit_005.json}{records[1].per_k[0].max_abs_distance}{g}{0} and \verify{tests/acceptance/results/exit_005.json}{records[1].per_k[1].max_abs_distance}{g}{0} (fly, $k=2.5$ and 3), and the pooled floors equal to the record's. The four planted defects each fail first;
for the mosquito record, without leave-one-out or without floors the rule's flags differ from the record's.
\paragraph{HR-4 (the noise baseline).} The floor null's flag rate per $(m,k)$ is compared with the fly's
\code{null\_drops.csv} (its cells with multiplier 1) and the own-spread null's expected drops per group and $k$ with
\code{null\_by\_group.csv}, each within 4 binomial standard errors at $R=4000$; two planted nulls (no leave-one-out, no
floors) must fail.
HR-4's group cells pool as \code{clean} does, through the same function. In the recorded run (phase \code{noise})
the floor null's rate at $m=3$, $k=2.5$ is \verify{tests/acceptance/results/noise_005.json}{floor_cells[0].ours}{.3f}{0.091} against the fly's \verify{tests/acceptance/results/noise_005.json}{floor_cells[0].theirs}{.3f}{0.090}, within \verify{tests/acceptance/results/noise_005.json}{se_multiple}{.0f}{4} binomial standard errors, as are all its floor cells and all its groups (phase status \verify{tests/acceptance/results/noise_005.json}{status}{}{pass}).
That run replays pilot P2 with the same seed, so it locks the construction rather than testing it afresh; with another
seed the floor cells stay within \verify{specs/005-data-plots-cli/review/evidence/devchecks_005.json}{n2_floor_seed1.max_z}{.2f}{2.28} standard errors (development check). The planted nulls fail: without leave-one-out the rate is \verify{tests/acceptance/results/noise_005.json}{planted[0].ours}{.3f}{0.032} against \verify{tests/acceptance/results/noise_005.json}{planted[0].theirs}{.3f}{0.090}, and a rule without floors judging scatter of size $\varphi$ flags \verify{tests/acceptance/results/noise_005.json}{planted[1].ours}{.3f}{0.514}, both far outside the 4-SE band.

\paragraph{Stated limit.} The cleaning is tested against \emph{stored records}, not side by side against the old code
(no new oracle, Db). A defect shared by a record and the code that wrote it would not show. Without the records the
exit reports \code{not\_applicable}, never a pass; a synthetic fixture (a planted outlier cage, ties, a three-cage
cohort, a zero-IQR cage) runs everywhere. The score is a pair of KM statistics, so the rule is only as good as those
definitions; changing \code{quartiles} changes which cages flag. The records hold small cohorts of cages judged on
both axes with pooled floors: the two-cage whole-group path, \code{axes="median"}, given floors and \code{min\_cages}
other than 10 are tested on synthetic data only. The fly record's \code{kept} column predates the fly redo's final
choices (its D19: median-only for 2--8 LONG, flag-only for SWITCH), which were not reproduced. The mosquito record's
scores come from cages already shifted (its D2) and spike-censored (D8). Verdicts near the radius are fragile: removing
one animal from a cage flips its own verdict for a few cages of each record (gate C11; not re-measured here).

% Section for specs/006-calibrator-capabilities (T070; FR-029). Paths are relative to src/sraging/ unless they start with
% tests/, tools/ or docs/. Measured numbers carry \verify markers against the acceptance results; a result not yet
% recorded carries a TODO(T073) comment.
\section{Fitting a time range: the window $W$ and $H$ (\texttt{profile\_jax2\_w}, \texttt{profile\_jax2\_h})}\label{sec:win}

\subsection{Purpose and order of steps}\label{sec:win-order}

A fit window $[\mathrm{lo},\mathrm{hi}]$ says which part of the lifetimes a cohort shows. \code{profile_jax2_w} ($W$) fits
it with no external hazard; \code{profile_jax2_h} ($H$) fits it with the external hazard $h$ as a nuisance, as jax2
shipped it. Every other calibrator refuses a window; the other jax2 names do so naming $W$ and $H$
\src{calibrators/profile_jax2.py:WINDOW_FITTERS, fits_window}, and \code{profile_v033} without naming them
\src{calibrators/profile_v033.py:fit}. \code{window_path} turns the cohort into what the frozen
code is given, in this order \src{calibrators/profile_jax2.py:window_path}:
\begin{enumerate}[nosep]
  \item $\mathrm{hi}=\infty$ with no recording resolution $\delta$ is refused by name: the frozen snap cannot place it.
  \item $\mathrm{hi}=\infty$ with $\delta$ is mapped to the first grid point past every record,
    $\mathrm{hi}=(\lfloor a_{\max}/\delta+10^{-9}\rfloor+1)\,\delta$ (\code{HI_INF_RULE}), and the mapping is recorded
    (\code{hi_mapped_from}, \code{hi_rule}).
  \item The tester's \code{apply_window} \src{compat/tester_f94ec2b/observation.py:apply_window}, verbatim: records with
    age below lo are \emph{dropped}, not censored (a censored record would say it was known alive at lo, which is not
    what happened: it was never seen); records with age above hi are censored at hi.
  \item The minimum-deaths rule: if fewer deaths are left to fit than the calibrator has free parameters (4 for $W$, 5
    for $H$), the fit is refused and the count is named. The count is the deaths the frozen \code{prepare} leaves, after
    its edge tie (Section~\ref{sec:win-snap}).
  \item The windowed cohort, with the window (lo, hi used), and the row map: the kept rows' indices in the cohort given.
    The fit, its floor report and its leverage check all see the windowed rows, and the rows they name are mapped back
    to the cohort given through the map \src{calibrators/profile_jax2.py:map_window_rows}.
  \item The frozen \code{prepare(window_on=True)}, whose snap record (the effective window and the edge tie) goes to the
    fit's diagnostics.
\end{enumerate}

\subsection{The frozen snap and the edge tie}\label{sec:win-snap}

With a recording resolution $\delta$ and ages on its grid, the frozen \code{window_snap} puts the window on the grid
\src{compat/jax2_101/observation_model.py:window_snap}:
\begin{itemize}[nosep]
  \item $\mathrm{hi}_{\rm eff}$ is hi floored to the grid;
  \item $\mathrm{lo}_{\rm eff}$ is lo ceiled to the grid; if $\mathrm{lo}\le\delta/2$ then $\mathrm{lo}_{\rm eff}=0$ and
    there is no left truncation (the $\delta/2$ clamp keeps every record);
  \item every record in the bin of $\mathrm{hi}_{\rm eff}$ or later is censored at $\mathrm{hi}_{\rm eff}$.
\end{itemize}
Without $\delta$, or with off-grid ages, the window is used as given and nothing is snapped; an empty snapped window is
also used as given. The snap never moves an age up, never turns a censored record into a death and never changes $n$.
It counts the deaths it censors (\code{n_deaths_censored}) and the records it censors at $\mathrm{hi}_{\rm eff}$
(\code{n_censored_at_hi_eff}).

\paragraph{The edge tie.} Ages are recorded as the floor of the age at death: a death recorded at $\mathrm{hi}_{\rm eff}$
happened in $[\mathrm{hi}_{\rm eff},\mathrm{hi}_{\rm eff}+\delta)$, after the window closes. It is therefore known alive at
$\mathrm{hi}_{\rm eff}$ and no more, and is censored there. On a grid window this equals the engine's rule for a death
recorded at the window close (register D-107). It is a convention of the data, not a rare case: in spec 006's
\verify{specs/006-calibrator-capabilities/review/evidence/windows_counts_006.json}{edge_tie.n_fits}{d}{260} gate fits it fired on \verify{specs/006-calibrator-capabilities/review/evidence/windows_counts_006.json}{edge_tie.fired}{d}{159}, its censored deaths came to more than 1\,\% of the records on \verify{specs/006-calibrator-capabilities/review/evidence/windows_counts_006.json}{edge_tie.above_1pc}{d}{67} and more than
10\,\% on \verify{specs/006-calibrator-capabilities/review/evidence/windows_counts_006.json}{edge_tie.above_10pc}{d}{4}, with a median of \verify{specs/006-calibrator-capabilities/review/evidence/windows_counts_006.json}{edge_tie.median_pc}{.2f}{0.13}\,\% and a maximum of \verify{specs/006-calibrator-capabilities/review/evidence/windows_counts_006.json}{edge_tie.max_pc}{.1f}{46.7}\,\% (\code{window-k2-p01-c1-s1}, where $W$'s $\eta$ and
$\varepsilon$ claims were lower bounds only; \code{review/evidence/windows_counts_006.json}). On v6 every windowed record
has both ends on the $\delta$ grid, so no age is moved and the $\delta/2$ clamp does not act (gate M5,
\code{docs/planning/spec006_seed_gate.md}).

\subsection{The windowed likelihood}\label{sec:win-lik}

The frozen likelihood conditions every at-risk row on being alive at lo \src{compat/jax2_101/fp_jax.py:make_loglik}. With
$\ell_i(\theta)$ the terms of the unwindowed likelihood on the same prepared data (deaths, censorings at their ages, the
edge tie as above), $n$ the rows at risk, and $S(\mathrm{lo})$ the survival at the window's start,
\begin{equation}\label{eq:win-cond}
  \ell_W(\theta)=\sum_i\ell_i(\theta)\;-\;n\,\log\operatorname{clip}\!\big(S(\mathrm{lo}),\,10^{-12},\,1\big),
\end{equation}
with $S(\mathrm{lo})$ replaced by $S(\mathrm{lo})\,e^{-h\,\mathrm{lo}}$ under $H$ (the nuisance hazard acts on the time
before lo too) \src{compat/jax2_101/fp_jax.py:850}. Here lo is $\mathrm{lo}_{\rm eff}$ of the snap. So the windowed value
is $J$'s terms plus one conditioning term; this decomposition is tested at four $\theta$, one near the clip (the terms
bit for bit, the value within a few ulp), with a planted wrong conditioning term failing it \src{tests/unit/test_profile_jax2_w.py:test_ws_windowed_value_is_js_terms_plus_the_conditioning_term}.

\paragraph{The clip.} Where $S(\mathrm{lo})<10^{-12}$ the conditioning term is constant in $\theta$, so it carries no
information about the time before lo there. Spec 006 reports it: the warning \code{WINDOW_S_LO_CLIPPED} is raised when $S(\mathrm{lo})$ at
$\hat\theta$ is below the clip \src{calibrators/profile_jax2.py:s_lo_warning; warnings.py}. The fit itself is the frozen one
and is not changed.

\subsection{$W$ as a composition}\label{sec:win-w}

$W$ is \code{MleProfileJax2W}, version \code{1.0.1+w1}: jax2 1.0.1's class $J$ with \code{WINDOW = True} and the hazard
off, under its own port so that its version never reads as a jax2 release
\src{calibrators/profile_jax2.py:MleProfileJax2W, PORT_W}. It is sraging's composition; jax2 never ran it, and
\code{compat/jax2_101/settings.py} stays verbatim. It is not the default, and it is registered at its own version only.
$H$ is the 1.0.1 class as shipped, whose window switch is already on.

Because no jax2 release ran $W$, it is checked against the oracle snapshot run with the same switches (the \emph{identity}
phase, \code{tools/acceptance_006.py}). The phase compares, side by side in one process on the same host: for the window
path, $W$'s or $H$'s public \code{window_path} and port run against the 1.0.1 oracle's \code{run_fit} on the tester's own
observation, after comparing the two observations first; and $W$ against $J$ on an unwindowed cohort, by $\hat\theta$,
log-likelihood, warning codes and the claims, each float by its hex.
The recorded phase passed (\verify{tests/acceptance/results/identity_006.json}{status}{}{pass}, at 966bbe7): $W$ equals $J$ on \verify{tests/acceptance/results/identity_006.json}{equal_by_rule.R-5}{d}{3}
unwindowed cohorts, and the window path equals the oracle's on \verify{tests/acceptance/results/identity_006.json}{equal_by_rule.R-6}{d}{4} pairs ($W$ and $H$ on two
window cohorts); none of the \verify{tests/acceptance/results/identity_006.json}{pairs}{d}{7} pairs has a differing path.

\subsection{The engine's window}\label{sec:win-engine}

The spec 002 engine has its own window \src{engine/observation.py}, which the fits above do not use. Spec 006 changed it
in two places.
\begin{itemize}[nosep]
  \item \textbf{D-601.} With $\mathrm{hi}=\infty$ no $S(\mathrm{hi})$ term is assembled, by a host-side switch. The value
    was already finite, but \code{jnp.interp} at $\infty$ has a NaN derivative, and with the hazard the unused term adds
    $0\times\mathrm{NaN}$, so the whole gradient was NaN. The value now equals a large finite hi's bit for bit, and the gradient to round-off
    (bit for bit in \code{tests/unit/engine/test_engine_window_006.py}; within $10^{-13}$ relative ($8.5\times10^{-14}$ measured) on the parity phase's two cases,
    where the two compiled programs order the sums differently).
  \item \textbf{D-108 (retired).} Numerics v1 now clips $S(\mathrm{lo})$ at $10^{-12}$, as v033, SRK and jax2 do. Before
    spec 006 its common $\log$ floor gave a spurious maximum where $S(\mathrm{lo})$ underflows: the windowed value was
    $0.0$ against $-34.2$ at a sensible $\theta$ in \code{tests/unit/engine/test_engine_clip_006.py}.
    \code{Evaluator.window_report} reports where the clip binds, with the warning \code{ENGINE_S_LO_CLIPPED}
    \src{engine/evaluator.py:window_report}.
\end{itemize}
D-107 (a death recorded in the bin of hi is censored at hi) is the engine's form of the edge tie of
Section~\ref{sec:win-snap}. D-214 records that jax2's mesh differs from the engine's.

\subsection{What a window cannot tell you}\label{sec:win-limits}

By eq.~(\ref{eq:win-cond}), a record that enters at lo is conditioned on being alive at lo, so the data say nothing about
survival before lo except through $S(\mathrm{lo})$: two parameter values with the same $S(\mathrm{lo})$ and the same
conditional survival after lo fit equally well. A narrow window leaves fewer deaths and constrains fewer combinations of
the parameters, and the minimum-deaths rule (Section~\ref{sec:win-order}) refuses only the case with fewer deaths than
parameters, not a window that is merely short.

\subsection{Gate}\label{sec:win-gate}

Whether the windowed fits cover the truth is the window arm's coverage in spec 006's acceptance run (\code{tools/gate_006.py};
the rules are \code{tests/acceptance/rules_006.toml}). The cohorts are generated, with known truth. Counts that the records
hold only as lists are in \code{specs/006-calibrator-capabilities/review/evidence/windows_counts_006.json}, built from the
records by the script beside it. A1 and A2 split the window arm at $n=25{,}000$: the panel's split, set before the run, not
a measured change (below).

On A1, the \verify{tests/acceptance/results/a1_006.json}{n_cohorts}{d}{196} cohorts with $n<25{,}000$ (\verify{specs/006-calibrator-capabilities/review/evidence/windows_counts_006.json}{truncation.a1_synthetic.eta.n_cohorts}{d}{112} at SR parameter grids and \verify{specs/006-calibrator-capabilities/review/evidence/windows_counts_006.json}{truncation.n_a1_real}{d}{84} planted at real life tables,
\verify{specs/006-calibrator-capabilities/review/evidence/windows_counts_006.json}{truncation.n_a1_human}{d}{22} of them human), $W$'s closed intervals cover the truth for $\eta$ on \verify{tests/acceptance/results/a1_006.json}{R-1.gated.eta.hits}{d}{172} of \verify{tests/acceptance/results/a1_006.json}{R-1.gated.eta.claimed}{d}{175} cohorts and for
$\varepsilon$ on \verify{tests/acceptance/results/a1_006.json}{R-1.gated.epsilon.hits}{d}{158} of \verify{tests/acceptance/results/a1_006.json}{R-1.gated.epsilon.claimed}{d}{160}, above the pre-registered cut (\verify{tests/acceptance/results/a1_006.json}{R-1.gated.eta.cut}{.3f}{0.913}). Its truth witness fails on none (R-3 passes:
\verify{tests/acceptance/results/a1_006.json}{R-3.passes}{}{True}). Against $H$ (R-2), the paired median of $\log_{10}(\text{width}_W/\text{width}_H)$ is at most 0.1 dex: \verify{tests/acceptance/results/a1_006.json}{R-2.eta.vs_H.median_dex}{.3f}{-0.006} dex for
$\eta$ over the \verify{tests/acceptance/results/a1_006.json}{R-2.eta.vs_H.n_paired}{d}{66} cohorts where both claims are closed and $W$'s upper end is not truncated, and \verify{tests/acceptance/results/a1_006.json}{R-2.epsilon.vs_H.median_dex}{.3f}{-0.051} for
$\varepsilon$ over \verify{tests/acceptance/results/a1_006.json}{R-2.epsilon.vs_H.n_paired}{d}{85}. $W$'s fit costs \verify{tests/acceptance/results/a1_006.json}{r.fit_over_j}{.2f}{1.28}, \verify{tests/acceptance/results/a2_006.json}{r.fit_over_j}{.2f}{0.52} and \verify{tests/acceptance/results/realfull_006.json}{r.fit_over_j}{.2f}{0.70}\,$\times$ $J$'s stored CPU on A1, A2 and realfull; with the
lead pair, which a default fit computes, a pilot cohort cost \verify{tests/acceptance/results/pilot_006.json}{r.element_over_j}{.3f}{3.627}\,$\times$, the gate's own evaluations included (the gate's
price \code{r_lead}).

\paragraph{A stale $\hat\theta$, more frequent on windows.} $J$ 1.0.1 keeps $\hat\theta$ when a profile walk finds a point
better by less than 1.92 nats (``stale: minor''), and stores that point in \code{diagnostics.estimate.best_point}. So
$\hat\theta$ can sit up to 1.92 nats below the best point the fit saw. The claims contain that point: $J$ widens each claim
to its union with the extents of every evaluated point within $\mathrm{thr}/2=1.92$ nats of the best log-likelihood
\src{compat/jax2_101/claims.py:sc_widen_claim}. The gate's R-4 asks that $W$'s $\hat\theta$ score within 0.1 nats of $H$'s
stored fit at $H$'s smallest $h$. It failed on \verify{specs/006-calibrator-capabilities/review/evidence/windows_counts_006.json}{r4.a1.failures}{d}{15} of A1's \verify{tests/acceptance/results/a1_006.json}{R-4.n}{d}{196} cohorts (largest gap \verify{tests/acceptance/results/a1_006.json}{R-4.max_gap}{.2f}{1.80} nats), \verify{specs/006-calibrator-capabilities/review/evidence/windows_counts_006.json}{r4.a2.failures}{d}{3} of A2's \verify{tests/acceptance/results/a2_006.json}{R-4.n}{d}{44} (\verify{tests/acceptance/results/a2_006.json}{R-4.max_gap}{.2f}{0.54}) and
\verify{specs/006-calibrator-capabilities/review/evidence/windows_counts_006.json}{r4.realfull.failures}{d}{2} of realfull's \verify{tests/acceptance/results/realfull_006.json}{R-4.n}{d}{20} (\verify{tests/acceptance/results/realfull_006.json}{R-4.max_gap}{.2f}{0.81}), every gap below 1.92 nats. The mechanism is $J$'s, and the window makes it more frequent: on A1's
\verify{specs/006-calibrator-capabilities/review/evidence/windows_counts_006.json}{r4.j_vs_w_synthetic.n}{d}{112} SR-grid cohorts $W$'s gap exceeds 0.1 nats on \verify{specs/006-calibrator-capabilities/review/evidence/windows_counts_006.json}{r4.j_vs_w_synthetic.w_above}{d}{11}, and $J$'s own gap on its stored, unwindowed fits of their pure twins
on \verify{specs/006-calibrator-capabilities/review/evidence/windows_counts_006.json}{r4.j_vs_w_synthetic.j_above}{d}{3} (\verify{specs/006-calibrator-capabilities/review/evidence/windows_counts_006.json}{r4.j_vs_w_synthetic.both_above}{d}{1} both; Fisher $p=\verify{specs/006-calibrator-capabilities/review/evidence/windows_counts_006.json}{r4.j_vs_w_synthetic.fisher_p}{.3f}{0.050}$). A Mac refit of one failing cohort per cell
(\code{review/evidence/r4_refits_006.json}) reads a better point in $W$'s own self-consistency check, ``stale:
\verify{specs/006-calibrator-capabilities/review/evidence/r4_refits_006.json}{refits.realfull-real-drosophila_217-c0-s0.stale}{}{minor}'', at about R-4's gap: \verify{specs/006-calibrator-capabilities/review/evidence/r4_refits_006.json}{refits.realfull-real-drosophila_217-c0-s0.self_consistency_gap_nats}{.2f}{0.81} against \verify{specs/006-calibrator-capabilities/review/evidence/r4_refits_006.json}{refits.realfull-real-drosophila_217-c0-s0.r4_gap_wexac}{.2f}{0.81} nats on realfull, \verify{specs/006-calibrator-capabilities/review/evidence/r4_refits_006.json}{refits.window-real-Sweden_M_1910_hetro-c1-s1.self_consistency_gap_nats}{.2f}{0.58} against \verify{specs/006-calibrator-capabilities/review/evidence/r4_refits_006.json}{refits.window-real-Sweden_M_1910_hetro-c1-s1.r4_gap_wexac}{.2f}{0.54} on A2, and \verify{specs/006-calibrator-capabilities/review/evidence/windows_counts_006.json}{refits_by_cell.a1.self_consistency_gap_nats}{.3f}{0.349} against \verify{specs/006-calibrator-capabilities/review/evidence/windows_counts_006.json}{refits_by_cell.a1.r4_gap_wexac}{.3f}{0.346} on A1.
Naveh shipped $W$ with the limit stated, and extended the decision to realfull and A2 (spec 006, Clarifications);
re-centring $\hat\theta$ on the best point goes to jax2's next release.

\paragraph{Host dependence.} $W$ inherits 1.0.1's host dependence on human cohorts. A2's worst R-4 cohort
(\code{window-real-Sweden_M_1910_hetro-c1-s1}) ran on an \verify{tests/acceptance/results/gate_006/a2/window-real-Sweden_M_1910_hetro-c1-s1.json}{cpu_model}{}{AMD EPYC 9754 128-Core Processor}. Its Mac refit found a $\hat\theta$
\verify{specs/006-calibrator-capabilities/review/evidence/windows_counts_006.json}{host.mac_minus_wexac_loglik}{.3f}{0.047} nats better, with $\eta$ \verify{specs/006-calibrator-capabilities/review/evidence/windows_counts_006.json}{host.eta_pc_apart}{.0f}{5}\,\% and $\kappa$ \verify{specs/006-calibrator-capabilities/review/evidence/windows_counts_006.json}{host.kappa_ratio}{.0f}{5}\,$\times$ apart, both inside WEXAC's claims, and $\kappa\ll1$ on both
(\verify{specs/006-calibrator-capabilities/review/evidence/r4_refits_006.json}{refits.window-real-Sweden_M_1910_hetro-c1-s1.wexac_theta.kappa}{.2g}{0.0028} and \verify{specs/006-calibrator-capabilities/review/evidence/r4_refits_006.json}{refits.window-real-Sweden_M_1910_hetro-c1-s1.mac_theta.kappa}{.2g}{0.014}: saturated removal). All three refits put the Mac's $\hat\theta$ inside WEXAC's claims (\verify{specs/006-calibrator-capabilities/review/evidence/windows_counts_006.json}{host.refits_inside_wexac_claims}{d}{3} of
\verify{specs/006-calibrator-capabilities/review/evidence/windows_counts_006.json}{host.n_refits}{d}{3}); their own claims were not recorded. On a human life table, compare claims across machines, not $\hat\theta$.

\paragraph{Large cohorts.} A2 is the \verify{tests/acceptance/results/a2_006.json}{n_cohorts}{d}{44} cohorts planted at human life tables, with $n\ge25{,}000$, judged under the
collapse rule (too few claims separate 0.95 from 0.88, so the rule asks only that coverage did not collapse to 0.70). $W$'s
closed intervals cover the truth for $\eta$ on \verify{tests/acceptance/results/a2_006.json}{R-1.gated.eta.hits}{d}{31} of \verify{tests/acceptance/results/a2_006.json}{R-1.gated.eta.claimed}{d}{32} cohorts and for $\varepsilon$ on \verify{tests/acceptance/results/a2_006.json}{R-1.gated.epsilon.hits}{d}{28} of \verify{tests/acceptance/results/a2_006.json}{R-1.gated.epsilon.claimed}{d}{29}, and the truth
witness fails on none (R-3 passes: \verify{tests/acceptance/results/a2_006.json}{R-3.passes}{}{True}). But the data did not set the upper ends of most of them:
\verify{tests/acceptance/results/a2_006.json}{truncation.eta.truncated}{d}{28} of the \verify{tests/acceptance/results/a2_006.json}{truncation.eta.closed}{d}{32} for $\eta$ (\verify{specs/006-calibrator-capabilities/review/evidence/windows_counts_006.json}{truncation.a2.eta.box_limited}{d}{14} at the box, \verify{specs/006-calibrator-capabilities/review/evidence/windows_counts_006.json}{truncation.a2.eta.widened}{d}{13} widened, \verify{specs/006-calibrator-capabilities/review/evidence/windows_counts_006.json}{truncation.a2.eta.contradicted}{d}{1} below the lead pair's walk) and \verify{tests/acceptance/results/a2_006.json}{truncation.epsilon.truncated}{d}{29} of
the \verify{tests/acceptance/results/a2_006.json}{truncation.epsilon.closed}{d}{29} for $\varepsilon$ (\verify{specs/006-calibrator-capabilities/review/evidence/windows_counts_006.json}{truncation.a2.epsilon.box_limited}{d}{13} at the box, \verify{specs/006-calibrator-capabilities/review/evidence/windows_counts_006.json}{truncation.a2.epsilon.widened}{d}{16} widened), against a pre-registered 10\,\%. So $W$'s windowed
fits of 25{,}000 rows or more carry the warning \code{W_LARGE_N}.

\paragraph{The life table, not $n$.} A claim is truncated above when its upper end sits at the
box, was widened, or ends below the lead pair's walk or more than 0.3 dex below $H$'s. The records show that the variable
is the life table (Table~\ref{tab:win-trunc}). A1 holds A2's human tables at $n=10{,}000$, and they truncate as A2 does.
Below the split, the real-table cohorts' $\eta$ share (\verify{specs/006-calibrator-capabilities/review/evidence/windows_counts_006.json}{truncation.a1_real.eta.truncated}{d}{56} of \verify{specs/006-calibrator-capabilities/review/evidence/windows_counts_006.json}{truncation.a1_real.eta.closed}{d}{67}) does not differ from A2's (\verify{specs/006-calibrator-capabilities/review/evidence/windows_counts_006.json}{truncation.a2.eta.truncated}{d}{28} of \verify{specs/006-calibrator-capabilities/review/evidence/windows_counts_006.json}{truncation.a2.eta.closed}{d}{32}; Fisher
$p=\verify{specs/006-calibrator-capabilities/review/evidence/windows_counts_006.json}{truncation.fisher_p_eta_a1_real_vs_a2}{.2f}{0.77}$), and inside A1 it does not rise with $n$ (\verify{specs/006-calibrator-capabilities/review/evidence/windows_counts_006.json}{truncation.a1_eta_by_n.0-1999.truncated}{d}{43} of \verify{specs/006-calibrator-capabilities/review/evidence/windows_counts_006.json}{truncation.a1_eta_by_n.0-1999.closed}{d}{61} below 2{,}000, \verify{specs/006-calibrator-capabilities/review/evidence/windows_counts_006.json}{truncation.a1_eta_by_n.2000-9999.truncated}{d}{36} of \verify{specs/006-calibrator-capabilities/review/evidence/windows_counts_006.json}{truncation.a1_eta_by_n.2000-9999.closed}{d}{67} at 2{,}000--9{,}999, \verify{specs/006-calibrator-capabilities/review/evidence/windows_counts_006.json}{truncation.a1_eta_by_n.10000-24999.truncated}{d}{28} of \verify{specs/006-calibrator-capabilities/review/evidence/windows_counts_006.json}{truncation.a1_eta_by_n.10000-24999.closed}{d}{47} at
10{,}000--24{,}999). The two main causes differ. A box-limited end (label \texttt{open above: box-limited}) is the
model's: on these windows $\eta$, $\varepsilon$ and $\beta/\kappa$ slide together along a ridge, the box ends it, and more
deaths do not close it. A widened end (label \texttt{edge not found}) is the code's: $J$ 1.0.1's frozen profile walk
stopped short of points the fit had already evaluated, and the end moved out to them. Such an end can sit below the
truth: A2's only $\eta$ miss and only $\varepsilon$ miss are truths \verify{tests/acceptance/results/a2_006.json}{misses.eta.above.window-real-Sweden_F_1910_hetro-c0-s1}{.2f}{0.28} and \verify{tests/acceptance/results/a2_006.json}{misses.epsilon.above.window-real-Sweden_F_1910_hetro-c0-s1}{.2f}{0.31} dex above a widened end
(\code{window-real-Sweden_F_1910_hetro-c0-s1}). The rule at every $n$: an upper end labelled either way is ``at least
this'', not the interval's end; read the lead pair. \code{W_LARGE_N}'s message says so, gives A2's counts by cause and
the real-table rate below 25{,}000, and calls 25{,}000 the panel's split \src{warnings.py:WINDOW_WARNINGS}.

\begin{table}[H]\centering\small
\begin{tabular}{@{}llll@{}}
\toprule
Cohorts & $n$ & $\eta$ truncated & $\varepsilon$ truncated\\\midrule
A1, SR parameter grids & $<25{,}000$ & \verify{specs/006-calibrator-capabilities/review/evidence/windows_counts_006.json}{truncation.a1_synthetic.eta.truncated}{d}{51} of \verify{specs/006-calibrator-capabilities/review/evidence/windows_counts_006.json}{truncation.a1_synthetic.eta.closed}{d}{108} (\verify{specs/006-calibrator-capabilities/review/evidence/windows_counts_006.json}{truncation.a1_synthetic.eta.box_limited}{d}{47}, \verify{specs/006-calibrator-capabilities/review/evidence/windows_counts_006.json}{truncation.a1_synthetic.eta.widened}{d}{3}) & \verify{specs/006-calibrator-capabilities/review/evidence/windows_counts_006.json}{truncation.a1_synthetic.epsilon.truncated}{d}{36} of \verify{specs/006-calibrator-capabilities/review/evidence/windows_counts_006.json}{truncation.a1_synthetic.epsilon.closed}{d}{109} (\verify{specs/006-calibrator-capabilities/review/evidence/windows_counts_006.json}{truncation.a1_synthetic.epsilon.box_limited}{d}{30}, \verify{specs/006-calibrator-capabilities/review/evidence/windows_counts_006.json}{truncation.a1_synthetic.epsilon.widened}{d}{3})\\
A1, real life tables & $<25{,}000$ & \verify{specs/006-calibrator-capabilities/review/evidence/windows_counts_006.json}{truncation.a1_real.eta.truncated}{d}{56} of \verify{specs/006-calibrator-capabilities/review/evidence/windows_counts_006.json}{truncation.a1_real.eta.closed}{d}{67} (\verify{specs/006-calibrator-capabilities/review/evidence/windows_counts_006.json}{truncation.a1_real.eta.box_limited}{d}{50}, \verify{specs/006-calibrator-capabilities/review/evidence/windows_counts_006.json}{truncation.a1_real.eta.widened}{d}{5}) & \verify{specs/006-calibrator-capabilities/review/evidence/windows_counts_006.json}{truncation.a1_real.epsilon.truncated}{d}{37} of \verify{specs/006-calibrator-capabilities/review/evidence/windows_counts_006.json}{truncation.a1_real.epsilon.closed}{d}{51} (\verify{specs/006-calibrator-capabilities/review/evidence/windows_counts_006.json}{truncation.a1_real.epsilon.box_limited}{d}{32}, \verify{specs/006-calibrator-capabilities/review/evidence/windows_counts_006.json}{truncation.a1_real.epsilon.widened}{d}{5})\\
A1, A2's human tables & $10{,}000$ & \verify{specs/006-calibrator-capabilities/review/evidence/windows_counts_006.json}{truncation.a1_same_tables_as_a2_n10000.eta.truncated}{d}{14} of \verify{specs/006-calibrator-capabilities/review/evidence/windows_counts_006.json}{truncation.a1_same_tables_as_a2_n10000.eta.closed}{d}{15} (\verify{specs/006-calibrator-capabilities/review/evidence/windows_counts_006.json}{truncation.a1_same_tables_as_a2_n10000.eta.box_limited}{d}{8}, \verify{specs/006-calibrator-capabilities/review/evidence/windows_counts_006.json}{truncation.a1_same_tables_as_a2_n10000.eta.widened}{d}{5}) & \verify{specs/006-calibrator-capabilities/review/evidence/windows_counts_006.json}{truncation.a1_same_tables_as_a2_n10000.epsilon.truncated}{d}{8} of \verify{specs/006-calibrator-capabilities/review/evidence/windows_counts_006.json}{truncation.a1_same_tables_as_a2_n10000.epsilon.closed}{d}{8} (\verify{specs/006-calibrator-capabilities/review/evidence/windows_counts_006.json}{truncation.a1_same_tables_as_a2_n10000.epsilon.box_limited}{d}{4}, \verify{specs/006-calibrator-capabilities/review/evidence/windows_counts_006.json}{truncation.a1_same_tables_as_a2_n10000.epsilon.widened}{d}{4})\\
A2, human life tables & $\ge25{,}000$ & \verify{specs/006-calibrator-capabilities/review/evidence/windows_counts_006.json}{truncation.a2.eta.truncated}{d}{28} of \verify{specs/006-calibrator-capabilities/review/evidence/windows_counts_006.json}{truncation.a2.eta.closed}{d}{32} (\verify{specs/006-calibrator-capabilities/review/evidence/windows_counts_006.json}{truncation.a2.eta.box_limited}{d}{14}, \verify{specs/006-calibrator-capabilities/review/evidence/windows_counts_006.json}{truncation.a2.eta.widened}{d}{13}) & \verify{specs/006-calibrator-capabilities/review/evidence/windows_counts_006.json}{truncation.a2.epsilon.truncated}{d}{29} of \verify{specs/006-calibrator-capabilities/review/evidence/windows_counts_006.json}{truncation.a2.epsilon.closed}{d}{29} (\verify{specs/006-calibrator-capabilities/review/evidence/windows_counts_006.json}{truncation.a2.epsilon.box_limited}{d}{13}, \verify{specs/006-calibrator-capabilities/review/evidence/windows_counts_006.json}{truncation.a2.epsilon.widened}{d}{16})\\
realfull & \verify{specs/006-calibrator-capabilities/review/evidence/windows_counts_006.json}{truncation.realfull_drosophila_n_range[0]}{,d}{400}--\verify{tests/acceptance/results/gate_006/realfull/realfull-real-combined_human_F-c0-s0.json}{n}{,d}{100,000} & \verify{specs/006-calibrator-capabilities/review/evidence/windows_counts_006.json}{truncation.realfull.eta.truncated}{d}{18} of \verify{specs/006-calibrator-capabilities/review/evidence/windows_counts_006.json}{truncation.realfull.eta.closed}{d}{18} (\verify{specs/006-calibrator-capabilities/review/evidence/windows_counts_006.json}{truncation.realfull.eta.box_limited}{d}{16}, \verify{specs/006-calibrator-capabilities/review/evidence/windows_counts_006.json}{truncation.realfull.eta.widened}{d}{2}) & \verify{specs/006-calibrator-capabilities/review/evidence/windows_counts_006.json}{truncation.realfull.epsilon.truncated}{d}{15} of \verify{specs/006-calibrator-capabilities/review/evidence/windows_counts_006.json}{truncation.realfull.epsilon.closed}{d}{16} (\verify{specs/006-calibrator-capabilities/review/evidence/windows_counts_006.json}{truncation.realfull.epsilon.box_limited}{d}{13}, \verify{specs/006-calibrator-capabilities/review/evidence/windows_counts_006.json}{truncation.realfull.epsilon.widened}{d}{2})\\
\bottomrule
\end{tabular}
\caption{Closed claims truncated above, of all closed claims, by population; in brackets, those at the box and those
widened (\texttt{windows\_counts\_006.json}, \texttt{truncation}).}\label{tab:win-trunc}
\end{table}

\paragraph{The lead pair on windows.} $W$'s lead pair was truth-scored on windows (\code{Dg} in each record). On A1,
$\eta/\varepsilon$'s closed intervals cover the truth on \verify{tests/acceptance/results/a1_006.json}{Dg.by_quantity.eta_over_eps.hits}{d}{52} of \verify{tests/acceptance/results/a1_006.json}{Dg.by_quantity.eta_over_eps.claimed}{d}{54} cohorts and $\bar r/\eta$'s on \verify{tests/acceptance/results/a1_006.json}{Dg.by_quantity.rbar_over_eta.hits}{d}{20} of \verify{tests/acceptance/results/a1_006.json}{Dg.by_quantity.rbar_over_eta.claimed}{d}{21}, plus \verify{tests/acceptance/results/a1_006.json}{Dg.by_quantity.rbar_over_eta.one_sided[0]}{d}{32} of \verify{tests/acceptance/results/a1_006.json}{Dg.by_quantity.rbar_over_eta.one_sided[1]}{d}{32}
one-sided. On A2 $\eta/\varepsilon$ covers on \verify{tests/acceptance/results/a2_006.json}{Dg.by_quantity.eta_over_eps.hits}{d}{11} of \verify{tests/acceptance/results/a2_006.json}{Dg.by_quantity.eta_over_eps.claimed}{d}{11} and $\bar r/\eta$ on \verify{tests/acceptance/results/a2_006.json}{Dg.by_quantity.rbar_over_eta.hits}{d}{11} of \verify{tests/acceptance/results/a2_006.json}{Dg.by_quantity.rbar_over_eta.claimed}{d}{11}, all closed, on the human tables whose
$\eta$ and $\varepsilon$ upper ends the box or the search set. That count passes the collapse rule (cut \verify{tests/acceptance/results/a2_006.json}{Dg.by_quantity.eta_over_eps.cut}{.3f}{0.777}) and is too
small to show 0.95.

\paragraph{realfull.} The \verify{tests/acceptance/results/realfull_006.json}{n_cohorts}{d}{20} realfull window-only cohorts are generated cohorts planted at real life tables
(\verify{specs/006-calibrator-capabilities/review/evidence/windows_counts_006.json}{leverage.human.realfull_human.n}{d}{12} human; \verify{specs/006-calibrator-capabilities/review/evidence/windows_counts_006.json}{truncation.realfull_drosophila.eta.n_cohorts}{d}{8} Drosophila, at $n=\verify{specs/006-calibrator-capabilities/review/evidence/windows_counts_006.json}{truncation.realfull_drosophila_n_range[0]}{d}{400}$--\verify{specs/006-calibrator-capabilities/review/evidence/windows_counts_006.json}{truncation.realfull_drosophila_n_range[1]}{,d}{1,116}), reported, not gated. Their closed intervals cover the
truth for $\eta$ on \verify{tests/acceptance/results/realfull_006.json}{R-1.gated.eta.hits}{d}{17} of \verify{tests/acceptance/results/realfull_006.json}{R-1.gated.eta.claimed}{d}{18} and for $\varepsilon$ on \verify{tests/acceptance/results/realfull_006.json}{R-1.gated.epsilon.hits}{d}{16} of \verify{tests/acceptance/results/realfull_006.json}{R-1.gated.epsilon.claimed}{d}{16}, but \verify{tests/acceptance/results/realfull_006.json}{truncation.eta.truncated}{d}{18} of the \verify{tests/acceptance/results/realfull_006.json}{truncation.eta.closed}{d}{18} closed $\eta$ intervals are
truncated above (\verify{specs/006-calibrator-capabilities/review/evidence/windows_counts_006.json}{truncation.realfull.eta.box_limited}{d}{16} at the box, \verify{specs/006-calibrator-capabilities/review/evidence/windows_counts_006.json}{truncation.realfull.eta.widened}{d}{2} widened) and \verify{tests/acceptance/results/realfull_006.json}{truncation.epsilon.truncated}{d}{15} of \verify{tests/acceptance/results/realfull_006.json}{truncation.epsilon.closed}{d}{16} for $\varepsilon$ (\verify{specs/006-calibrator-capabilities/review/evidence/windows_counts_006.json}{truncation.realfull.epsilon.box_limited}{d}{13}, \verify{specs/006-calibrator-capabilities/review/evidence/windows_counts_006.json}{truncation.realfull.epsilon.widened}{d}{2}). The data set
$\eta$'s upper end on \verify{tests/acceptance/results/realfull_006.json}{shapes.eta.capped_by_data}{d}{0} of the \verify{tests/acceptance/results/realfull_006.json}{shapes.eta.n}{d}{20} cohorts and $\varepsilon$'s on \verify{tests/acceptance/results/realfull_006.json}{shapes.epsilon.capped_by_data}{d}{1}.

\paragraph{The leverage check on windows.} It raised no flag on the \verify{specs/006-calibrator-capabilities/review/evidence/windows_counts_006.json}{leverage.assessed_total}{d}{209} fits it assessed (\verify{specs/006-calibrator-capabilities/review/evidence/windows_counts_006.json}{leverage.flags_total}{d}{0} flags; \verify{specs/006-calibrator-capabilities/review/evidence/windows_counts_006.json}{leverage.slides_total}{d}{52} of
them flat-valley slides). No case was planted on windows, so its detection there is untested. On cohorts planted at
human life tables it could not assess \verify{specs/006-calibrator-capabilities/review/evidence/windows_counts_006.json}{leverage.human.not_assessed_total}{d}{47} of \verify{specs/006-calibrator-capabilities/review/evidence/windows_counts_006.json}{leverage.human.n_total}{d}{78} fits (\verify{specs/006-calibrator-capabilities/review/evidence/windows_counts_006.json}{leverage.human.a1_human_n10000.not_assessed}{d}{13} of \verify{specs/006-calibrator-capabilities/review/evidence/windows_counts_006.json}{leverage.human.a1_human_n10000.n}{d}{22} at $n=10{,}000$ in A1, \verify{specs/006-calibrator-capabilities/review/evidence/windows_counts_006.json}{leverage.human.a2.not_assessed}{d}{26} of \verify{specs/006-calibrator-capabilities/review/evidence/windows_counts_006.json}{leverage.human.a2.n}{d}{44} in A2, \verify{specs/006-calibrator-capabilities/review/evidence/windows_counts_006.json}{leverage.human.realfull_human.not_assessed}{d}{8} of \verify{specs/006-calibrator-capabilities/review/evidence/windows_counts_006.json}{leverage.human.realfull_human.n}{d}{12}
in realfull): $\hat\theta$ was not stationary there.


\end{document}
